% ============================================================================
%  Plataforma Unooria — PU-PRD-001
%  Mapa de produto: épicos, features e histórias de usuário.
%
%  Documento de PRODUTO (do PO para o time), não de engenharia: fala de
%  pessoas, telas e regras de negócio. O detalhe técnico mora na PU-DOC-001
%  (../../main.tex); a ponte entre os dois é o Apêndice B deste documento.
%
%  Compilação (a partir desta pasta):
%     tectonic -k mapa-de-produto.tex
%  Reaproveita o preâmbulo da PU-DOC-001 (../../preambulo) para ter as mesmas
%  cores, caixas e tabelas.
% ============================================================================
\documentclass[11pt, a4paper, oneside, openany]{report}

% ============================================================================
%  estilo.tex — identidade visual, tipografia e componentes do documento
%
%  Tipografia: Roboto (texto e titulos) + Roboto Mono (codigo). As duas vem do
%  pacote `roboto`, entao o documento compila com pdfLaTeX, XeLaTeX ou Tectonic
%  sem fontspec explicito.
%
%  Paleta: os tokens `--unooria-*` do @unooria/design-system, os MESMOS valores
%  que `packages/ui/styles/plataforma.css` usa como fallback. Cor nova neste
%  documento so entra se existir como token no pacote.
% ============================================================================

% --- Codificacao e idioma ---------------------------------------------------
\usepackage[T1]{fontenc}
\usepackage[utf8]{inputenc}
\usepackage[shorthands=off,brazil]{babel}   % shorthands OFF: com babel-pt o `"` fica ativo e quebra codigo em linha
\usepackage{csquotes}

% --- Tipografia -------------------------------------------------------------
\usepackage[sfdefault]{roboto}            % Roboto como familia padrao
\usepackage[scaled=0.90]{roboto-mono}     % Roboto Mono para codigo
\usepackage{microtype}
\usepackage{textcomp}

% --- Pagina -----------------------------------------------------------------
\usepackage[
  a4paper,
  top=2.6cm, bottom=2.5cm, left=2.6cm, right=2.4cm,
  headheight=15pt, headsep=18pt, footskip=26pt
]{geometry}
\usepackage{pdflscape}
\usepackage{setspace}
\onehalfspacing

% --- Cores canônicas do Brandbook Unooria ("BLUNOORIA") -----------------------
% Regra semântica: "A cor não decora. Ela organiza."
%   - Azul Profundo (#0E1039): Estrutura, corpo, regras e títulos
%   - Lavanda (#8B85AC): Inteligência, decisões arquiteturais (ADR), IA
%   - Azul Claro (#7FA5E2): Interface, navegação, referências e notas
\usepackage{xcolor}
\definecolor{unooNavy}        {HTML}{0E1039}  % Azul Profundo — Estrutura (C100 M96 Y44 K56)
\definecolor{unooLavender}    {HTML}{8B85AC}  % Lavanda — Inteligência (C52 M48 Y15 K2)
\definecolor{unooBlue}        {HTML}{7FA5E2}  % Azul Claro — Interface (C54 M29 Y0 K0)
\definecolor{unooBlueLight}   {HTML}{B8D0F5}  % Variação clara de interface
\definecolor{unooBlueWash}    {HTML}{F0F4FC}  % Fundo suave de interface
\definecolor{unooLavenderWash}{HTML}{F3F2F8}  % Fundo suave de inteligência
\definecolor{unooLavenderSoft}{HTML}{5D577D}  % Lavanda escurecido p/ alto contraste tipográfico
\definecolor{unooInk}         {HTML}{0E1039}  % neutral-950 / texto primário ancorado na marca
\definecolor{unooInkSoft}     {HTML}{3A3C59}  % neutral-700 / texto secundário
\definecolor{unooMuted}       {HTML}{6D6E82}  % neutral-500 / legendas e metadados
\definecolor{unooGrid}        {HTML}{E3E5EE}  % Linhas de grade e bordas sutis
\definecolor{unooAxis}        {HTML}{CACDDC}  % Réguas e separadores
\definecolor{unooSurface}     {HTML}{FFFFFF}  % Fundo puro
\definecolor{unooPaper}       {HTML}{F7F8FC}  % Fundo alternado / zebra
% Tons funcionais (alertas de sistema e status):
\definecolor{unooRed}         {HTML}{C92A2A}  % danger-600
\definecolor{unooRedWash}     {HTML}{FDF2F2}
\definecolor{unooOrange}      {HTML}{D97706}  % warning-600
\definecolor{unooOrangeWash}  {HTML}{FFFBEB}
\definecolor{unooGreen}       {HTML}{2F855A}  % Status funcional aprovado
\definecolor{unooGreenWash}   {HTML}{F0FDF4}
\definecolor{unooYellow}      {HTML}{B45309}
\definecolor{unooYellowWash}  {HTML}{FEF3C7}
\definecolor{unooPink}        {HTML}{9D4EDD}

% --- Tabelas ----------------------------------------------------------------
\usepackage{booktabs}
\usepackage{tabularx}
\usepackage{longtable}
\usepackage{array}
\usepackage{multirow}
\usepackage{makecell}
\usepackage{colortbl}
\usepackage{ltablex}          % tabularx com quebra de pagina
\keepXColumns
\renewcommand{\arraystretch}{1.22}
\arrayrulecolor{unooAxis}
\setlength{\aboverulesep}{0pt}
\setlength{\belowrulesep}{0pt}
\setlength{\extrarowheight}{2pt}

% Tabelas em \small por decisao de composicao: as celulas deste documento
% misturam prosa e identificadores monoespacados longos (caminhos de arquivo,
% nomes de coluna), que NAO hifenizam. Em corpo 11 elas transbordavam a caixa;
% em \small cabem, e a tabela continua legivel porque a densidade e alta de
% proposito. `\emergencystretch` cobre o residuo.
\usepackage{etoolbox}
\AtBeginEnvironment{tabularx}{\small\setlength{\emergencystretch}{1.2em}}
\AtBeginEnvironment{longtable}{\small\setlength{\emergencystretch}{1.2em}}

\newcolumntype{L}[1]{>{\raggedright\arraybackslash}p{#1}}
\newcolumntype{C}[1]{>{\centering\arraybackslash}p{#1}}
\newcolumntype{R}[1]{>{\raggedleft\arraybackslash}p{#1}}
\newcolumntype{Y}{>{\raggedright\arraybackslash}X}

% Colunas X sem justificacao: as celulas misturam prosa e codigo monoespacado,
% que nao hifeniza — justificar produziria espacos internos muito largos.
\renewcommand{\tabularxcolumn}[1]{>{\raggedright\arraybackslash}p{#1}}

\newcommand{\thd}[1]{\textcolor{white}{\textbf{#1}}}
\newcommand{\headrow}{\rowcolor{unooNavy}}
\newcommand{\zebra}{\rowcolor{unooPaper}}

% --- Listas -----------------------------------------------------------------
\usepackage{enumitem}
\setlist[itemize]{leftmargin=1.35em, itemsep=2pt, topsep=4pt, parsep=0pt,
                  label=\textcolor{unooBlue}{\raisebox{0.18ex}{\rule{0.40em}{0.40em}}}}
\setlist[enumerate]{leftmargin=1.7em, itemsep=2pt, topsep=4pt, parsep=0pt,
                    font=\bfseries\color{unooNavy}}
\setlist[description]{leftmargin=1.4em, style=nextline,
                      font=\bfseries\color{unooNavy}}

% --- Matematica e unidades --------------------------------------------------
\usepackage{amsmath}
\usepackage{amssymb}
\usepackage{siunitx}
\sisetup{output-decimal-marker={,}, group-separator={.}, group-minimum-digits=4,
         per-mode=symbol, detect-weight=true, detect-family=true}
\DeclareSIUnit{\tokens}{tokens}
\DeclareSIUnit{\req}{req}

% --- Figuras e legendas -----------------------------------------------------
\usepackage{graphicx}
\usepackage{float}
\usepackage{adjustbox}
\usepackage[font=small, labelfont={bf,color=unooNavy}, textfont=normalfont,
            labelsep=endash, justification=raggedright, singlelinecheck=false,
            skip=8pt]{caption}
\usepackage{subcaption}
\captionsetup[table]{position=top, skip=6pt}

% Ajusta a figura a largura do texto apenas quando ela excede a caixa.
\newcommand{\fitwidth}[1]{\adjustbox{max width=\linewidth, center}{#1}}
\newcommand{\fitpage}[1]{\adjustbox{max width=\linewidth, max height=0.82\textheight, center}{#1}}

% --- Blocos de destaque (tcolorbox) ----------------------------------------
\usepackage[most]{tcolorbox}
\tcbuselibrary{skins,breakable}

\tcbset{
  unoobase/.style={
    enhanced, breakable, sharp corners,
    boxrule=0pt, leftrule=2.6pt, arc=0pt,
    left=10pt, right=10pt, top=7pt, bottom=7pt,
    fonttitle=\bfseries\small,
    before skip=10pt, after skip=10pt
  }
}

% `colbacktitle` e OBRIGATORIO aqui: sem ele o tcolorbox pinta a faixa do
% titulo com a cor de `colframe`, e o titulo (que usa a MESMA cor em
% `coltitle`) fica invisivel — texto vermelho sobre faixa vermelha. O sintoma
% e uma barra colorida solida no lugar do rotulo do bloco.
\newtcolorbox{nota}[1][]{unoobase, colback=unooBlueWash, colframe=unooBlue,
  colbacktitle=unooBlueWash, coltitle=unooNavy, titlerule=0pt, title={NOTA}, #1}
\newtcolorbox{atencao}[1][]{unoobase, colback=unooOrangeWash, colframe=unooOrange,
  colbacktitle=unooOrangeWash, coltitle=unooOrange, titlerule=0pt,
  title={ATENÇÃO}, #1}
\newtcolorbox{decisao}[1][]{unoobase, colback=unooLavenderWash, colframe=unooLavender,
  colbacktitle=unooLavenderWash, coltitle=unooNavy, titlerule=0pt,
  title={DECISÃO DE PROJETO}, #1}
\newtcolorbox{limitacao}[1][]{unoobase, colback=unooRedWash, colframe=unooRed,
  colbacktitle=unooRedWash, coltitle=unooRed, titlerule=0pt,
  title={LIMITAÇÃO CONHECIDA}, #1}
\newtcolorbox{regra}[1][]{unoobase, colback=unooRedWash, colframe=unooRed,
  colbacktitle=unooRedWash, coltitle=unooRed, titlerule=0pt,
  title={REGRA INVIOLÁVEL}, #1}
\newtcolorbox{destaque}[1][]{unoobase, colback=unooPaper, colframe=unooLavender, #1}
\newtcolorbox{especificacao}[1]{unoobase, colback=unooSurface, colframe=unooNavy,
  colbacktitle=unooSurface, coltitle=unooNavy, titlerule=0pt, title={#1}}

% --- Codigo-fonte -----------------------------------------------------------
\usepackage{listings}

\definecolor{lstBg}     {HTML}{FAFBFB}
\definecolor{lstFrame}  {HTML}{E4E9EB}
\definecolor{lstKeyword}{HTML}{2F4F7F}
\definecolor{lstString} {HTML}{0F6B52}
\definecolor{lstComment}{HTML}{617379}
\definecolor{lstEmph}   {HTML}{A9490F}

\lstdefinestyle{unoo}{
  backgroundcolor=\color{lstBg},
  basicstyle=\ttfamily\scriptsize\color{unooInk},
  keywordstyle=\color{lstKeyword}\bfseries,
  keywordstyle=[2]\color{lstEmph}\bfseries,
  keywordstyle=[3]\color{unooBlue},
  stringstyle=\color{lstString},
  commentstyle=\color{lstComment}\itshape,
  numberstyle=\ttfamily\tiny\color{unooMuted},
  identifierstyle=\color{unooInk},
  numbers=left, numbersep=9pt, stepnumber=1,
  frame=single, framesep=6pt, rulecolor=\color{lstFrame},
  xleftmargin=17pt, xrightmargin=2pt,
  breaklines=true, breakatwhitespace=false, breakindent=0pt,
  postbreak=\mbox{\textcolor{unooBlue}{$\hookrightarrow$}\space},
  showstringspaces=false, tabsize=2,
  columns=fullflexible, keepspaces=true, upquote=true,
  aboveskip=10pt, belowskip=10pt,
  literate=
    {á}{{\'a}}1 {à}{{\`a}}1 {ã}{{\~a}}1 {â}{{\^a}}1
    {é}{{\'e}}1 {ê}{{\^e}}1 {í}{{\'i}}1
    {ó}{{\'o}}1 {õ}{{\~o}}1 {ô}{{\^o}}1
    {ú}{{\'u}}1 {ü}{{\"u}}1 {ç}{{\c c}}1
    {Á}{{\'A}}1 {À}{{\`A}}1 {Ã}{{\~A}}1 {Â}{{\^A}}1
    {É}{{\'E}}1 {Ê}{{\^E}}1 {Í}{{\'I}}1
    {Ó}{{\'O}}1 {Õ}{{\~O}}1 {Ô}{{\^O}}1
    {Ú}{{\'U}}1 {Ç}{{\c C}}1
    {°}{{\textdegree}}1 {º}{{\textordmasculine}}1
    {µ}{{$\mu$}}1 {×}{{$\times$}}1 {≈}{{$\approx$}}1
    {≥}{{$\geq$}}1 {≤}{{$\leq$}}1 {≠}{{$\neq$}}1
    {→}{{$\rightarrow$}}1 {←}{{$\leftarrow$}}1 {⇒}{{$\Rightarrow$}}1
    {—}{{---}}3 {–}{{--}}2 {…}{{\ldots}}3
    {“}{{\textquotedblleft}}1 {”}{{\textquotedblright}}1
    {‘}{{\textquoteleft}}1 {’}{{\textquoteright}}1
    {ª}{{\textordfeminine}}1 {§}{{\S}}1
}
\lstset{style=unoo}

\lstdefinelanguage{tsunoo}[]{Java}{
  morekeywords={const,let,export,default,import,from,as,interface,type,await,
                async,function,return,string,number,boolean,readonly,unknown,
                never,satisfies,enum,declare,namespace,implements},
  morekeywords=[2]{z,Prisma,useEffect,postMessage},
  sensitive=true,
  morecomment=[l]{//},
  morecomment=[s]{/*}{*/},
  morestring=[b]",
  morestring=[b]',
  morestring=[b]`
}
\lstdefinelanguage{sqlunoo}[]{SQL}{
  morekeywords={POLICY,ENABLE,FORCE,ROW,LEVEL,SECURITY,USING,PARTITION,RANGE,
                CONFLICT,DO,NOTHING,JSONB,TIMESTAMPTZ,current_setting,
                set_config,ATTACH,DEFAULT,GENERATED,ALWAYS},
  sensitive=false,
  morecomment=[l]{--},
  morestring=[b]'
}
\lstdefinelanguage{yamlunoo}{
  morekeywords={true,false,null,services,image,build,command,environment,
                volumes,ports,depends_on,healthcheck,condition,runs-on,steps,
                uses,needs,jobs,name,on,with,env,if,run,concurrency,
                permissions,workflow_call,workflow_dispatch},
  sensitive=true,
  morecomment=[l]{\#},
  morestring=[b]",
  morestring=[b]'
}
\lstdefinelanguage{jsonunoo}{
  morekeywords={true,false,null},
  keywordstyle=\color{lstEmph}\bfseries,
  sensitive=true,
  morestring=[b]"
}
\lstdefinelanguage{mermaid}{
  morekeywords={flowchart,graph,sequenceDiagram,erDiagram,classDiagram,gantt,
            subgraph,end,participant,actor,alt,else,loop,par,opt,note,activate,
            deactivate,direction,TB,LR,RL,BT,Person,Person_Ext,System,System_Ext,
            SystemDb,Container,ContainerDb,ContainerQueue,Component,ComponentDb,
            Rel,Rel_D,Rel_U,Rel_L,Rel_R,System_Boundary,Container_Boundary,
            Enterprise_Boundary,Boundary,Deployment_Node,UpdateLayoutConfig,
            UpdateRelStyle,classDef,class,style,state,C4Context,C4Container,
            C4Component,C4Deployment,stateDiagram,stateDiagram-v2},
  sensitive=true,
  morecomment=[l]{\%\%},
  morestring=[b]"
}
\lstdefinelanguage{bashunoo}[]{bash}{
  morekeywords={pnpm,docker,gcloud,node,psql,tectonic,turbo,prisma,gsutil},
  sensitive=true
}

% Codigo em linha.
%
% Uso obrigatorio de \texttt e \textcolor (e nao de {\ttfamily\color{...}}):
% os comandos \text... emitem \leavevmode antes de trocar a fonte, o que faz o
% paragrafo comecar na fonte do texto corrente. Sem isso, uma celula de tabela
% que inicia com codigo teria a primeira linha de base calculada com as metricas
% da fonte monoespacada e ficaria visivelmente desalinhada das demais colunas.
\newcommand{\cd}[1]{\texttt{\textcolor{unooNavy}{#1}}}
\newcommand{\arq}[1]{\texttt{\textcolor{unooInkSoft}{#1}}}
\newcommand{\rota}[1]{\texttt{\textcolor{unooGreen}{#1}}}
\newcommand{\env}[1]{\texttt{\textcolor{lstEmph}{#1}}}
% Ponto de quebra de linha sem hifen, para caminho longo em celula estreita:
% identificador monoespacado nao hifeniza, e sem isto ele transborda a caixa.
\newcommand{\B}{\discretionary{}{}{}}

% --- Titulos ----------------------------------------------------------------
\usepackage{titlesec}
\usepackage{titletoc}

\titleformat{\chapter}[display]
  {\normalfont}
  {\raggedleft\fontsize{62}{62}\selectfont\color{unooBlueLight}\textbf{\thechapter}}
  {-10pt}
  {\raggedright\fontsize{23}{27}\selectfont\bfseries\color{unooNavy}}
  [\vspace{5pt}{\color{unooBlue}\rule{\linewidth}{2.2pt}}]
\titlespacing*{\chapter}{0pt}{-30pt}{26pt}

\titleformat{name=\chapter,numberless}[display]
  {\normalfont}
  {}
  {0pt}
  {\raggedright\fontsize{23}{27}\selectfont\bfseries\color{unooNavy}}
  [\vspace{5pt}{\color{unooBlue}\rule{\linewidth}{2.2pt}}]

\titleformat{\section}
  {\normalfont\Large\bfseries\color{unooNavy}}{\thesection}{0.7em}{}
\titleformat{\subsection}
  {\normalfont\large\bfseries\color{unooInk}}{\thesubsection}{0.7em}{}
\titleformat{\subsubsection}
  {\normalfont\normalsize\bfseries\color{unooInkSoft}}{\thesubsubsection}{0.7em}{}
\titleformat{\paragraph}[runin]
  {\normalfont\normalsize\bfseries\color{unooNavy}}{}{0pt}{}[\;---\;]

\titlespacing*{\section}{0pt}{20pt}{8pt}
\titlespacing*{\subsection}{0pt}{15pt}{6pt}
\titlespacing*{\subsubsection}{0pt}{12pt}{5pt}
\titlespacing*{\paragraph}{0pt}{9pt}{0.4em}

\setcounter{secnumdepth}{3}
\setcounter{tocdepth}{2}

% --- Sumario ----------------------------------------------------------------
\titlecontents{chapter}[0em]
  {\addvspace{9pt}\bfseries\color{unooNavy}}
  {\contentslabel[\thecontentslabel]{1.5em}}
  {}
  {\hfill\contentspage}
\titlecontents{section}[1.7em]
  {\addvspace{1.5pt}}
  {\contentslabel{2.2em}}
  {}
  {\titlerule*[0.7pc]{\color{unooAxis}.}\contentspage}
\titlecontents{subsection}[4.0em]
  {\addvspace{0.5pt}\small\color{unooInkSoft}}
  {\contentslabel{2.8em}}
  {}
  {\titlerule*[0.7pc]{\color{unooGrid}.}\contentspage}

% --- Cabecalho e rodape -----------------------------------------------------
\usepackage{fancyhdr}
\usepackage{lastpage}

\newcommand{\docid}{PU-DOC-001}
\newcommand{\docversao}{1.2}

\renewcommand{\headrule}{{\color{unooGrid}\hrule height \headrulewidth}}
\renewcommand{\footrule}{{\color{unooGrid}\hrule height \footrulewidth}}

\fancypagestyle{unoomain}{%
  \fancyhf{}
  \fancyhead[L]{\footnotesize\color{unooInkSoft}\nouppercase{\leftmark}}
  \fancyhead[R]{\raisebox{-0.15ex}{\includegraphics[height=1.8ex]{figuras/marca/simbolo_azul_profundo.png}}\hspace{4pt}\footnotesize\bfseries\color{unooNavy}Unooria}
  \fancyfoot[L]{\footnotesize\color{unooMuted}\docid{} \textperiodcentered{} v\docversao}
  \fancyfoot[R]{\footnotesize\color{unooInkSoft}\thepage\,/\,\pageref{LastPage}}
  \renewcommand{\headrulewidth}{0.5pt}
  \renewcommand{\footrulewidth}{0.5pt}
}
% Pre-textual: numeracao romana, sem total de paginas (o total do
% \pageref{LastPage} refere-se a numeracao arabe do corpo do documento).
\fancypagestyle{unoofront}{%
  \fancyhf{}
  \fancyhead[L]{\footnotesize\color{unooInkSoft}\nouppercase{\leftmark}}
  \fancyhead[R]{\raisebox{-0.15ex}{\includegraphics[height=1.8ex]{figuras/marca/simbolo_azul_profundo.png}}\hspace{4pt}\footnotesize\bfseries\color{unooNavy}Unooria}
  \fancyfoot[L]{\footnotesize\color{unooMuted}\docid{} \textperiodcentered{} v\docversao}
  \fancyfoot[R]{\footnotesize\color{unooInkSoft}\thepage}
  \renewcommand{\headrulewidth}{0.5pt}
  \renewcommand{\footrulewidth}{0.5pt}
}
\fancypagestyle{plain}{%
  \fancyhf{}
  \fancyfoot[R]{\footnotesize\color{unooInkSoft}\thepage}
  \renewcommand{\headrulewidth}{0pt}
  \renewcommand{\footrulewidth}{0pt}
}
\pagestyle{unoofront}
\renewcommand{\chaptermark}[1]{\markboth{\thechapter\quad #1}{}}

% --- Referencias cruzadas e PDF --------------------------------------------
\usepackage[
  colorlinks=true,
  linkcolor=unooNavy,
  citecolor=unooBlue,
  urlcolor=unooBlue,
  filecolor=unooBlue,
  bookmarksnumbered=true,
  bookmarksopen=true,
  bookmarksopenlevel=1,
  pdfdisplaydoctitle=true
]{hyperref}
\usepackage{bookmark}
\usepackage{url}
\urlstyle{sf}

% --- Bibliografia (ABNT NBR 6023) ------------------------------------------
%
% `abntex2cite` + BibTeX, e nao biblatex+biber: o pipeline de compilacao
% (Tectonic) chama o BibTeX classico automaticamente, sem ferramenta externa.
% A saida sai formatada em ABNT do mesmo jeito.
\usepackage[alf]{abntex2cite}

% --- Semantica de referencia do projeto ------------------------------------
\newcommand{\rf}[1]{\textbf{\textcolor{unooNavy}{#1}}}
\newcommand{\adrr}[1]{\textbf{\textcolor{unooLavenderSoft}{#1}}}
\newcommand{\ucr}[1]{\textbf{\textcolor{unooBlue}{#1}}}
\newcommand{\drr}[1]{\textbf{\textcolor{unooOrange}{#1}}}
\newcommand{\yes}{\textcolor{unooGreen}{\textbf{\checkmark}}}
\newcommand{\no}{\textcolor{unooOrange}{\textbf{--}}}

% Setas: a Roboto NAO tem U+2192 nem U+21D2 (medido — o XeTeX avisa
% "Missing character" e o glifo simplesmente nao e desenhado). Em prosa,
% sempre estes comandos; nunca a seta literal.
\newcommand{\seta}{\ensuremath{\rightarrow}}
\newcommand{\implica}{\ensuremath{\Rightarrow}}
\newcommand{\setabi}{\ensuremath{\leftrightarrow}}

% Espacamento de paragrafo
\setlength{\parindent}{0pt}
\setlength{\parskip}{7pt}

\widowpenalty=10000
\clubpenalty=10000
\flushbottom

\graphicspath{{../../}}
\usetikzlibrary{positioning, arrows.meta}
\usepackage{xurl}
\raggedbottom

% Só os capítulos (épicos) são numerados. Features e histórias levam o
% próprio identificador (F01.1, H01.1.1): numeração do LaTeX ao lado dele
% seria um segundo número para a mesma coisa.
\setcounter{secnumdepth}{0}
\setcounter{tocdepth}{1}

\renewcommand{\docid}{PU-PRD-001}
\renewcommand{\docversao}{0.1}

\hypersetup{
  pdftitle={Mapa de produto da Plataforma Unooria},
  pdfauthor={Lucas Salomao},
  pdfsubject={Epicos, features e historias de usuario do que a plataforma ja faz e do que falta},
  pdfcreator={Tectonic / XeTeX}
}

% ------------------------------------------------------------ vocabulário
% Selo de estado. \colorbox e não tcbox: o selo vai dentro do título de
% outra caixa, e caixa dentro de caixa quebra a linha de base.
\newcommand{\selo}[3]{{\setlength{\fboxsep}{1.6pt}\colorbox{#2}{\scriptsize\bfseries\textcolor{#3}{\,#1\,}}}}
\newcommand{\pronto}{\selo{PRONTO}{unooGreenWash}{unooGreen}}
\newcommand{\parcial}{\selo{PARCIAL}{unooOrangeWash}{unooOrange}}
\newcommand{\afazer}{\selo{A FAZER}{unooLavenderWash}{unooLavenderSoft}}
\newcommand{\id}[1]{\textbf{\textcolor{unooLavenderSoft}{#1}}}

% Ficha do épico: uma caixa com os campos em lista de descrição.
\newtcolorbox{fichaepico}{unoobase, colback=unooPaper, colframe=unooNavy,
  colbacktitle=unooPaper, coltitle=unooNavy, titlerule=0pt,
  title={FICHA DO ÉPICO}}
\newlist{campos}{description}{1}
\setlist[campos]{font=\bfseries\color{unooNavy}, leftmargin=30mm,
  labelindent=0pt, labelwidth=28mm, labelsep=2mm, align=left, itemsep=3pt,
  topsep=2pt}

% Feature: seção com identificador e estado logo abaixo do título.
\newcommand{\feature}[3]{%
  \section[#1\quad #2]{\textcolor{unooLavenderSoft}{#1}\quad #2}%
  \noindent{\small\color{unooMuted}Estado da feature:}\enspace #3\par\medskip}

% História de usuário: cartão com identificador, título e estado.
\newtcolorbox{historia}[3]{unoobase, colback=unooSurface, colframe=unooBlue,
  colbacktitle=unooBlueWash, coltitle=unooNavy, titlerule=0pt,
  fonttitle=\bfseries\small,
  title={\textcolor{unooLavenderSoft}{#1}\enspace #2\hfill #3}}
\newcommand{\hu}[3]{\textbf{Como} #1, \textbf{quero} #2, \textbf{para} #3.\par}
\newcommand{\rotulo}[1]{\par\smallskip{\footnotesize\bfseries\color{unooLavenderSoft}#1}\par\nopagebreak}
\newcommand{\aceite}{\rotulo{CRITÉRIOS DE ACEITE}}
\newcommand{\aceiteproposta}{\rotulo{CRITÉRIOS DE ACEITE --- PROPOSTA A VALIDAR}}
\newcommand{\pendencia}[1]{\par\smallskip{\small\textcolor{unooOrange}{\textbf{Pendência:}}\enspace #1}\par}
\setlist[itemize]{nosep, leftmargin=*}

\begin{document}

% ------------------------------------------------------------------- capa
\begin{titlepage}
\thispagestyle{empty}
\begin{tikzpicture}[remember picture, overlay]
  \fill[unooNavy] (current page.north west) rectangle
        ([yshift=-12mm]current page.north east);
  \fill[unooLavender] ([yshift=-12mm]current page.north west) rectangle
        ([yshift=-13.8mm]current page.north east);
  \fill[unooBlue] ([yshift=-13.8mm]current page.north west) rectangle
        ([yshift=-15.0mm]current page.north east);
\end{tikzpicture}
\vspace*{14mm}

\noindent\includegraphics[width=50mm]{figuras/marca/logo_azul_profundo.png}\par
\vspace{18mm}

{\footnotesize\bfseries\color{unooLavenderSoft} DOCUMENTO DE PRODUTO \quad\textperiodcentered\quad PU-PRD-001\par}
\vspace{4mm}
{\fontsize{26}{31}\selectfont\bfseries\color{unooNavy}
 Mapa de produto da\\Plataforma Unooria\par}
\vspace{6mm}
{\large\color{unooInkSoft}
 O que a plataforma é, para quem existe e o que já faz --- organizado em
 épicos, features e histórias de usuário, com o que está pronto, o que está
 pela metade e o que falta\par}

\vfill
\begin{tabularx}{\linewidth}{@{}L{38mm}X@{}}
\toprule
Documento & PU-PRD-001, versão 0.1\\
Data & 25 de setembro de 2026\\
Autor & Lucas Salomão\\
Para quem & Produto (PO) e time de desenvolvimento\\
Natureza & Retrato do produto \emph{como ele está hoje}, levantado a partir
  das telas, das regras e dos testes de ponta a ponta da plataforma. Não
  substitui a documentação técnica (PU-DOC-001); fala a língua de produto.\\
\bottomrule
\end{tabularx}
\end{titlepage}

\pagenumbering{roman}
\tableofcontents
\cleardoublepage
\pagenumbering{arabic}
\setcounter{page}{1}
\pagestyle{unoomain}

% ===========================================================================
%  1 — Como ler este documento
% ===========================================================================
\chapter{Como ler este documento}
\label{cap:como-ler}

A Plataforma Unooria cresceu depressa: hoje ela tem nove missões para
estudantes, três ferramentas para professores, uma área inteira de gestão da
rede e um painel para a equipe da Unooria. Tudo isso foi construído, testado e
documentado do ponto de vista técnico --- mas nunca tinha sido descrito do
ponto de vista de \emph{produto}: quem usa, para quê, e o que conta como
\enquote{pronto}.

Este documento preenche essa lacuna. Ele é o \textbf{retrato do produto em
25 de setembro de 2026}, levantado a partir das telas, das regras de negócio e
dos testes de ponta a ponta que a plataforma já passa. Serve para três coisas:

\begin{enumerate}
  \item \textbf{Dar a todos a mesma visão do produto} --- o que ele é, para
    quem existe e quais promessas ele faz (Capítulo~\ref{cap:produto}).
  \item \textbf{Ser o backlog de partida} --- cada história tem estado, e o que
    está pela metade ou por fazer está consolidado no
    Capítulo~\ref{cap:backlog}.
  \item \textbf{Registrar as perguntas que só o produto pode responder} ---
    o Capítulo~\ref{cap:decisoes} lista as decisões em aberto encontradas no
    levantamento.
\end{enumerate}

\section{A hierarquia: épico, feature, história}

\begin{tabularx}{\linewidth}{@{}L{24mm}X L{30mm}@{}}
\toprule
\headrow \thd{Nível} & \thd{O que é} & \thd{Identificador}\\
\midrule
\textbf{Épico} & Um grande objetivo do produto, que atravessa várias telas e
  papéis. Ex.: \enquote{Cadastrar pessoas e turmas}. & \id{E01} a \id{E14}\\
\zebra \textbf{Feature} & Uma capacidade concreta dentro do épico, que alguém
  consegue usar e reconhecer. Ex.: \enquote{Importar pessoas por planilha}. &
  \id{F03.1} (épico 03, feature 1)\\
\textbf{História} & Uma necessidade de uma pessoa, contada do ponto de vista
  dela, com critérios que dizem quando está atendida. & \id{H03.1.2} (feature
  F03.1, história 2)\\
\bottomrule
\end{tabularx}

\section{O formato de cada história}

Toda história segue o mesmo molde, e o molde é o mesmo para o que já existe e
para o que falta:

\begin{historia}{H00.0.0}{Título curto da necessidade}{\pronto}
\hu{\emph{quem} (o papel, com o nome que a tela usa)}{\emph{o quê} (a ação, em
palavras de quem usa)}{\emph{por quê} (o valor para essa pessoa ou para a rede)}
\aceite
\begin{itemize}
  \item Cada critério é algo que dá para \textbf{verificar olhando a tela}.
  \item Os textos entre aspas são os que a interface mostra hoje, palavra por
    palavra.
\end{itemize}
\pendencia{quando algo da história ainda falta, aparece aqui.}
\end{historia}

Nas histórias marcadas \afazer, os critérios de aceite são uma
\textbf{proposta} para o PO validar, não uma decisão tomada --- o rótulo deixa
isso explícito.

\section{Os três estados}

\begin{tabularx}{\linewidth}{@{}L{24mm}X@{}}
\toprule
\headrow \thd{Estado} & \thd{Significado}\\
\midrule
\pronto & Funciona de ponta a ponta hoje, e há um teste automatizado que
  percorre a tela como uma pessoa faria.\\
\zebra \parcial & Funciona, mas falta uma parte declarada: um material
  (áudio, foto), um ajuste de tela, ou uma regra ainda divergente. A
  pendência está escrita na própria história.\\
\afazer & Ainda não existe. Está no documento porque apareceu como lacuna,
  risco ou evolução planejada.\\
\zebra \arquivada & Existiu e saiu do produto das escolas, mas ficou GUARDADA:
  outra missão ocupou o lugar dela na trilha, e ela não abre para ninguém. Só a
  equipe da Unooria vê a lista, e os dados das turmas continuam guardados.\\
\bottomrule
\end{tabularx}

\section{O que este documento \emph{não} é}

\begin{itemize}
  \item \textbf{Não é especificação técnica.} Não há aqui nomes de tabela,
    endereços internos nem desenho de sistema. Quem precisar desse nível
    encontra tudo na documentação técnica (PU-DOC-001); o
    Apêndice~\ref{ap:rastreio} diz em que capítulo dela cada épico mora.
  \item \textbf{Não é plano de lançamento.} A ordem dos épicos segue o caminho
    de uso da plataforma, não prioridade. A prioridade está no
    Capítulo~\ref{cap:backlog}.
  \item \textbf{Não é definitivo.} É a versão 0.1: o retrato do que existe,
    feito para ser corrigido pelo PO e mantido junto com o produto.
\end{itemize}

\begin{nota}
\textbf{Como manter vivo.} Uma história nova entra com o próximo número da
feature; uma história que muda de estado troca o selo. Nenhum identificador é
reaproveitado: se uma história for descartada, ela sai do documento e o número
fica vago. Assim, um \id{H04.2.3} citado num card ou numa conversa sempre aponta
para a mesma coisa.
\end{nota}

% ===========================================================================
%  2 — O produto
% ===========================================================================
\chapter{O produto}
\label{cap:produto}

\section{O problema}

Crianças e adolescentes vivem dentro de redes sociais, vídeos curtos,
aplicativos de mensagem e assistentes de inteligência artificial --- e a
escola quase nunca consegue falar sobre isso \emph{de dentro}. Falar de
\enquote{fake news} numa aula expositiva é fácil; fazer a turma
\emph{sentir} como um boato viaja, como um feed decide o que cada um vê, ou
por que um vídeo com muitos inscritos não é necessariamente verdadeiro, é
difícil. E o professor da rede pública raramente tem tempo para preparar
esse tipo de aula do zero.

Ao mesmo tempo, as redes públicas precisam cumprir as habilidades da BNCC de
Computação, planejar o ano letivo, cadastrar milhares de pessoas e prestar
contas do que acontece nas escolas --- sem expor dado de criança.

\section{A proposta}

A Plataforma Unooria é uma plataforma educacional vendida para \textbf{redes
públicas de ensino} (secretarias estaduais e municipais). Ela entrega:

\begin{itemize}
  \item \textbf{Missões imersivas para os estudantes}: atividades digitais que
    imitam os aplicativos que eles conhecem (uma rede social de fotos, um
    feed de vídeos curtos, um grupo de mensagens, um assistente de IA, um jogo
    em 3D). O estudante age como agiria no próprio celular; depois, a missão
    revela o que estava acontecendo por trás, e a turma discute.
  \item \textbf{Ferramentas para o professor}: preparar e conduzir a aula,
    gerar plano de aula e planejamento de ensino alinhados à BNCC, e cuidar do
    próprio bem-estar.
  \item \textbf{Gestão da rede}: cadastrar escolas, pessoas e turmas,
    acompanhar a adoção e controlar o que cada rede contratou.
\end{itemize}

\section{Quem compra e quem usa}

O cliente \textbf{não é a escola nem a família}: é a rede. Uma rede contrata a
plataforma e recebe uma instalação própria, separada das instalações de outras
redes. Dentro dela, cem ou mais escolas compartilham a mesma plataforma, e
cada escola só enxerga os próprios dados.

\begin{tabularx}{\linewidth}{@{}L{46mm}X@{}}
\toprule
\headrow \thd{Grandeza, por rede} & \thd{Alvo}\\
\midrule
Estudantes cadastrados & cerca de \num{10000}\\
\zebra Pessoas ao mesmo tempo no horário de pico & \num{2000} a \num{4000}, entre
  7h50 e 8h10 de um dia letivo\\
Escolas & 100 ou mais\\
\zebra Faixa atendida & do 1º ano do Ensino Fundamental à 3ª série do Ensino
  Médio\\
Missões no catálogo alvo & 48 (12 anos $\times$ 4 unidades)\\
\zebra Entrada em produção & novembro de 2026\\
\bottomrule
\end{tabularx}

\section{Os quatro produtos}

O que a rede contrata são \textbf{produtos}. Cada produto aparece como um
cartão no início da área do professor, e a Unooria pode ligar ou desligar cada
um por rede, regional ou escola.

\begin{tabularx}{\linewidth}{@{}L{40mm}L{30mm}X@{}}
\toprule
\headrow \thd{Produto (nome na tela)} & \thd{Chamada} & \thd{O que é}\\
\midrule
\textbf{Entre Redes} & Letramento digital & A trilha de missões: 4 unidades
  por ano, uma missão por unidade, do 1º ano EF à 3ª série EM, mais missões
  \enquote{para qualquer ano}. Hoje: 8 das 48 unidades publicadas, mais a
  Vila do Lume.\\
\zebra \textbf{Gerador de Plano de Aula} & Autoria & Gera um plano anual a
  partir do currículo e do material didático que a rede já adota.\\
\textbf{Planejamento de Ensino} & Currículo e planejamento & Parte da BNCC e
  desce do plano anual até o plano de cada semana, com as datas reais do
  calendário.\\
\zebra \textbf{Como você está?} & Bem-estar & Um acompanhamento curto e
  privado do bem-estar do próprio professor.\\
\bottomrule
\end{tabularx}

\begin{nota}
\textbf{Instanooria, UnooTube e unooGPT não são produtos.} São os nomes dos
aplicativos que as missões imitam (os \enquote{clones}), e vivem todos dentro
do Entre Redes. O estudante nunca vê a palavra \enquote{produto}: ele vê só as
missões da turma dele.
\end{nota}

\section{O caminho de uma aula}
\label{sec:caminho}

A plataforma é um encadeamento de responsabilidades. Cada pessoa faz a sua
parte, e a parte seguinte só acontece se a anterior foi feita. A
Figura~\ref{fig:caminho} mostra o caminho completo de uma missão, da
contratação até o relatório da diretoria.

\begin{figure}[H]
\centering
\fitwidth{%
\begin{tikzpicture}[
  node distance=5mm and 6mm,
  passo/.style={draw=unooBlue, fill=unooBlueWash, rounded corners=2pt,
    minimum width=38mm, minimum height=19mm, align=center, font=\small,
    text width=36mm},
  quem/.style={font=\scriptsize\bfseries, text=unooLavenderSoft},
  liga/.style={-{Stealth[length=2.4mm]}, thick, draw=unooNavy}
]
  \node[passo] (p1) {\textbf{1. Disponibilizar}\\o produto e as missões\\existem nesta rede};
  \node[passo, right=of p1] (p2) {\textbf{2. Cadastrar}\\escolas, pessoas\\e turmas};
  \node[passo, right=of p2] (p3) {\textbf{3. Preparar e liberar}\\a missão para\\a turma};
  \node[passo, below=14mm of p3] (p4) {\textbf{4. Jogar}\\a missão, no ritmo\\da aula};
  \node[passo, left=of p4] (p5) {\textbf{5. Conduzir e ler}\\a turma: debate,\\revelação, resultados};
  \node[passo, left=of p5] (p6) {\textbf{6. Acompanhar}\\a adoção, sem\\identificar criança};

  \node[quem, above=1mm of p1] {Administrador Unooria};
  \node[quem, above=1mm of p2] {Coordenação / Adm. do contrato};
  \node[quem, above=1mm of p3] {Professor};
  \node[quem, below=1mm of p4] {Estudante};
  \node[quem, below=1mm of p5] {Professor};
  \node[quem, below=1mm of p6] {Coordenação regional / Adm. do contrato};

  \draw[liga] (p1) -- (p2);
  \draw[liga] (p2) -- (p3);
  \draw[liga] (p3) -- (p4);
  \draw[liga] (p4) -- (p5);
  \draw[liga] (p5) -- (p6);
\end{tikzpicture}}
\caption[O caminho de uma aula]{O caminho de uma aula. Cada passo depende do
  anterior: sem o produto disponível a missão nem aparece; sem turma o
  estudante não vê missão nenhuma; sem liberação a missão fica com cadeado.}
\label{fig:caminho}
\end{figure}

\begin{atencao}
\textbf{Disponível não é liberado.} São duas perguntas diferentes, feitas por
pessoas diferentes, e confundir as duas é o erro mais caro de suporte:
\begin{itemize}
  \item \textbf{Disponibilidade} --- \enquote{este recurso existe NESTA REDE?}
    Quem responde é só a Unooria, no painel dela. Por padrão, tudo está
    ligado; desligar é exceção (não contratado, ou problema numa versão).
  \item \textbf{Liberação} --- \enquote{esta missão está aberta NESTA TURMA?}
    Quem responde é o professor, na sala. Por padrão, nada está aberto.
\end{itemize}
Um produto desligado tem, para quem está na escola, o mesmo sintoma de um
estudante mal enturmado: a missão não aparece.
\end{atencao}

\section{As pessoas}
\label{sec:pessoas}

São seis papéis com conta, mais o visitante do site. Os nomes abaixo são os
que a interface mostra; nas telas de gestão diz-se sempre
\enquote{Coordenação}, nunca \enquote{gestor}.

\begin{tabularx}{\linewidth}{@{}L{36mm}X X@{}}
\toprule
\headrow \thd{Papel (nome na tela)} & \thd{Quem é} & \thd{O que mais precisa}\\
\midrule
\textbf{Estudante} & Criança ou adolescente da rede, do 1º ano EF à 3ª série
  EM. É titular de dado de menor. & Entrar sem atrito, achar a missão do dia e
  jogar sem perder o que fez.\\
\zebra \textbf{Professor} & Dá aula em uma ou mais turmas, às vezes em mais de
  uma escola. & Preparar a aula rápido, conduzir o momento coletivo, ler a
  turma e planejar o ano.\\
\textbf{Coordenação da escola} & A secretaria ou coordenação pedagógica de uma
  escola. & Cadastrar e manter pessoas e turmas da escola.\\
\zebra \textbf{Coordenação regional} & Diretoria de ensino, polo ou secretaria
  municipal que agrupa escolas. & Ver as escolas da regional, trabalhar numa
  por vez e acompanhar a adoção.\\
\textbf{Administrador do contrato} & A equipe da secretaria que contratou a
  plataforma. & Organizar a rede inteira e ler o agregado, sem entrar na sala
  de aula.\\
\zebra \textbf{Administrador Unooria} & A equipe da Unooria. & Decidir o que
  existe em cada rede, controlar versões e limites de IA, receber
  interessados.\\
\textbf{Visitante} & Alguém de uma rede que ainda não é cliente. & Entender a
  proposta e deixar contato.\\
\bottomrule
\end{tabularx}

\medskip
\textbf{Quem cadastra quem} segue uma cadeia, sempre de cima para baixo:
Administrador Unooria \seta{} Administrador do contrato \seta{} Coordenação
regional \seta{} Coordenação da escola \seta{} Professor e Estudante. Cada
nível cria os que estão abaixo dele, nunca o próprio nível nem um acima.

\section{As nossas promessas}
\label{sec:promessas}

Algumas regras não são detalhe de tela: são o que a plataforma promete à rede,
à escola e à família. Elas valem para todas as missões que existem e para todas
as que vierem, e toda história deste documento respeita todas elas.

\begin{enumerate}
  \item \textbf{A missão só abre na turma em que o professor a liberou.}
    Esconder o cartão não basta: digitar o endereço da missão também não abre.
  \item \textbf{A rede controla o que existe nela --- e desligar nunca apaga.}
    Tudo pode ser desligado por rede, regional ou escola, com motivo e com um
    histórico que a própria rede consegue ler.
  \item \textbf{Nota e correção são da plataforma, nunca do aparelho do
    estudante.} A resposta certa não chega ao tablet antes da hora.
  \item \textbf{Não existe ranking de crianças.} Nenhuma tela de professor ou
    de gestão ordena estudantes por acerto.
  \item \textbf{O relatório não rotula criança.} Ele não afirma o que não mediu:
    \enquote{não teve a oportunidade} nunca é lido como \enquote{escolheu não},
    e um valor com uma única medida não recebe número.
  \item \textbf{Dado de criança é protegido.} Uma escola nunca vê outra; o que
    a criança escreve em texto livre é lido pelo professor dela, não pela
    gestão; os registros de uso não carregam nome.
  \item \textbf{A IA tem limite e tem garantia.} Nenhuma aula depende da IA
    estar no ar; nada gerado chega à turma sem o professor aprovar; a IA nunca
    corrige e nunca recebe dado pessoal; cada pessoa tem um teto diário.
  \item \textbf{O bem-estar do professor é só dele.} Não existe tela de gestão,
    agregado nem exportação desse dado.
  \item \textbf{Tudo em português do Brasil}, com datas no formato brasileiro
    --- exceto dentro dos aplicativos imitados, que falam como o aplicativo
    real (\enquote{Feed}, \enquote{Shorts}, \enquote{Stories}).
\end{enumerate}

\section{Glossário de produto}
\label{sec:glossario}

\begin{tabularx}{\linewidth}{@{}L{34mm}X@{}}
\toprule
\headrow \thd{Termo} & \thd{Significado}\\
\midrule
\textbf{Rede} / contrato & A secretaria cliente. Um contrato é uma rede, com
  instalação própria.\\
\zebra \textbf{Regional} & O nível que agrupa escolas dentro da rede
  (diretoria de ensino, polo, secretaria municipal). Toda escola pertence a
  uma.\\
\textbf{Escola} & A unidade. Tudo que é dado de estudante pertence a uma
  escola.\\
\zebra \textbf{Vínculo} & A ligação de uma pessoa com uma escola, com papel e
  data. Um professor pode ter vários; um estudante tem um ativo por vez.\\
\textbf{Turma} & A sala. O estudante tem uma \textbf{turma de referência}
  (a principal), que é a que as missões usam.\\
\zebra \textbf{Enturmamento} & O ato de pôr pessoas em turmas.\\
\textbf{Ano letivo} / período & Ano, semestre, bimestre ou trimestre. O
  \enquote{8A de 2026} e o \enquote{8A de 2027} são turmas diferentes, o que
  preserva a história.\\
\zebra \textbf{Missão} & O que o estudante joga, visto do ponto de vista da
  aula. É a palavra que professor e estudante leem na tela.\\
\textbf{Unidade} & Cada ano tem 4 unidades; cada unidade tem uma missão.\\
\zebra \textbf{Trilha} & O conjunto das missões do Entre Redes, organizado
  por ano e unidade.\\
\textbf{Clone} / disfarce & A tela que imita um aplicativo real. Dentro dela,
  tudo fala a língua do aplicativo imitado.\\
\zebra \textbf{Revelação} & O momento em que a missão deixa de fingir ser o
  aplicativo e passa a ser a escola falando: a explicação, o \enquote{Meu
  percurso}, o debate.\\
\textbf{Ficha da missão} & A página que o professor lê antes: como a missão
  funciona e o que ela \emph{não} conta ao estudante.\\
\zebra \textbf{Guia da aula} & O roteiro de condução: minuto a minuto,
  perguntas de debate, mediação, avaliação e habilidade da BNCC.\\
\textbf{Liberar} / encerrar & Abrir e fechar a missão para uma turma.\\
\zebra \textbf{Zerar a turma} & Apagar os dados de uma missão numa turma, para
  repeti-la. Não toca outras missões.\\
\textbf{Disponibilidade} & Se um produto ou missão existe numa rede, regional
  ou escola. Quem decide é a Unooria.\\
\zebra \textbf{Limite diário de IA} & O quanto cada pessoa pode usar de IA por
  dia, por perfil (estudante ou professor). Na tela aparece como
  \enquote{Tokens por dia}.\\
\textbf{Perfil portátil} & O documento assinado que o estudante leva ao sair
  da rede, a pedido da família.\\
\zebra \textbf{Primeiro acesso} & A primeira entrada, em que a pessoa troca a
  senha provisória pela dela.\\
\textbf{Equipe} & Quem divide um computador numa missão: dois ou três
  estudantes diante da mesma tela, com o login de um deles, que representa o
  grupo. Vale para uma missão numa turma; não é grupo fixo da turma.\\
\zebra \textbf{Atividade} & Uma tarefa dentro de uma missão, definida pela
  própria missão --- não existe, por ora, uma atividade comum a todas. Hoje só o
  Instanooria tem atividades.\\
\textbf{Rodada} & Cada vez que uma atividade é feita. Reiniciar a atividade
  abre uma rodada nova, e a anterior fica guardada; o resultado é a melhor
  rodada de cada habilidade.\\
\zebra \textbf{Régua pedagógica} / ficha pedagógica & O que uma missão mede:
  quais habilidades da BNCC, em quais momentos, qual pergunta e quanto vale
  cada alternativa. Escrita e validada pela equipe pedagógica.\\
\textbf{Evidência} & O nível que um estudante demonstrou numa habilidade da
  BNCC, numa atividade, numa rodada: atingiu, em processo ou não atingiu. Sem
  registro, é \enquote{Não avaliado}, que nunca conta como não atingiu.\\
\bottomrule
\end{tabularx}

% ===========================================================================
%  3 — Mapa dos épicos
% ===========================================================================
\chapter{Mapa dos épicos}
\label{cap:mapa}

São \textbf{14 épicos, 75 features e 229 histórias}. Os épicos seguem o
caminho de uso da plataforma (Figura~\ref{fig:caminho}): da entrada à
organização da rede, da sala de aula à missão, dos resultados às ferramentas,
as garantias que valem para tudo, a sala em que os estudantes dividem o
computador e, por fim, o que cada estudante demonstrou de cada habilidade da
BNCC.

\begin{destaque}
\textbf{Em 16 de outubro de 2026:} \textbf{155} histórias prontas
(\textasciitilde67\%), \textbf{57} parciais, \textbf{12} a fazer e \textbf{6}
arquivadas (as três do UnooTube, que saiu da trilha sem ser apagado, as perguntas do fim do Unoovista, as do resultado da Missão Tempestade, e a rodada
de desvio do Júri Estadual). Nove das
parciais são do agente de IA do professor (\id{F08.4} e \id{H10.1.4}),
construído e ainda não no ar, e oito são do Júri Estadual de Projetos
(\id{F06.14}, com o podcast dele e as imagens da mostra pronta), construído com
os quarenta estandes e com o teste de ponta a ponta à espera de ser atualizado
e corrido. Quatro são da Ajuda com atendente (\id{F01.6} e
\id{H10.1.5}), construída e em teste interno, desligada nas redes. Mais da
metade do que é parcial é a evidência pedagógica (\id{E14}) e as três missões
com régua, a versão nova de Quem mandou isso? (\id{F06.11}), A Seleção
(\id{F06.12}, com o podcast dela) e o Jornal da Escola (\id{H06.10.4}) ---
construídas e percorridas pelo teste de ponta a ponta; a primeira à espera de o
PO revisar os textos, a segunda à espera da revisão da voz e das imagens
geradas, da conferência dos níveis que a IA dá à legenda, da revisão dos
textos e da foto da Amara pelo PO, e a terceira à espera das imagens do banco padrão, da revisão
do critério do livro do 4º ano, da menção à IA dela nos Termos de Uso e da
primeira corrida verde do teste da aprovação das imagens (\id{H06.10.6}). O resto
é, em quase todos os casos, material que falta (áudio, foto, revisão
pedagógica) ou ajuste de tela. O que é a fazer concentra-se em recuperação de
senha, prazos de guarda de dados, acessibilidade, nas equipes do Unoovista e
nas missões que faltam. Oito itens --- de teste de carga a contrato --- bloqueiam a
entrada em produção (Capítulo~\ref{cap:backlog}).
\end{destaque}

\section{Os catorze épicos}

\begin{tabularx}{\linewidth}{@{}L{10mm}X L{12mm}L{9mm}L{9mm}L{9mm}L{17mm}@{}}
\toprule
\headrow \thd{ID} & \thd{Épico} & \thd{Feat.} & \thd{Pr.} & \thd{Parc.} & \thd{A f.} & \thd{Estado}\\
\midrule
\endhead
\id{E01} & Entrar e circular na plataforma (p.~\pageref{ep:e01}) & 6 & 10 & 3 & 2 & \parcial\\
\zebra \id{E02} & Organizar a rede (p.~\pageref{ep:e02}) & 4 & 7 & 0 & 1 & \parcial\\
\id{E03} & Cadastrar pessoas e turmas (p.~\pageref{ep:e03}) & 8 & 16 & 1 & 0 & \pronto\\
\zebra \id{E04} & Sala de aula: o professor abre a missão (p.~\pageref{ep:e04}) & 4 & 8 & 0 & 1 & \pronto\\
\id{E05} & A experiência do estudante na plataforma (p.~\pageref{ep:e05}) & 2 & 6 & 0 & 0 & \pronto\\
\zebra \id{E06} & As missões do Entre Redes (p.~\pageref{ep:e06}) & 12 & 43 & 28 & 1 & \pronto\\
\id{E07} & Acompanhar a aprendizagem e a rede (p.~\pageref{ep:e07}) & 3 & 6 & 1 & 1 & \parcial\\
\zebra \id{E08} & Ferramentas do professor (p.~\pageref{ep:e08}) & 4 & 10 & 8 & 0 & \pronto\\
\id{E09} & Inteligência artificial com limites (p.~\pageref{ep:e09}) & 4 & 5 & 0 & 2 & \parcial\\
\zebra \id{E10} & Painel da Unooria (p.~\pageref{ep:e10}) & 7 & 15 & 2 & 1 & \pronto\\
\id{E11} & Site público e captação (p.~\pageref{ep:e11}) & 3 & 3 & 1 & 0 & \parcial\\
\zebra \id{E12} & Confiança: privacidade, segurança e operação (p.~\pageref{ep:e12}) & 7 & 3 & 3 & 1 & \parcial\\
\id{E13} & Equipes na missão (p.~\pageref{ep:e13}) & 4 & 13 & 0 & 2 & \parcial\\
\zebra \id{E14} & Evidência pedagógica (p.~\pageref{ep:e14}) & 7 & 10 & 10 & 0 & \parcial\\
\midrule
 & \textbf{Total} & \textbf{75} & \textbf{155} & \textbf{57} & \textbf{12} & \\
\bottomrule
\end{tabularx}

\medskip
{\small\color{unooMuted}Pr. = prontas; Parc. = parciais; A f. = a fazer. As 6
histórias \arquivada{} (o UnooTube, \id{F06.2}, as perguntas do fim do Unoovista, \id{H06.6.8}, as do resultado da Missão Tempestade, \id{H06.8.5}, e a rodada de desvio do Júri,
\id{H06.14.8}) não entram nas colunas. O
estado do épico é \pronto{} quando o que falta é ajuste ou evolução e não
impede o uso; é \parcial{} quando falta uma capacidade que alguém vai sentir
falta.}

\section{Quem aparece em cada épico}

\begin{tabularx}{\linewidth}{@{}L{10mm}X *{7}{>{\centering\arraybackslash}p{9mm}}@{}}
\toprule
\headrow \thd{ID} & \thd{Épico} & \thd{Est.} & \thd{Prof.} & \thd{C.\,esc.} & \thd{C.\,reg.} & \thd{Adm.\,c.} & \thd{Uno.} & \thd{Vis.}\\
\midrule
\id{E01} & Entrar e circular & \yes & \yes & \yes & \yes & \yes & \yes & \no\\
\zebra \id{E02} & Organizar a rede & \no & \no & \yes & \yes & \yes & \yes & \no\\
\id{E03} & Pessoas e turmas & \no & \no & \yes & \yes & \yes & \yes & \no\\
\zebra \id{E04} & Sala de aula & \no & \yes & \yes & \no & \no & \no & \no\\
\id{E05} & Experiência do estudante & \yes & \no & \no & \no & \no & \no & \no\\
\zebra \id{E06} & Missões & \yes & \yes & \no & \no & \no & \no & \no\\
\id{E07} & Acompanhamento & \no & \yes & \no & \yes & \yes & \yes & \no\\
\zebra \id{E08} & Ferramentas do professor & \no & \yes & \no & \no & \no & \no & \no\\
\id{E09} & IA com limites & \yes & \yes & \no & \no & \no & \yes & \no\\
\zebra \id{E10} & Painel da Unooria & \no & \no & \no & \no & \yes & \yes & \no\\
\id{E11} & Site público & \no & \no & \no & \no & \no & \yes & \yes\\
\zebra \id{E12} & Confiança & \yes & \yes & \yes & \yes & \yes & \yes & \no\\
\id{E13} & Equipes na missão & \yes & \yes & \yes & \yes & \yes & \yes & \no\\
\zebra \id{E14} & Evidência pedagógica & \yes & \yes & \yes & \yes & \yes & \yes & \no\\
\bottomrule
\end{tabularx}

\medskip
{\small\color{unooMuted}Est. = Estudante; Prof. = Professor; C.\,esc. =
Coordenação da escola; C.\,reg. = Coordenação regional; Adm.\,c. =
Administrador do contrato; Uno. = Administrador Unooria; Vis. = Visitante. Em
\id{E10}, o Administrador do contrato só lê o histórico; em \id{E13}, as
coordenações e as administrações só consultam as equipes; em \id{E14}, elas só
leem, e o estudante responde sem ver o próprio resultado.}

\section{O que cada papel pode, em resumo}

\begin{tabularx}{\linewidth}{@{}X *{6}{>{\centering\arraybackslash}p{10.5mm}}@{}}
\toprule
\headrow \thd{Capacidade} & \thd{Est.} & \thd{Prof.} & \thd{C.\,esc.} & \thd{C.\,reg.} & \thd{Adm.\,c.} & \thd{Uno.}\\
\midrule
\endhead
Jogar as missões liberadas & \yes & \no & \no & \no & \no & \no\\
\zebra Liberar e encerrar missão na turma & \no & \yes & vê & vê & vê & vê\\
Montar as equipes de uma missão na turma & \no & \yes & vê & vê & vê & vê\\
\zebra Confirmar e ajustar a evidência por habilidade da BNCC & \no & \yes & vê & vê & vê & vê\\
Ler os resultados da turma e de cada estudante & o seu & \yes & \no & \no & \no & \no\\
\zebra Ler o agregado (sem identificar criança) & \no & \no & \no & \yes & \yes & \yes\\
Plano de aula, planejamento, bem-estar, agente de IA & \no & \yes & \no & \no & \no & \no\\
\zebra Cadastrar pessoas, importar, resetar senha & \no & \no & \yes & \yes & \yes & \yes\\
Turmas, enturmamento, ano letivo & \no & lê & \yes & \yes & \yes & \yes\\
\zebra Perfil portátil & \no & \no & \yes & \yes & \yes & \yes\\
Transferir estudante de escola & \no & \no & \no & \no & \yes & \yes\\
\zebra Regionais, escolas, vínculos em lote & \no & \no & \no & \no & \yes & \yes\\
Criar Coordenação regional & \no & \no & \no & \no & \yes & \yes\\
\zebra Criar Administrador do contrato & \no & \no & \no & \no & \no & \yes\\
Ligar e desligar recursos, versões, limites de IA & \no & \no & \no & \no & lê & \yes\\
\zebra Prospecção & \no & \no & \no & \no & \no & \yes\\
Tirar dúvidas com o atendente da Ajuda & \no & \yes & \yes & \yes & \yes & \yes\\
\bottomrule
\end{tabularx}

\medskip
{\small\color{unooMuted}Ausências deliberadas: o bem-estar do professor não
tem nenhuma visão de gestão; a gestão não tem acesso à linha de uma criança
nas missões --- a exceção declarada é a evidência pedagógica, em que ela lê o
nível de cada habilidade, nunca texto (\id{H14.4.4}, \id{F14.7}); o Administrador do
contrato lê a rede em agregado, nunca a sala de aula --- e, no painel de
habilidades, desce da rede ao estudante como as coordenações (\id{H14.7.5}); e o
estudante não vê a própria evidência pedagógica.}

% ===========================================================================
%  E01 — Entrar e circular na plataforma
% ===========================================================================
\chapter{E01 \textperiodcentered{} Entrar e circular na plataforma}
\label{ep:e01}

\begin{fichaepico}
\begin{campos}
  \item[Objetivo] Cada pessoa entra com uma conta só, cai direto na área do
    próprio papel e circula sem se perder --- inclusive quem trabalha em mais
    de uma escola.
  \item[Para quem] Todos os seis papéis com conta.
  \item[Valor] Às 8h da manhã, milhares de crianças entram ao mesmo tempo. Cada
    segundo de atrito na entrada é tempo de aula perdido.
  \item[Estado] \parcial{} --- entrar, primeiro acesso e navegação estão
    prontos; faltam a recuperação de senha e o login único da rede, e a
    Ajuda com atendente está construída, em teste interno.
\end{campos}
\end{fichaepico}

\begin{tabularx}{\linewidth}{@{}L{16mm}X L{20mm}@{}}
\toprule
\headrow \thd{ID} & \thd{Feature} & \thd{Estado}\\
\midrule
\id{F01.1} & Entrar com usuário e senha & \pronto\\
\zebra \id{F01.2} & Cada pessoa na sua área & \pronto\\
\id{F01.3} & Trabalhar em mais de uma escola & \pronto\\
\zebra \id{F01.4} & Cuidar da própria senha & \parcial\\
\id{F01.5} & Entrar pelo login único da rede & \afazer\\
\zebra \id{F01.6} & Tirar dúvidas com o atendente da Ajuda & \parcial\\
\bottomrule
\end{tabularx}

% ---------------------------------------------------------------------------
\feature{F01.1}{Entrar com usuário e senha}{\pronto}

Uma única tela de entrada serve a todos os papéis. É o papel da pessoa que
decide para onde ela vai depois.

\begin{historia}{H01.1.1}{Entrar e cair na minha área}{\pronto}
\hu{qualquer pessoa com conta na rede}{entrar com usuário e senha e cair
direto na minha área}{não ter de procurar por onde começar}
\aceite
\begin{itemize}
  \item A tela se chama \enquote{Entrar na plataforma}, com os campos Usuário
    e Senha e o botão \enquote{Entrar}.
  \item Depois de entrar: o Estudante vai para \enquote{Minhas missões}; o
    Professor para o Início dele; as três coordenações e o Administrador do
    contrato para \enquote{Usuários}; o Administrador Unooria para
    \enquote{Recursos}.
  \item Senha errada, usuário inexistente e conta desativada recebem
    \emph{a mesma} mensagem: \enquote{Usuário ou senha incorretos.} --- a tela
    nunca diz qual das duas metades errou, para não confirmar que um usuário
    existe.
  \item Depois de muitas tentativas seguidas, aparece \enquote{Muitas
    tentativas. Aguarde um momento e tente de novo.}
  \item Se a plataforma estiver com problema, aparece \enquote{Não foi possível
    entrar agora. Tente de novo em alguns instantes.}
\end{itemize}
\end{historia}

\begin{historia}{H01.1.2}{Criar minha senha no primeiro acesso}{\pronto}
\hu{pessoa que recebeu uma senha provisória (da importação ou de um
reajuste)}{ser levada a criar a minha própria senha antes de qualquer
coisa}{que só eu conheça a minha senha}
\aceite
\begin{itemize}
  \item A tela \enquote{Crie sua senha} aparece logo depois de entrar, sem menu
    e sem saída.
  \item Campos: \enquote{Senha atual}, \enquote{Nova senha} (pelo menos 8
    caracteres, diferente da atual) e \enquote{Repita a nova senha}.
  \item Enquanto a senha não for trocada, nada funciona: nem as telas nem as
    missões. Digitar o endereço de uma missão não pula esta etapa.
  \item Vale para todos os papéis, inclusive o Administrador Unooria.
  \item Trocada a senha, a pessoa segue para a área dela normalmente.
\end{itemize}
\end{historia}

\begin{historia}{H01.1.3}{Não cair no meio da aula}{\pronto}
\hu{Estudante}{que a minha sessão não expire enquanto estou numa
missão}{não perder a aula tendo de entrar de novo}
\aceite
\begin{itemize}
  \item A sessão expira depois de 2 horas \textbf{sem uso}, igual para todos os
    papéis.
  \item Ficar dentro de uma missão conta como uso.
\end{itemize}
\end{historia}

\begin{historia}{H01.1.4}{Sair}{\pronto}
\hu{qualquer pessoa}{sair da plataforma}{que ninguém use a minha conta num
computador compartilhado da escola}
\aceite
\begin{itemize}
  \item O menu do avatar, no cabeçalho, tem o item \enquote{Sair}.
  \item Depois de sair, voltar a qualquer tela interna leva ao login.
\end{itemize}
\end{historia}

% ---------------------------------------------------------------------------
\feature{F01.2}{Cada pessoa na sua área}{\pronto}

\begin{historia}{H01.2.1}{Voltar para a minha área sem ver erro}{\pronto}
\hu{Estudante que digitou o endereço da área do professor}{ser levado de
volta à minha área, sem mensagem de erro}{não ficar perdido nem assustado}
\aceite
\begin{itemize}
  \item Quem entra no endereço de outra área volta para a própria, sem aviso de
    \enquote{sem permissão}.
  \item A proteção não depende disso: os dados de outra área também são
    recusados, mesmo que alguém tente por outro caminho.
\end{itemize}
\end{historia}

\begin{historia}{H01.2.2}{Saber onde estou}{\pronto}
\hu{qualquer pessoa}{ver onde estou e voltar um nível com um clique}{não me
perder numa plataforma com muitas telas}
\aceite
\begin{itemize}
  \item O cabeçalho é o mesmo em todas as áreas; a marca leva ao início do
    papel.
  \item Toda tela interna tem a trilha de navegação (\enquote{Início \seta{}
    Entre Redes \seta{} Missão\ldots}), exceto o início de cada papel.
  \item Dentro de uma missão, o cabeçalho mostra \enquote{\ensuremath{\leftarrow}
    Missões}.
  \item O cabeçalho mostra nome, uma linha de apoio (a série do estudante ou o
    papel) e um avatar com as iniciais.
  \item As telas de gestão têm menu lateral filtrado pelo papel.
  \item Os menus laterais da gestão e da Unooria mostram só o que leva a algum
    lugar: o item \enquote{Iniciativas}, que não fazia nada, saiu.
\end{itemize}
\end{historia}

% ---------------------------------------------------------------------------
\feature{F01.3}{Trabalhar em mais de uma escola}{\pronto}

\begin{historia}{H01.3.1}{Trocar a escola ativa}{\pronto}
\hu{Professor ou Coordenação com vínculo em mais de uma escola}{trocar a
escola em que estou trabalhando pelo cabeçalho}{ver e registrar só o que é
daquela escola}
\aceite
\begin{itemize}
  \item Com vínculo em 2 ou mais escolas aparece o seletor \enquote{Escola
    ativa}; com uma escola só, ele não aparece.
  \item Trocar de escola muda tudo o que a pessoa vê, sem recarregar a página.
  \item As ferramentas abertas (plano de aula, planejamento) recomeçam na
    escola nova, para nada ser gravado na escola errada.
  \item No celular o seletor continua no cabeçalho, mais estreito; tocá-lo
    abre a lista do próprio telefone, com o nome inteiro de cada escola.
\end{itemize}
\end{historia}

% ---------------------------------------------------------------------------
\feature{F01.4}{Cuidar da própria senha}{\parcial}

Hoje a senha só muda quando uma coordenação usa \enquote{Resetar senha}
(\id{H03.3.4}). O item \enquote{Trocar senha} ficou fora do menu de propósito,
para não ser um link que não leva a nada.

\begin{historia}{H01.4.1}{Trocar minha senha quando eu quiser}{\pronto}
\hu{qualquer pessoa}{trocar a minha senha a qualquer momento}{me proteger se
alguém descobriu a senha}
\aceite
\begin{itemize}
  \item O menu do avatar tem \enquote{Trocar senha}, com as mesmas regras do
    primeiro acesso: a senha atual, pelo menos 8 caracteres e diferente da
    atual.
  \item A tela se chama \enquote{Trocar senha} e tem \enquote{Cancelar}, que
    volta para a área da pessoa --- ela não fica presa ali, como no primeiro
    acesso.
  \item Depois de salvar, a tela confirma \enquote{Senha trocada} e
    \enquote{Voltar} leva de volta à área da pessoa. Na próxima entrada vale a
    senha nova.
\end{itemize}
\end{historia}

\begin{historia}{H01.4.2}{Recuperar a senha esquecida}{\afazer}
\hu{Professor ou Coordenação}{recuperar a minha senha sem depender de
alguém}{não ficar sem acesso no dia da aula}
\aceiteproposta
\begin{itemize}
  \item Na tela de entrada, \enquote{Esqueci minha senha} envia um link de uso
    único para o e-mail cadastrado.
  \item Para o estudante (que muitas vezes não tem e-mail), o caminho continua
    sendo a coordenação --- ver a decisão \id{D05} no
    Capítulo~\ref{cap:decisoes}.
\end{itemize}
\end{historia}

% ---------------------------------------------------------------------------
\feature{F01.5}{Entrar pelo login único da rede}{\afazer}

A rede usa um provedor de identidade central (WSO2), o mesmo dos outros
sistemas da secretaria. A plataforma já está preparada para ele: a troca é de
configuração. Falta validar num ambiente de homologação com a rede.

\begin{historia}{H01.5.1}{Usar a conta que eu já tenho na rede}{\afazer}
\hu{qualquer pessoa da rede}{entrar com a mesma conta que uso nos outros
sistemas da secretaria}{não ter mais uma senha para decorar}
\aceiteproposta
\begin{itemize}
  \item A tela de entrada leva ao login da rede; ao voltar, a pessoa cai na
    mesma área de hoje.
  \item Nada muda dentro das missões nem nas telas.
  \item O primeiro acesso com senha provisória deixa de existir para quem
    entra pelo login da rede.
\end{itemize}
\end{historia}

% ---------------------------------------------------------------------------
\feature{F01.6}{Tirar dúvidas com o atendente da Ajuda}{\parcial}

Um atendente que conhece a plataforma. A pessoa abre
\enquote{Ajuda} no menu do próprio nome, no alto da tela, e pergunta falando
ou escrevendo; o atendente responde em voz e o texto aparece embaixo dele.
São dois atendentes: o \textbf{mascote da Unooria}, um personagem 3D que fala
com a pessoa (\id{H01.6.4}), e o \textbf{atendente em vídeo}, uma pessoa
virtual de um serviço contratado. Quem atende primeiro é escolha do
Administrador Unooria (\id{H10.4.3}); hoje, o mascote.
Ele sabe o que este mapa diz que está pronto, os guias da aula e as missões
--- e mais nada: não vê turma, estudante nem resultado.

\begin{nota}
\textbf{O tempo de conversa é limitado.} Cada conversa e cada dia têm um teto
de minutos, que o Administrador Unooria define (\id{H10.4.2}). Começa em 10
minutos por conversa e 30 por dia, por pessoa. A Ajuda não gasta o limite
diário de IA de ninguém (\id{F09.4}). Os minutos são os mesmos para os dois
atendentes: trocar de um para o outro não dá tempo novo.
\end{nota}

\begin{historia}{H01.6.1}{Perguntar ao atendente sem sair da tela}{\parcial}
\hu{Professor, Coordenação ou Administrador}{abrir a Ajuda de qualquer tela e
perguntar ao atendente como fazer alguma coisa na plataforma}{resolver a
dúvida na hora, sem esperar alguém do suporte}
\aceite
\begin{itemize}
  \item \enquote{Ajuda} é o segundo item do menu do nome, para todos os papéis
    menos o Estudante. Abre uma janela por cima da tela: o vídeo do atendente
    em cima, a conversa escrita embaixo.
  \item Enquanto conecta, a janela diz \enquote{Chamando o atendente\ldots} e
    depois \enquote{Conectando o vídeo\ldots}; o atendente se apresenta e
    espera a pergunta.
  \item O atendente responde sobre o que a plataforma faz hoje, no papel de
    quem pergunta. O que não existe ainda, ele diz que não existe.
  \item Junto do campo, fixo: \enquote{Não diga nem escreva nomes de
    estudantes.} Nada da conversa é guardado: fechar a janela encerra a
    conversa e o vídeo.
  \item Se o vídeo não conectar, a janela diz que a rede da escola pode estar
    bloqueando a conexão e oferece \enquote{Tentar de novo}.
\end{itemize}
\pendencia{a Ajuda nasce desligada em toda rede, enquanto está em teste
interno: só aparece onde o Administrador Unooria a ligou (\id{H10.1.5}).
Falta o teste de ponta a ponta, e a voz do atendente fala português com
sotaque --- aceito pelo dono do produto.}
\end{historia}

\begin{historia}{H01.6.2}{Escrever quando não posso falar}{\parcial}
\hu{Professor ou Coordenação}{escrever a minha pergunta em vez de falar}{usar
a Ajuda numa sala com gente, ou num computador sem microfone}
\aceite
\begin{itemize}
  \item O microfone pode ser silenciado e religado a qualquer momento; o
    estado aparece no botão (\enquote{Silenciar microfone},
    \enquote{Ligar microfone}, \enquote{Microfone bloqueado}).
  \item Sem microfone, a janela diz \enquote{Sem microfone: escreva a sua
    pergunta abaixo.}, e a conversa segue só por escrito; o atendente responde
    em voz e por escrito do mesmo jeito.
  \item A pergunta escrita e a falada entram na mesma conversa, na ordem em
    que aconteceram.
\end{itemize}
\pendencia{a mesma de \id{H01.6.1}.}
\end{historia}

\begin{historia}{H01.6.3}{Saber quando o meu tempo de Ajuda acabou}{\parcial}
\hu{qualquer pessoa que usa a Ajuda}{saber quando a conversa acabou por
tempo, e quando o tempo do dia acabou}{não achar que o atendente travou}
\aceite
\begin{itemize}
  \item Ao chegar no limite da conversa, o atendente se despede, e a janela
    mostra \enquote{A conversa terminou}, com \enquote{Começar de novo}.
  \item Com os minutos do dia gastos, a Ajuda não abre e diz \enquote{Você já
    usou todos os minutos de Ajuda de hoje.} No dia seguinte, o tempo volta.
  \item Quando a resposta demora, o atendente avisa (\enquote{Um momento,
    estou vendo isso para você.}) em vez de ficar em silêncio; se não
    conseguir responder, pede que a pessoa repita a pergunta.
\end{itemize}
\pendencia{a mesma de \id{H01.6.1}.}
\end{historia}

\begin{historia}{H01.6.4}{Conversar com o mascote da Unooria}{\pronto}
\hu{Professor, Coordenação ou Administrador}{tirar a minha dúvida conversando
com o mascote da Unooria}{ser atendido por um personagem da própria
plataforma, que responde na hora, em voz e por escrito}
\aceite
\begin{itemize}
  \item Quando o mascote é quem atende, a janela da Ajuda mostra o mascote da
    Unooria em cima e a conversa escrita embaixo, no lugar do vídeo. Enquanto
    ele chega, a janela diz \enquote{Chamando o mascote\ldots} e depois
    \enquote{Conectando a conversa\ldots}.
  \item O mascote fala em voz alta, e a boca dele acompanha a fala. O que ele
    diz aparece por escrito na conversa, como no atendente em vídeo; o que a
    pessoa fala também.
  \item A pessoa pergunta falando ou escrevendo, do mesmo jeito que no
    atendente em vídeo: o microfone começa desligado e pode ser ligado e
    silenciado a qualquer momento; sem microfone, a conversa segue por escrito.
    Se a pessoa fala por cima, o mascote para e escuta.
  \item O mascote sabe o mesmo que o atendente em vídeo, responde no papel de
    quem pergunta e diz quando algo ainda não existe. Nada da conversa é
    guardado, e o aviso \enquote{Não diga nem escreva nomes de estudantes.}
    continua fixo junto do campo.
  \item Os minutos são os mesmos da Ajuda (\id{H01.6.3}): no fim do tempo da
    conversa, a janela mostra \enquote{A conversa terminou}, com
    \enquote{Começar de novo}.
  \item Se o Administrador Unooria deixou trocar de atendente, aparece um botão
    na janela --- \enquote{Conversar com o atendente em vídeo} ou
    \enquote{Conversar com o mascote} ---, e a conversa recomeça com o outro.
    Se o atendente escolhido não está disponível, a Ajuda abre com o outro;
    sem nenhum dos dois, diz \enquote{A Ajuda com atendente não está
    disponível agora.} e oferece \enquote{Tentar de novo}.
  \item Como toda a Ajuda, só aparece onde o Administrador Unooria a ligou
    (\id{H10.1.5}).
\end{itemize}
\end{historia}

% ===========================================================================
%  E02 — Organizar a rede
% ===========================================================================
\chapter{E02 \textperiodcentered{} Organizar a rede}
\label{ep:e02}

\begin{fichaepico}
\begin{campos}
  \item[Objetivo] A rede se organiza na plataforma do mesmo jeito que se
    organiza na vida real: regionais, escolas, profissionais em várias escolas
    e o calendário do ano letivo.
  \item[Para quem] Administrador do contrato; Coordenação regional; quem tem
    vínculo em várias escolas.
  \item[Valor] Sem a estrutura da rede não há onde pôr pessoas nem turmas, e a
    diretoria não consegue acompanhar as escolas dela.
  \item[Estado] \parcial{} --- regionais, escolas, vínculos e panorama
    prontos; a virada do ano ainda é manual.
\end{campos}
\end{fichaepico}

\begin{tabularx}{\linewidth}{@{}L{16mm}X L{20mm}@{}}
\toprule
\headrow \thd{ID} & \thd{Feature} & \thd{Estado}\\
\midrule
\id{F02.1} & Regionais & \pronto\\
\zebra \id{F02.2} & Escolas e vínculos de profissionais & \pronto\\
\id{F02.3} & Panorama das minhas escolas & \pronto\\
\zebra \id{F02.4} & Ano letivo e virada do ano & \parcial\\
\bottomrule
\end{tabularx}

% ---------------------------------------------------------------------------
\feature{F02.1}{Regionais}{\pronto}

\begin{historia}{H02.1.1}{Cadastrar e editar as regionais}{\pronto}
\hu{Administrador do contrato}{cadastrar as regionais da rede (diretoria de
ensino, polo, secretaria)}{organizar as escolas como a rede já se organiza}
\aceite
\begin{itemize}
  \item A tela \enquote{Regionais} tem busca por \enquote{Nome ou id} e, em cada
    regional, \enquote{Ver escolas}, \enquote{Editar} e \enquote{Excluir}.
  \item \enquote{Cadastrar regional} pede Código interno e Nome.
  \item Uma rede sem divisão ganha automaticamente a \enquote{Regional padrão},
    com nome editável.
  \item Toda escola pertence a uma regional, e a rede sempre tem pelo menos
    uma.
\end{itemize}
\end{historia}

\begin{historia}{H02.1.2}{Excluir uma regional sem perder escolas}{\pronto}
\hu{Administrador do contrato}{excluir uma regional que deixou de
existir}{que a estrutura na plataforma continue igual à da rede}
\aceite
\begin{itemize}
  \item \enquote{Excluir} sempre abre um diálogo de confirmação.
  \item Se a regional tem escolas, o diálogo exige dizer para onde elas vão
    (\enquote{Transferir e excluir}), numa operação só.
  \item Escola nunca é excluída: ela guarda a história de crianças.
\end{itemize}
\end{historia}

% ---------------------------------------------------------------------------
\feature{F02.2}{Escolas e vínculos de profissionais}{\pronto}

\begin{historia}{H02.2.1}{Cadastrar e editar as escolas}{\pronto}
\hu{Administrador do contrato}{cadastrar as escolas da rede com os dados
oficiais}{encontrar cada uma pelo código que a rede já usa}
\aceite
\begin{itemize}
  \item \enquote{Cadastrar escola} pede Código interno, Nome, Regional, Código
    INEP, Município e UF.
  \item A lista tem busca por \enquote{Nome, id, INEP ou município}, filtro por
    regional e as colunas Turmas, Estudantes e Professores.
  \item Em cada escola: \enquote{Ver turmas}, \enquote{Ver pessoas} e
    \enquote{Editar}.
\end{itemize}
\end{historia}

\begin{historia}{H02.2.2}{Vincular um profissional a várias escolas de uma vez}{\pronto}
\hu{Administrador do contrato}{procurar um professor ou coordenador na rede
inteira e vinculá-lo a várias escolas de uma vez}{não repetir o cadastro
escola por escola}
\aceite
\begin{itemize}
  \item Seleciono as escolas na lista e uso \enquote{Vincular professor ou
    gestor}.
  \item A busca procura na rede inteira por nome, usuário, matrícula ou CPF ---
    e o resultado não mostra o CPF.
  \item O relatório final diz o que foi criado, reaberto ou ignorado; quem já
    estava vinculado aparece marcado e não é vinculado de novo (\enquote{Nada
    a criar}).
  \item Professor e coordenação podem ter vários vínculos; o estudante tem uma
    escola ativa por vez.
\end{itemize}
\end{historia}

% ---------------------------------------------------------------------------
\feature{F02.3}{Panorama das minhas escolas}{\pronto}

\begin{historia}{H02.3.1}{Ver a minha regional de uma vez}{\pronto}
\hu{Coordenação regional}{ver as escolas da minha regional numa tela só, com
os números somados}{saber onde a implantação está atrasada}
\aceite
\begin{itemize}
  \item \enquote{Minhas escolas} mostra uma linha por escola (Turmas,
    Estudantes, Professores, \enquote{Sem 1º acesso}) e a soma da regional.
  \item \enquote{Trocar e abrir} leva para dentro de uma escola.
  \item Aparecem só as escolas da própria regional, e trabalha-se numa escola
    por vez.
\end{itemize}
\end{historia}

\begin{historia}{H02.3.2}{Ver as escolas em que trabalho}{\pronto}
\hu{Coordenação com vínculo em várias escolas}{ver uma linha por escola em que
trabalho}{escolher em qual entrar}
\aceite
\begin{itemize}
  \item Uma linha por escola, com \enquote{Abrir} ou \enquote{Trocar e abrir}.
  \item Aqui os números \emph{não} são somados: são escolas independentes.
\end{itemize}
\end{historia}

% ---------------------------------------------------------------------------
\feature{F02.4}{Ano letivo e virada do ano}{\parcial}

\begin{historia}{H02.4.1}{Manter o calendário do ano letivo}{\pronto}
\hu{Coordenação}{registrar os períodos do ano letivo}{que as turmas de um ano
não se misturem com as do ano seguinte}
\aceite
\begin{itemize}
  \item \enquote{Ano letivo} lista Período, Tipo (ano, semestre, bimestre,
    trimestre), Turmas e Situação.
  \item \enquote{Novo período letivo} pede Tipo, Título, Ano, Início e Fim.
  \item O \enquote{8A de 2026} e o \enquote{8A de 2027} são turmas diferentes:
    a história do ano anterior não é reescrita.
\end{itemize}
\pendencia{a tela diz \enquote{O calendário é do contrato, não da escola}, mas
a Coordenação da escola consegue criar período. Ver a decisão \id{D01}.}
\end{historia}

\begin{historia}{H02.4.2}{Virar o ano com ajuda}{\afazer}
\hu{Coordenação da escola}{um passo a passo para a virada do ano}{não
refazer tudo à mão em fevereiro}
\aceiteproposta
\begin{itemize}
  \item Um assistente guia: criar o novo período, encerrar as turmas do ano
    anterior e reenturmar a partir da planilha nova.
  \item Nada do ano anterior é apagado.
\end{itemize}
\pendencia{hoje a virada é manual (criar o período e reimportar). A primeira
virada real acontece em fevereiro de 2027.}
\end{historia}

% ===========================================================================
%  E03 — Cadastrar pessoas e turmas
% ===========================================================================
\chapter{E03 \textperiodcentered{} Cadastrar pessoas e turmas}
\label{ep:e03}

\begin{fichaepico}
\begin{campos}
  \item[Objetivo] Pôr todas as pessoas da rede na plataforma, nas turmas
    certas, sem digitar uma por uma --- e mantê-las em dia ao longo do ano.
  \item[Para quem] Coordenação da escola; Coordenação regional;
    Administrador do contrato.
  \item[Valor] Sem cadastro não há aula: o estudante sem turma não vê missão
    nenhuma. E cadastro errado vira dado de uma criança atribuído a outra.
  \item[Estado] \pronto{} na função, com ajustes de tela pendentes (a escola
    digitada à mão, filtros incompletos).
\end{campos}
\end{fichaepico}

\begin{tabularx}{\linewidth}{@{}L{16mm}X L{20mm}@{}}
\toprule
\headrow \thd{ID} & \thd{Feature} & \thd{Estado}\\
\midrule
\id{F03.1} & Importar pessoas por planilha & \pronto\\
\zebra \id{F03.2} & Receber os dados do sistema de gestão da rede & \pronto\\
\id{F03.3} & Balcão de usuários & \pronto\\
\zebra \id{F03.4} & Quem pode cadastrar quem & \pronto\\
\id{F03.5} & Turmas e enturmamento & \pronto\\
\zebra \id{F03.6} & Transferir estudante dentro da rede & \pronto\\
\id{F03.7} & Limpar contas sem matrícula & \pronto\\
\zebra \id{F03.8} & Perfil portátil: o estudante que sai da rede & \pronto\\
\bottomrule
\end{tabularx}

% ---------------------------------------------------------------------------
\feature{F03.1}{Importar pessoas por planilha}{\pronto}

\begin{historia}{H03.1.1}{Importar a escola inteira de uma vez}{\pronto}
\hu{Coordenação da escola}{importar estudantes, professores e coordenação por
uma planilha}{começar a usar a plataforma sem cadastrar pessoa por pessoa}
\aceite
\begin{itemize}
  \item A aba \enquote{Planilha da escola} tem três passos: \enquote{1. Baixe o
    modelo}, \enquote{2. Escolha a senha inicial}, \enquote{3. Envie a
    planilha}.
  \item Colunas: matrícula, nome, usuário, e-mail, papel, turma, série, turno.
    Só nome e papel são obrigatórios; para estudante, também turma e série.
  \item Usuário em branco é gerado a partir do nome e devolvido no relatório.
  \item Turma que ainda não existe é criada pela própria planilha.
  \item Coluna desconhecida é erro --- a planilha não é aceita pela metade sem
    aviso.
  \item O relatório mostra criados, atualizados, enturmados e os erros linha a
    linha.
  \item Aceita planilhas grandes (até 50 mil linhas).
\end{itemize}
\end{historia}

\begin{historia}{H03.1.2}{Escolher como nasce a senha}{\pronto}
\hu{Coordenação da escola}{escolher se o lote recebe uma senha única ou uma
senha por pessoa}{equilibrar praticidade e segurança conforme a idade da
turma}
\aceite
\begin{itemize}
  \item Opções: \enquote{Uma senha para todo o lote} ou \enquote{Uma senha
    diferente por pessoa (recomendado para estudantes)}.
  \item As senhas geradas aparecem uma única vez, no relatório (\enquote{não
    aparecem de novo}).
  \item A senha nunca vem na planilha.
  \item Toda senha nascida aqui é provisória: a pessoa cria a dela no primeiro
    acesso (\id{H01.1.2}).
\end{itemize}
\pendencia{com senha única no lote, um estudante consegue entrar na conta de
um colega antes do primeiro acesso dele. É risco aceito e avisado na tela; a
recomendação é senha por pessoa.}
\end{historia}

\begin{historia}{H03.1.3}{Reimportar sem bagunçar}{\pronto}
\hu{Coordenação da escola}{reenviar a planilha atualizada no meio do
ano}{corrigir dados sem prejudicar quem já está usando}
\aceite
\begin{itemize}
  \item Reimportar atualiza os dados e não troca a senha de quem já entrou.
  \item Um estudante que já existe em outra escola da rede (mesmo CPF) vira
    \textbf{transferência}; um professor, um \textbf{vínculo a mais}.
  \item Não existe \enquote{fundir contas}: isso atribuiria o trabalho de uma
    criança a outra.
\end{itemize}
\end{historia}

% ---------------------------------------------------------------------------
\feature{F03.2}{Receber os dados do sistema de gestão da rede}{\pronto}

Redes que já têm sistema de gestão escolar podem enviar os dados num padrão
internacional de troca (OneRoster 1.2), em vez de preencher a planilha.

\begin{historia}{H03.2.1}{Manter a plataforma igual ao sistema oficial}{\pronto}
\hu{Administrador do contrato de uma rede que já tem sistema de gestão}{enviar
os arquivos exportados desse sistema}{não manter dois cadastros à mão}
\aceite
\begin{itemize}
  \item A aba \enquote{OneRoster 1.2} lista os 7 arquivos na ordem, cada um com
    modelo em branco e exemplo.
  \item Modos: \enquote{Delta --- só o que mudou desde o último envio} ou
    \enquote{Bulk --- lote completo da rede}.
  \item Arquivo que falta ou vem fora de ordem é apontado no relatório.
  \item Pessoa sem papel definido entra como estudante (o menor acesso).
  \item Senha que vier no arquivo é ignorada.
  \item Uma carga incompleta nunca desativa ninguém por omissão, e
    \enquote{excluir} desativa ou encerra --- nunca apaga.
\end{itemize}
\end{historia}

% ---------------------------------------------------------------------------
\feature{F03.3}{Balcão de usuários}{\pronto}

\begin{historia}{H03.3.1}{Encontrar uma pessoa}{\parcial}
\hu{Coordenação}{achar rápido uma pessoa da escola}{resolver o problema dela
no balcão}
\aceite
\begin{itemize}
  \item \enquote{Usuários} busca por \enquote{Nome, CPF, usuário ou
    matrícula}.
  \item Filtros por Papel e por Situação (\enquote{Ativos},
    \enquote{Inativos}, \enquote{Sem primeiro acesso}).
  \item Mostra só as pessoas da escola ativa.
\end{itemize}
\pendencia{o filtro Papel oferece só Estudantes, Professores e Coordenação ---
faltam Coordenação regional e Administrador do contrato.}
\end{historia}

\begin{historia}{H03.3.2}{Cadastrar uma pessoa no balcão}{\pronto}
\hu{Coordenação da escola}{cadastrar uma pessoa que chegou no meio do
ano}{ela começar a usar hoje, sem esperar a próxima planilha}
\aceite
\begin{itemize}
  \item \enquote{Novo usuário} pede Papel, \enquote{Nome completo}, CPF,
    \enquote{Matrícula da rede}, Usuário, E-mail, Turma, Turno e
    \enquote{Senha do primeiro acesso} (\enquote{Gerar automaticamente} ou
    \enquote{Definir agora}).
  \item Usuário em branco é gerado a partir do nome; a turma é escolhida numa
    lista, nunca digitada.
  \item Ao salvar: \enquote{Acesso criado}, e a pessoa já nasce na turma.
  \item Se a pessoa já existe, o cadastro é recusado (não atualiza ninguém sem
    querer).
  \item Se o CPF já existe em outra escola, a mensagem \emph{não} revela quem é
    nem onde está, e orienta usar a importação.
\end{itemize}
\end{historia}

\begin{historia}{H03.3.3}{Corrigir, desativar e reativar}{\pronto}
\hu{Coordenação}{corrigir dados e desativar quem saiu}{manter a lista da escola
verdadeira}
\aceite
\begin{itemize}
  \item \enquote{Editar} permite corrigir Nome, E-mail e CPF.
  \item \enquote{Desativar} tira o acesso sem apagar nada; \enquote{Reativar}
    devolve.
\end{itemize}
\end{historia}

\begin{historia}{H03.3.4}{Resetar a senha de quem esqueceu}{\pronto}
\hu{Coordenação}{gerar uma senha nova para quem esqueceu a dele}{a pessoa
voltar a usar a plataforma no mesmo dia}
\aceite
\begin{itemize}
  \item \enquote{Resetar senha} mostra a senha temporária uma única vez:
    \enquote{Anote agora --- ela não aparece de novo.}
  \item No próximo acesso, a pessoa é obrigada a criar a própria senha.
\end{itemize}
\end{historia}

% ---------------------------------------------------------------------------
\feature{F03.4}{Quem pode cadastrar quem}{\pronto}

\begin{historia}{H03.4.1}{Só ver os papéis que posso atribuir}{\pronto}
\hu{Coordenação}{que a tela me ofereça só os papéis que eu posso dar}{não
tentar algo que não tenho autoridade para fazer}
\aceite
\begin{itemize}
  \item A cadeia é: Administrador Unooria \seta{} Administrador do contrato
    \seta{} Coordenação regional \seta{} Coordenação da escola \seta{}
    Professor e Estudante.
  \item Cada nível cria os que estão abaixo dele, nunca o próprio nível nem um
    acima. A Coordenação da escola, por exemplo, só oferece Professor e
    Estudante.
  \item O papel que não pode ser dado não aparece na tela --- e, se alguém
    tentar por outro caminho, a plataforma recusa.
  \item Ao criar uma Coordenação regional, a janela troca Turma e Turno pelo
    campo Regional.
  \item Só o Administrador Unooria cria Administrador do contrato.
\end{itemize}
\end{historia}

\begin{historia}{H03.4.2}{Escolher a escola numa lista}{\pronto}
\hu{Administrador do contrato}{escolher a escola numa lista quando trabalho em
Usuários, Turmas, Importar e Análise de produto}{não ter de saber de cor o
código interno de cada escola}
\aceite
\begin{itemize}
  \item Nas quatro telas, \enquote{Buscar escola} filtra por nome, código INEP
    ou município --- sem diferenciar acento nem maiúscula --- a lista
    \enquote{Escola}, em que cada escola aparece com o município e o código
    INEP.
  \item Abaixo da lista, a tela diz quantas escolas a busca encontrou, ou que
    nenhuma combina.
  \item O campo \enquote{Escola (id)}, que pedia o código interno digitado,
    deixou de existir.
  \item A escolha aparece só para quem atende a rede inteira (Administrador do
    contrato e Administrador Unooria).
\end{itemize}
\end{historia}

% ---------------------------------------------------------------------------
\feature{F03.5}{Turmas e enturmamento}{\pronto}

\begin{historia}{H03.5.1}{Criar as turmas do ano}{\pronto}
\hu{Coordenação da escola}{criar as turmas da escola}{enturmar estudantes e
professores}
\aceite
\begin{itemize}
  \item \enquote{Nova turma} pede Nome, Série e Turno.
  \item Nome repetido no mesmo ano letivo é recusado.
  \item A lista mostra Série, Turno, Ano letivo, Estudantes e Professores, com
    filtro por ano letivo.
\end{itemize}
\end{historia}

\begin{historia}{H03.5.2}{Pôr, mover e tirar pessoas da turma}{\pronto}
\hu{Coordenação da escola}{pôr pessoas na turma, mudar de turma e
tirar}{que cada estudante esteja na sala certa}
\aceite
\begin{itemize}
  \item Na turma, abas Estudantes, Professores e Missões;
    \enquote{Adicionar estudante} e \enquote{Adicionar professor} por busca.
  \item \enquote{Mover} leva para outra turma da escola num passo só --- o
    estudante nunca fica \enquote{sem turma} no meio do caminho.
  \item \enquote{Remover} encerra a matrícula sem apagar o que foi feito.
  \item \enquote{Tornar referência} define a turma principal do estudante,
    que é a que as missões usam.
  \item Existem turmas regulares e turmas de componente; o estudante pode estar
    nas duas.
\end{itemize}
\end{historia}

\begin{historia}{H03.5.3}{Ver as missões da turma sem mexer}{\pronto}
\hu{Coordenação da escola}{ver quais missões estão abertas numa
turma}{responder ao professor ou à família sem precisar perguntar}
\aceite
\begin{itemize}
  \item A aba Missões mostra o que está aberto e quem abriu, sem permitir
    mudar.
  \item Aviso na aba: \enquote{Quem abre a missão para a turma é quem dá aula
    nela}.
\end{itemize}
\end{historia}

% ---------------------------------------------------------------------------
\feature{F03.6}{Transferir estudante dentro da rede}{\pronto}

\begin{historia}{H03.6.1}{Transferir um estudante de escola}{\pronto}
\hu{Administrador do contrato}{transferir um estudante para outra escola da
rede}{ele seguir na plataforma sem perder a identidade}
\aceite
\begin{itemize}
  \item Na turma, \enquote{Transferir de escola}.
  \item O que o estudante produziu fica na escola de origem; a identidade o
    acompanha.
  \item A escola nova vê só o \enquote{Histórico} de participações, nunca as
    respostas dele.
\end{itemize}
\pendencia{o botão aparece só para o Administrador do contrato; para a
Coordenação, a transferência acontece indiretamente, importando o CPF. Ver a
decisão \id{D02}.}
\end{historia}

% ---------------------------------------------------------------------------
\feature{F03.7}{Limpar contas sem matrícula}{\pronto}

\begin{historia}{H03.7.1}{Desativar as contas de aparelho antigas}{\pronto}
\hu{Coordenação}{ver as contas que não têm matrícula na rede (as antigas
contas \enquote{de tablet})}{desativá-las sem apagar trabalho de criança}
\aceite
\begin{itemize}
  \item A seção \enquote{Contas sem matrícula da rede} lista essas contas com
    quanto dado cada uma acumulou (até 200 por vez).
  \item Posso desativar; não existe \enquote{fundir} com a conta verdadeira.
\end{itemize}
\end{historia}

% ---------------------------------------------------------------------------
\feature{F03.8}{Perfil portátil: o estudante que sai da rede}{\pronto}

A família tem direito de levar o histórico do estudante quando ele sai da
rede. O perfil portátil é esse histórico em forma de documento assinado.

\begin{historia}{H03.8.1}{Emitir o perfil a pedido da família}{\pronto}
\hu{Coordenação}{emitir o percurso do estudante num documento
assinado}{atender o pedido da família quando ele muda de rede}
\aceite
\begin{itemize}
  \item Na turma, \enquote{Perfil portátil} \seta{} Emitir.
  \item Aviso: \enquote{Isto é para sair da REDE, não da escola}.
  \item O documento nunca inclui dado de bem-estar.
\end{itemize}
\end{historia}

\begin{historia}{H03.8.2}{Receber o perfil que veio de outra rede}{\pronto}
\hu{Coordenação}{receber o documento que o estudante trouxe de outra
rede}{conhecer o percurso dele sem misturar com o que ele fizer aqui}
\aceite
\begin{itemize}
  \item A plataforma confere a assinatura: documento adulterado é recusado; o
    que não pôde ser conferido aparece como \enquote{Recebido, mas não
    conferido}.
  \item O que chega de fora nunca vira nota nem progresso: fica anexado ao
    estudante, identificado pelo CPF.
  \item A lista \enquote{Recebidos} mostra os documentos já recebidos.
\end{itemize}
\end{historia}

% ===========================================================================
%  E04 — Sala de aula: o professor abre a missão
% ===========================================================================
\chapter{E04 \textperiodcentered{} Sala de aula: o professor abre a missão}
\label{ep:e04}

\begin{fichaepico}
\begin{campos}
  \item[Objetivo] O professor conhece as missões, escolhe a da aula e a abre
    para a turma certa, no dia certo.
  \item[Para quem] Professor. A Coordenação consulta, sem mexer.
  \item[Valor] A missão acontece \emph{na aula}, com o professor presente. Se
    ela abrisse sozinha, o estudante jogaria em casa, sem o debate --- e o
    debate é onde a aprendizagem acontece.
  \item[Estado] \pronto{} --- falta a liberação com data marcada.
\end{campos}
\end{fichaepico}

\begin{tabularx}{\linewidth}{@{}L{16mm}X L{20mm}@{}}
\toprule
\headrow \thd{ID} & \thd{Feature} & \thd{Estado}\\
\midrule
\id{F04.1} & Minhas turmas & \pronto\\
\zebra \id{F04.2} & Liberar e encerrar missões na sala & \pronto\\
\id{F04.3} & Conhecer a trilha e a ficha da missão & \pronto\\
\zebra \id{F04.4} & Guia da aula & \pronto\\
\bottomrule
\end{tabularx}

% ---------------------------------------------------------------------------
\feature{F04.1}{Minhas turmas}{\pronto}

\begin{historia}{H04.1.1}{Começar pelo meu dia}{\pronto}
\hu{Professor}{ver no início os produtos que a minha rede contratou e as
minhas turmas}{ir direto ao que vou usar hoje}
\aceite
\begin{itemize}
  \item O Início (\enquote{Olá, \{nome\}}) mostra os cartões dos produtos:
    Gerador de Plano de Aula, Planejamento de Ensino, Entre Redes e Como você
    está?. O do bem-estar mostra a situação do dia (ex.: \enquote{Você já
    respondeu hoje \textperiodcentered{} 2 de 7}).
  \item Produto desligado para a rede não aparece --- nem como cartão
    desativado.
  \item \enquote{Minhas turmas} mostra as 4 primeiras, com \enquote{Ver
    todas}.
  \item \enquote{Continue de onde parou} mostra as atividades recentes do
    Instanooria.
\end{itemize}
\end{historia}

\begin{historia}{H04.1.2}{Ver só as turmas em que dou aula}{\pronto}
\hu{Professor}{ver as minhas turmas}{não me confundir com as turmas dos
colegas}
\aceite
\begin{itemize}
  \item \enquote{Minhas turmas} lista nome, série, turno, número de estudantes
    e período --- só das turmas em que estou matriculado.
  \item Sem turma, a tela orienta: \enquote{Quem faz o enturmamento é a
    coordenação da escola}.
  \item A turma de outro professor, aberta pelo endereço, não é encontrada.
\end{itemize}
\end{historia}

% ---------------------------------------------------------------------------
\feature{F04.2}{Liberar e encerrar missões na sala}{\pronto}

\begin{historia}{H04.2.1}{Liberar uma missão para a turma}{\pronto}
\hu{Professor}{liberar a missão para a minha turma no dia da aula}{que só essa
turma entre, e só quando eu estiver com ela}
\aceite
\begin{itemize}
  \item Na sala da turma, a aba Missões mostra a grade da série, com a
    etiqueta de cada missão: \enquote{aberta}, \enquote{fechada} ou
    \enquote{encerrada}.
  \item \enquote{Liberar} abre a missão, e a linha mostra \enquote{aberta por
    \{nome\} em \{data\}}.
  \item Só quem dá aula na turma consegue liberar.
  \item Sem liberação, o estudante vê o cartão com cadeado --- e digitar o
    endereço da missão também não abre.
  \item Sem nenhuma missão liberada, a sala avisa: \enquote{Nenhuma missão está
    aberta para esta turma}.
  \item Cada missão tem um único caminho a partir da sala: \enquote{Conduzir
    com esta turma}.
\end{itemize}
\end{historia}

\begin{historia}{H04.2.2}{Encerrar uma missão}{\pronto}
\hu{Professor}{encerrar a missão quando a sequência de aulas acabar}{a turma
seguir para a próxima}
\aceite
\begin{itemize}
  \item \enquote{Encerrar} fecha a entrada, mas não expulsa quem está jogando
    naquele momento.
  \item Nada do que a turma fez é apagado; a missão continua visível como
    \enquote{encerrada}.
\end{itemize}
\end{historia}

\begin{historia}{H04.2.3}{Programar a abertura e o fechamento}{\afazer}
\hu{Professor}{deixar a missão programada para abrir e fechar em datas
certas}{não depender de lembrar de liberar na hora da aula}
\aceiteproposta
\begin{itemize}
  \item Ao liberar, posso escolher \enquote{agora} ou uma janela (abre em
    \ldots{}, fecha em \ldots{}).
  \item A sala mostra a janela programada.
\end{itemize}
\end{historia}

\begin{historia}{H04.2.4}{Ver os estudantes da turma}{\pronto}
\hu{Professor}{ver quem está na turma e quem ainda não entrou}{ajudar quem está
com problema antes da aula}
\aceite
\begin{itemize}
  \item A aba Estudantes lista nome, usuário e matrícula, com as etiquetas
    \enquote{inativo} e \enquote{sem primeiro acesso}.
  \item O nome do estudante abre o perfil dele, com as habilidades da BNCC que
    ele já demonstrou (\id{H14.6.1}); ao lado, \enquote{Relatório da Vila do
    Lume} abre o relatório dele na Vila.
  \item O bloco \enquote{Também nesta turma} mostra os outros professores.
\end{itemize}
\end{historia}

% ---------------------------------------------------------------------------
\feature{F04.3}{Conhecer a trilha e a ficha da missão}{\pronto}

\begin{historia}{H04.3.1}{Consultar a trilha por ano}{\pronto}
\hu{Professor}{ver quais missões existem para cada ano}{planejar quando usar
cada uma}
\aceite
\begin{itemize}
  \item A trilha Entre Redes é organizada por etapa (Anos Iniciais, Anos
    Finais, Ensino Médio); o bloco \enquote{Para qualquer ano} vem primeiro.
  \item Anos com missão mostram \enquote{N de 4 unidades publicadas}; anos sem
    missão aparecem numa linha, \enquote{4 unidades em produção} --- sem
    título inventado.
  \item Faixa fixa: \enquote{Aqui você conhece as missões; quem abre para a
    turma é você, na sala}.
\end{itemize}
\end{historia}

\begin{historia}{H04.3.2}{Ler a ficha antes de levar para a turma}{\pronto}
\hu{Professor}{entender a missão antes de rodá-la}{conduzir sem estragar a
surpresa da revelação}
\aceite
\begin{itemize}
  \item A ficha tem \enquote{Como a missão funciona}, \enquote{Antes de rodar
    com a turma}, \enquote{O que o software NÃO conta ao estudante} e o
    roteiro.
  \item Botões \enquote{Abrir guia da aula} e \enquote{Abrir a aula e os
    resultados}.
  \item Vindo da sala, aparece a faixa \enquote{Conduzindo com a turma
    \{nome\}}.
  \item Unidade ainda não publicada: \enquote{\{missão\} ainda está em
    produção}.
\end{itemize}
\end{historia}

% ---------------------------------------------------------------------------
\feature{F04.4}{Guia da aula}{\pronto}

\begin{historia}{H04.4.1}{Ter um roteiro pronto para a aula}{\pronto}
\hu{Professor}{um roteiro de condução pronto para cada missão}{dar a aula sem
precisar prepará-la do zero}
\aceite
\begin{itemize}
  \item O guia tem nove seções fixas: objetivos e habilidade da BNCC
    Computação; preparo; o que não dizer; roteiro minuto a minuto; banco de
    perguntas; mediação; rubrica; adaptações; e o conteúdo real que a turma
    vê.
  \item Pode ser impresso.
  \item Existe para as nove missões publicadas.
\end{itemize}
\pendencia{o guia do unooGPT ainda não passou por revisão pedagógica em
sala.}
\end{historia}

% ===========================================================================
%  E05 — A experiência do estudante na plataforma
% ===========================================================================
\chapter{E05 \textperiodcentered{} A experiência do estudante na plataforma}
\label{ep:e05}

\begin{fichaepico}
\begin{campos}
  \item[Objetivo] O estudante acha a missão do dia, joga sem atrito e não perde
    o que fez --- em qualquer missão, antes mesmo de ela começar.
  \item[Para quem] Estudante, do 1º ano EF à 3ª série EM.
  \item[Valor] O que acontece \emph{dentro} de cada missão está no épico
    \id{E06}. Este épico é o que é igual em todas: a porta de entrada.
  \item[Estado] \pronto.
\end{campos}
\end{fichaepico}

\begin{tabularx}{\linewidth}{@{}L{16mm}X L{20mm}@{}}
\toprule
\headrow \thd{ID} & \thd{Feature} & \thd{Estado}\\
\midrule
\id{F05.1} & Minhas missões & \pronto\\
\zebra \id{F05.2} & Jogar a missão dentro da plataforma & \pronto\\
\bottomrule
\end{tabularx}

% ---------------------------------------------------------------------------
\feature{F05.1}{Minhas missões}{\pronto}

\begin{historia}{H05.1.1}{Ver as missões da minha turma}{\pronto}
\hu{Estudante}{ver as missões do meu ano e saber qual está aberta}{começar a
missão que o professor pediu}
\aceite
\begin{itemize}
  \item \enquote{Minhas missões} (\enquote{Olá, \{nome\}}) mostra a turma
    (ex.: \enquote{Turma 8A \textperiodcentered{} 8º ano \textperiodcentered{}
    manhã}) e um cartão por missão da grade do ano.
  \item Cartão aberto: \enquote{Começar}. Cartão fechado: cadeado e
    \enquote{Ainda não abriu}. Missão encerrada: etiqueta \enquote{encerrada}
    e \enquote{Esta você já fechou}.
  \item Aparecem só as missões da série da turma, mais as \enquote{para
    qualquer ano}.
  \item A missão fechada \textbf{não some}: o estudante vê que ela existe.
  \item Missão desligada pela Unooria para a rede some da grade.
\end{itemize}
\end{historia}

\begin{historia}{H05.1.2}{Entender quando não há nada para fazer}{\pronto}
\hu{Estudante}{uma mensagem clara quando não tenho missão}{não achar que a
plataforma quebrou}
\aceite
\begin{itemize}
  \item Mensagens diferentes para cada caso: \enquote{Você ainda não está em
    uma turma}, \enquote{Nenhuma missão por aqui ainda} e \enquote{Não
    conseguimos carregar suas missões}.
\end{itemize}
\end{historia}

% ---------------------------------------------------------------------------
\feature{F05.2}{Jogar a missão dentro da plataforma}{\pronto}

\begin{historia}{H05.2.1}{Abrir a missão}{\pronto}
\hu{Estudante}{abrir a missão com um toque}{começar a aula junto com a turma}
\aceite
\begin{itemize}
  \item A missão abre dentro de uma moldura da plataforma --- em formato de
    celular ou de tela cheia, conforme a missão --- com \enquote{\ensuremath{\leftarrow}{}
    Missões} no cabeçalho.
  \item Missão não liberada, de outra série ou desligada não abre, nem pelo
    endereço digitado.
  \item Para a criança, a mensagem de missão desligada é a mesma de missão
    ainda não aberta, para não quebrar a ficção.
\end{itemize}
\end{historia}

\begin{historia}{H05.2.2}{Não perder o que fiz}{\pronto}
\hu{Estudante}{que o que eu fiz na missão fique guardado}{não recomeçar se a
página recarregar ou a internet oscilar}
\aceite
\begin{itemize}
  \item O que a criança fez fica guardado na plataforma e volta depois de
    recarregar a página.
  \item Quando a janela da missão perde o foco, a missão pausa (os cronômetros
    congelam) e retoma ao voltar.
  \item A sessão não expira enquanto a criança está na missão (\id{H01.1.3}).
\end{itemize}
\pendencia{se a conexão cair de vez, os registros de uso ainda não enviados se
perdem (o progresso da missão, não). Guardar esses registros sem conexão é
evolução planejada.}
\end{historia}

\begin{historia}{H05.2.3}{Esperar o professor sem confusão}{\pronto}
\hu{Estudante}{saber quando estou esperando o professor}{não achar que travou}
\aceite
\begin{itemize}
  \item Várias missões param a turma até o professor soltar: a largada do
    relógio (unooGPT), a Aula 2 (Eco no Grupo, Rastro do Boato), o debate
    (Manifesto da IA), a roda de conversa (O Que a Foto Conta) e a revelação,
    a síntese e a reflexão (Quem mandou isso?).
  \item Nesses momentos a tela mostra uma espera clara, nunca um erro.
\end{itemize}
\end{historia}

\begin{historia}{H05.2.4}{Nunca ser comparado com os colegas}{\pronto}
\hu{família do estudante}{que nenhuma tela mostre nota ou ranking da
criança}{que a missão seja espaço de aprender, não de competir}
\aceite
\begin{itemize}
  \item Nenhuma tela do estudante mostra nota ou posição em ranking de
    colegas reais.
  \item Quando a missão tem ranking como parte da ficção (UnooTube), os
    colegas do ranking são fictícios.
  \item A correção acontece na plataforma; a resposta certa não chega ao
    aparelho antes da hora.
\end{itemize}
\end{historia}

% ===========================================================================
%  E06 — As missões do Entre Redes
% ===========================================================================
\chapter{E06 \textperiodcentered{} As missões do Entre Redes}
\label{ep:e06}

\begin{fichaepico}
\begin{campos}
  \item[Objetivo] Cada missão faz a turma viver por dentro como as redes
    funcionam --- e depois discutir o que viveu.
  \item[Para quem] Estudante (joga) e Professor (prepara, conduz e lê a
    turma).
  \item[Valor] É o coração do produto Entre Redes e a razão de ele existir:
    letramento digital pela experiência, não pela aula expositiva.
  \item[Estado] \pronto{} para 8 missões (7 da trilha + Vila do Lume), e
    \parcial{} para a versão nova de Quem mandou isso? (\id{F06.11}), que
    devolveu ao 5º ano a sua missão; 40 unidades da trilha em produção. A
    primeira versão (\id{F06.3}) saiu em 25 de setembro de 2026 para dar lugar a
    ela.
\end{campos}
\end{fichaepico}

\section{Como uma missão é feita}

Todas as missões seguem o mesmo desenho, e é ele que dá identidade ao
produto:

\begin{enumerate}
  \item \textbf{O disfarce.} O estudante entra num aplicativo que conhece
    (rede social de fotos, vídeos curtos, grupo de mensagens, assistente de
    IA) ou num jogo, e age como agiria no próprio celular.
  \item \textbf{A decisão.} Em algum momento ele escolhe --- o que repassar,
    em quem confiar, o que salvar --- sem saber que aquilo será olhado depois.
  \item \textbf{A revelação.} A missão deixa de fingir e mostra o que estava
    acontecendo por trás. Em algumas, a revelação é da própria tela; em
    outras, acontece no debate conduzido pelo professor.
  \item \textbf{O debate.} Os aparelhos congelam e a turma conversa. É aqui
    que a aprendizagem acontece, e por isso quase toda missão tem um momento
    que só o professor destrava.
  \item \textbf{O percurso.} O estudante registra o que levou da missão (o
    \enquote{Meu percurso}), sem nota.
\end{enumerate}

\begin{tabularx}{\linewidth}{@{}L{16mm}X L{24mm}L{26mm}L{18mm}@{}}
\toprule
\headrow \thd{ID} & \thd{Missão} & \thd{Ano / unidade} & \thd{Aplicativo imitado} & \thd{Estado}\\
\midrule
\id{F06.1} & O Que a Foto Conta & 1º EF, U1 & rede social de fotos & \parcial\\
\zebra \id{F06.2} & UnooTube & 4º EF, U1 & vídeos curtos & \pronto\\
\id{F06.11} & Quem Mandou Isso? (versão nova) & 5º EF, U1 & mensageiro & \parcial\\
\zebra \id{F06.4} & O Rastro do Boato & 6º EF, U2 & mensageiro & \pronto\\
\id{F06.5} & O Eco no Grupo & 7º EF, U1 & mensageiro & \pronto\\
\zebra \id{F06.6} & Instanooria & 8º EF, U1 & rede social de fotos & \pronto\\
\id{F06.7} & O Manifesto da IA & 8º EF, U4 & assistente de IA & \pronto\\
\zebra \id{F06.8} & unooGPT (Missão Tempestade) & 2ª EM, U1 & assistente de IA & \pronto\\
\id{F06.9} & Vila do Lume & qualquer ano & jogo 3D próprio & \pronto\\
\zebra \id{F06.10} & As missões que faltam & 40 unidades & --- & \afazer\\
\bottomrule
\end{tabularx}

% ---------------------------------------------------------------------------
\feature{F06.1}{O Que a Foto Conta}{\parcial}

\textbf{1º ano EF, Unidade 1} \textperiodcentered{} imita uma rede social de
fotos (o \enquote{mural da turma}) \textperiodcentered{} uma aula
\textperiodcentered{} IA só no preparo do professor.

\textbf{Proposta.} Letramento visual para quem está se alfabetizando: toda
foto é uma escolha de enquadramento e legenda, e o que fica de fora também
conta a história. É 100\% formativa: não há foto errada nem nota.

\begin{historia}{H06.1.1}{Publicar uma foto no mural da turma}{\parcial}
\hu{Estudante do 1º ano}{escolher uma foto e uma legenda e publicar no
mural}{descobrir que cada foto mostra só uma parte da história}
\aceite
\begin{itemize}
  \item Nina dá as boas-vindas em áudio e a criança escolhe a sua marca
    (Coruja Azul, Estrela Amarela ou Sol Laranja).
  \item Explora um mural com 3 fotos e curte com toque duplo; parada por 30
    segundos, uma seta indica o caminho.
  \item Recebe duas fotos do mesmo pátio (plano geral com peças espalhadas
    \emph{versus} detalhe do brinquedo montado) e duas legendas; escolhe uma
    de cada e toca em COMPARTILHAR.
  \item Depois de publicar, os colegas simulados Leo e Sofia comentam o que a
    foto mostrou e o que escondeu.
  \item Todo texto em CAIXA ALTA, com no máximo 8 palavras, pictograma e
    áudio.
\end{itemize}
\pendencia{as fotos ainda são desenhos provisórios e as locuções
profissionais não foram gravadas (a tela mostra texto e ícone). Falta também
um botão maior de repetir o áudio.}
\end{historia}

\begin{historia}{H06.1.2}{Refletir sobre o que a foto mostra}{\pronto}
\hu{Estudante do 1º ano}{responder o que a foto mostra depois da roda de
conversa}{guardar a ideia de que olhar é escolher}
\aceite
\begin{itemize}
  \item Todos os tablets congelam em \enquote{HORA DA RODA DE CONVERSA!}, uma
    tela sem nenhum botão.
  \item Depois, a reflexão \enquote{O QUE A FOTO MOSTRA?}: a resposta Lupa
    conclui com confete e a insígnia \enquote{Olhar de Detetive}; a resposta
    Globo recebe um lembrete gentil e a criança tenta de novo, sem limite.
\end{itemize}
\end{historia}

\begin{historia}{H06.1.3}{Adaptar a missão ao tema da semana}{\pronto}
\hu{Professor}{gerar legendas e comentários sobre o tema que a turma está
estudando}{que a missão converse com a minha aula}
\aceite
\begin{itemize}
  \item Temas: Pátio do Recreio, Horta, Sala de Aula, Parquinho.
  \item Gero o pacote com IA, confiro e aprovo; nada chega à turma sem
    aprovação.
  \item Se eu não fizer nada, roda o pacote pronto \enquote{O Mistério das
    Peças no Pátio}.
  \item A IA nunca roda no tablet da criança e nunca a avalia.
\end{itemize}
\end{historia}

\begin{historia}{H06.1.4}{Conduzir a roda e ver quem hesitou}{\pronto}
\hu{Professor}{conduzir a roda de conversa e saber quem teve dificuldade de
leitura}{ajudar quem precisa, sem expor ninguém}
\aceite
\begin{itemize}
  \item Projeto as duas fotos lado a lado e conduzo a roda; depois toco em
    \enquote{LIBERAR REFLEXÃO FINAL}.
  \item Vejo a lista de \enquote{hesitação de leitura} (mais de 60 segundos na
    decisão), com nome, só para mim.
  \item \enquote{Zerar a turma} arquiva a tentativa, em vez de apagar.
\end{itemize}
\end{historia}

% ---------------------------------------------------------------------------
\feature{F06.2}{UnooTube --- Em quais fontes você confiaria?}{\pronto}

\textbf{4º ano EF, Unidade 1} \textperiodcentered{} imita um aplicativo de
vídeos curtos, em tela de computador \textperiodcentered{} uma aula
\textperiodcentered{} sem IA.

\textbf{Proposta.} Distinguir fonte confiável de sensacionalismo e perceber
que ter muitos inscritos não significa contar a verdade.

\begin{historia}{H06.2.1}{Salvar as três fontes em que eu confiaria}{\pronto}
\hu{Estudante do 4º ano}{assistir aos vídeos e salvar os três em que eu
confiaria}{fazer uma escolha de verdade antes da revelação}
\aceite
\begin{itemize}
  \item Um feed de 21 vídeos curtos de História, em ordem aleatória; posso
    ligar o som, abrir a descrição, curtir e salvar.
  \item A pesquisa do aplicativo está bloqueada de propósito.
  \item Cabem 3 vídeos na bandeja; ao tentar um quarto, o aviso \enquote{sua
    playlist tem três lugares} oferece trocar.
  \item No terceiro salvo, habilita \enquote{Concluir missão} --- quem decide
    quando fechar é o estudante.
\end{itemize}
\end{historia}

\begin{historia}{H06.2.2}{Ver o ranking virar}{\pronto}
\hu{Estudante do 4º ano}{comparar as minhas escolhas com as da turma}{perceber
que popularidade não é credibilidade}
\aceite
\begin{itemize}
  \item O Ranking da turma ordena por inscritos, com colegas fictícios: quem
    escolheu vídeos falsos aparece no topo.
  \item \enquote{Adicionar credibilidade} revela a coluna de fontes confiáveis,
    e o ranking vira.
  \item Depois disso os salvos congelam; reiniciar descongela.
  \item Os links de fonte aparecem como texto: o estudante não sai da
    plataforma.
\end{itemize}
\end{historia}

\begin{historia}{H06.2.3}{Ler a turma}{\pronto}
\hu{Professor}{ver quantas fontes confiáveis a turma escolheu}{saber por onde
começar o debate}
\aceite
\begin{itemize}
  \item Não há nada a configurar além de liberar.
  \item Resultados: \enquote{Estudantes na turma}, \enquote{Média de fontes
    confiáveis} e a distribuição de 0 a 3 fontes confiáveis.
  \item O guia traz a tabela dos 21 vídeos, com quais são confiáveis.
\end{itemize}
\pendencia{os vídeos precisam estar publicados no armazenamento de cada rede;
sem isso, o feed aparece sem vídeo.}
\end{historia}

% F06.3 (a primeira versão de Quem mandou isso?) saiu do produto em 25 de
% setembro de 2026: a versão nova da missão ocupa a Unidade 1 do 5º ano no
% lugar dela (H06.10.3, detalhada em F06.11). Pela regra do documento, o número
% fica vago e não é reaproveitado, e as histórias H06.3.1 a H06.3.4 também.

% ---------------------------------------------------------------------------
% F06.11 é a versão nova de Quem mandou isso?. O número é novo porque o F06.3
% (a primeira versão) não se reaproveita; ela fica aqui, na posição do 5º ano.
\feature{F06.11}{Quem Mandou Isso? --- versão nova}{\parcial}

\textbf{5º ano EF, Unidade 1} \textperiodcentered{} imita um mensageiro, em
tema escuro \textperiodcentered{} uma aula de 50 minutos
\textperiodcentered{} IA só no preparo do professor \textperiodcentered{}
aceita equipes.

\textbf{Proposta.} Perceber que o jeito de uma mensagem --- séria ou alarmada
--- não diz se ela é verdadeira, e montar a própria regra para conferir antes
de repassar. É a primeira missão que registra o que cada estudante demonstrou
das habilidades da BNCC (\id{E14}).

\begin{historia}{H06.11.1}{Preparar o pacote de mensagens da turma}{\parcial}
\hu{Professor}{preparar as oito mensagens do grupo sobre um assunto da minha
escola}{que a turma leia uma confusão que parece a dela, e não uma qualquer}
\aceite
\begin{itemize}
  \item Em \enquote{A aula e os resultados}, escolho o componente curricular,
    escrevo o tema --- ele completa a fala da Nina, \enquote{Rolou uma confusão
    sobre\ldots}, com até 60 caracteres e 8 palavras ---, listo de 3 a 6 fatos
    da aula e, se quiser, o contexto da escola.
  \item \enquote{Gerar Atividade com IA} escreve as oito mensagens; sem saldo
    ou com a IA fora do ar, a tela avisa e já traz o pacote pronto da Unooria,
    \enquote{O Mistério do Passeio ao Museu}. \enquote{Usar Pacote Curado
    Padrão} o traz também, sem gerar nada.
  \item Vejo as oito em quatro grupos --- verdadeira ou falsa, emocional ou
    neutra ---, com os selos que só eu vejo, e posso reescrever qualquer texto;
    a resposta de cada mensagem não se edita.
  \item Toda turma recebe duas mensagens de cada grupo, quatro textos, duas
    imagens e dois links --- e a imagem e o link têm uma verdadeira e uma falsa
    cada, para o formato não entregar a resposta.
  \item Nada chega à turma sem o meu clique em \enquote{Aprovar Pacote}; sem
    ele, a criança abre a missão e vê a lista de conversas vazia, com um aviso
    de espera.
  \item Aprovar um pacote diferente pede confirmação: a turma recomeça com as
    mensagens novas, e o que ela fez com o anterior fica guardado.
  \item A tela diz se a missão já está liberada na sala da turma, e leva à
    ficha para montar as equipes.
\end{itemize}
\pendencia{textos escritos pela equipe --- a prévia do link do pacote pronto,
os fatos dele e a lista de componentes curriculares --- aguardam a revisão do
PO; o equilíbrio de imagem e link no pacote pronto espera a decisão \id{D14}.}
\end{historia}

\begin{historia}{H06.11.2}{Escolher exatamente três mensagens para o Mural}{\parcial}
\hu{Estudante do 5º ano}{ler o grupo da turma e mandar ao Mural da Turma as
três mensagens que explicam o que aconteceu}{fazer o que eu faria no meu
celular, sem ninguém me dizer em que desconfiar}
\aceite
\begin{itemize}
  \item A missão abre na lista de conversas; a Nina pede ajuda sobre o tema da
    turma e manda o convite do grupo \enquote{Turma 5ºB}.
  \item No grupo, oito mensagens de colegas, nenhuma marcada; só
    administradores escrevem, e a única ação é a seta de encaminhar, que abre
    \enquote{Encaminhar para\ldots} com um destino só: o Mural da Turma.
  \item Encaminho exatamente três. A quarta seta diz \enquote{Limite de 3
    mensagens atingido. Desmarque uma mensagem para escolher outra.}, e a faixa
    do topo, com menos de três, diz o que falta.
  \item No Mural as três aparecem como \enquote{Encaminhada}, sem o nome de
    quem mandou. Posso voltar ao grupo e trocar até tocar em \enquote{Confirmar
    envio}; depois, não muda mais.
  \item Parado 90 segundos sem encaminhar, a Nina lembra como encaminhar --- o
    gesto, nunca a escolha.
  \item Se eu recarregar a página, volto onde estava, com as escolhas
    guardadas.
\end{itemize}
\pendencia{o teste que percorre a missão como o estudante e como o professor
foi escrito e ainda não rodou.}
\end{historia}

\begin{historia}{H06.11.3}{Dizer o que pesou e descobrir os dois selos}{\parcial}
\hu{Estudante do 5º ano}{dizer por que escolhi aquelas mensagens e depois ver o
que cada uma era}{perceber que mensagem séria pode ser falsa, e alarmada pode
ser verdadeira}
\aceite
\begin{itemize}
  \item Confirmado o Mural, respondo à enquete \enquote{O que mais te fez
    escolher essas mensagens?}, com quatro opções, e espero a turma.
  \item Quando o professor libera a revelação, cada uma das minhas três ganha
    dois selos --- verdadeira ou falsa, emocional ou neutra ---, com o motivo
    que eu votei, um lembrete do bimestre e as duas barras da turma, sem nome
    de ninguém. Chega sozinho, sem recarregar a página.
  \item Antes do meu voto e da revelação, a resposta de nenhuma mensagem chega
    ao meu aparelho.
  \item Quem ainda estava escolhendo quando o professor liberou continua,
    confirma, vota e vê os próprios selos depois.
  \item A síntese chega ao aparelho quando o professor a apresenta, com o aviso
    de participar da conversa na sala.
\end{itemize}
\pendencia{o teste que percorre a missão foi escrito e ainda não rodou.}
\end{historia}

\begin{historia}{H06.11.4}{Montar a minha regra e responder à Nina}{\parcial}
\hu{Estudante do 5º ano}{escrever a minha regra para antes de repassar e
responder às perguntas da Nina}{levar a missão para o grupo da minha família}
\aceite
\begin{itemize}
  \item Com a reflexão liberada, o aplicativo fecha: na tela da escola, a Nina
    conta do áudio do tio no grupo da família, e eu monto \enquote{Se a
    mensagem for\ldots, então eu\ldots}.
  \item O \enquote{então} pede pelo menos 15 caracteres --- espaço não conta ---
    e até 240; texto sem sentido é recusado com um pedido de mais detalhe.
  \item A regra não é corrigida: vai para o professor da turma e para mais
    ninguém, e a tela confirma \enquote{Sua regra foi enviada para o
    professor.}
  \item De volta ao chat da Nina, respondo a duas perguntas rápidas com a mesma
    enquete do Mural; depois de cada uma, vejo só \enquote{Resposta enviada} ---
    nunca se acertei.
  \item A Nina se despede, e a tela da escola mostra \enquote{Missão
    Concluída} com \enquote{Sua Regra Pessoal de Checagem}.
\end{itemize}
\pendencia{a despedida da Nina e a confirmação da regra são textos da equipe e
aguardam a revisão do PO; o teste que percorre a missão ainda não rodou.}
\end{historia}

\begin{historia}{H06.11.5}{Conduzir a aula e ler a turma}{\parcial}
\hu{Professor}{conduzir a aula em três momentos e ver o que a turma
escolheu}{abrir a roda de conversa com o retrato da sala, sem expor ninguém}
\aceite
\begin{itemize}
  \item Três botões, na ordem --- \enquote{Liberar Revelação},
    \enquote{Apresentar Síntese} e \enquote{Liberar Reflexão} ---, cada um com
    confirmação e sempre para a frente; a síntese pode ser pulada.
  \item O painel se atualiza sozinho: quantos entraram no grupo, confirmaram o
    Mural, votaram, enviaram a regra e concluíram.
  \item Depois da revelação, o Termômetro da Turma: quantos repassaram uma
    falsa de formato sério, quantas verdadeiras exageradas ficaram de fora e o
    cruzamento com o motivo votado. \enquote{Projetar o Termômetro} abre uma
    tela cheia sem nome nenhum; com ninguém que confirmou, \enquote{sem dados},
    nunca 0\%.
  \item O roteiro da roda de conversa, com três perguntas.
  \item A folha da turma: onde cada um está, as três mensagens com o grupo de
    cada uma, o motivo e a regra --- o texto da regra, só para quem dá aula na
    turma. Sem coluna de acerto.
  \item O bloco \enquote{Habilidades da BNCC} mostra a EF05CO08 pela escolha do
    Mural, a EF05CO03 e a EF05CO04 pelas duas perguntas da Nina, e é ali que
    confirmo e ajusto (\id{E14}).
  \item \enquote{Zerar Turma para Nova Aplicação} recomeça a turma numa rodada
    nova, com o mesmo pacote, depois de uma confirmação deliberada; nada é
    apagado.
  \item Se a turma divide computadores, monto as equipes na ficha da missão, e
    o painel conta por computador (\id{E13}).
  \item O guia da aula traz o pacote pronto com as respostas, o roteiro e as
    perguntas da Nina com o que cada alternativa vale.
\end{itemize}
\pendencia{o texto do botão da síntese é da equipe e aguarda o PO; a régua das
habilidades aguarda a validação da equipe pedagógica (\id{H14.1.2}); o teste
que percorre a aula ainda não rodou.}
\end{historia}


% ---------------------------------------------------------------------------
\feature{F06.4}{O Rastro do Boato}{\pronto}

\textbf{6º ano EF, Unidade 2} \textperiodcentered{} imita um mensageiro, em
tema escuro \textperiodcentered{} duas aulas \textperiodcentered{} sem IA.

\textbf{Proposta.} Ler os dados de envio de uma mensagem antes de repassar, e
entender que cada encaminhamento --- ou silêncio --- tem consequência na rede.

\begin{historia}{H06.4.1}{Auditar as mensagens antes de agir}{\pronto}
\hu{Estudante do 6º ano}{abrir os dados de envio das mensagens sobre a água da
escola}{saber de onde o boato veio}
\aceite
\begin{itemize}
  \item Quatro conversas (Grêmio, turma, bairro, Tomás); cada mensagem mostra
    saltos, autoria e primeira aparição.
  \item Preciso auditar 3 mensagens \emph{diferentes} antes de decidir;
    reabrir a mesma não conta.
\end{itemize}
\end{historia}

\begin{historia}{H06.4.2}{Decidir diante do áudio}{\pronto}
\hu{Estudante do 6º ano}{escolher o que faço com o áudio que convoca um
protesto}{assumir uma decisão sem rede de proteção}
\aceite
\begin{itemize}
  \item Uma de 3 respostas, sem volta e sem dica.
  \item Parado por 90 segundos, aparece uma frase sobre o custo de não agir.
  \item Ao fim da Aula 1, uma tela de pausa com o progresso salvo.
\end{itemize}
\end{historia}

\begin{historia}{H06.4.3}{Ver a consequência da minha decisão}{\pronto}
\hu{Estudante do 6º ano}{ver o que a minha decisão fez na rede}{entender que
cada escolha se propaga}
\aceite
\begin{itemize}
  \item O professor toca em \enquote{Liberar Aula 2} (e pode \enquote{Pausar
    simulação}).
  \item Um gráfico mostra o aviso correndo 26 vezes mais rápido que a nota da
    direção; o mapa da rede mostra o resultado da minha decisão.
  \item Monto os 4 blocos do ciclo do boato, entrego o \enquote{Meu
    percurso} e leio o episódio de fechamento.
  \item Quem já entregou não perde o que fez se o professor voltar uma etapa.
\end{itemize}
\end{historia}

\begin{historia}{H06.4.4}{Ler a turma}{\parcial}
\hu{Professor}{ver como a turma respondeu}{conduzir o debate da Aula 2}
\aceite
\begin{itemize}
  \item Distribuição das respostas da turma e material de preparo (as 3
    intervenções e o ciclo correto).
  \item \enquote{Zerar a turma}.
\end{itemize}
\pendencia{o texto do \enquote{Meu percurso} ainda é visível à gestão ---
diverge da regra de \id{H07.1.4}. Faltam também o áudio
gravado e a definição formal de autor e revisor do gabarito.}
\end{historia}

% ---------------------------------------------------------------------------
\feature{F06.5}{O Eco no Grupo}{\pronto}

\textbf{7º ano EF, Unidade 1} \textperiodcentered{} imita um mensageiro, em
tema escuro \textperiodcentered{} duas aulas \textperiodcentered{} IA no
preparo e na conversa com a personagem.

\textbf{Proposta.} Reconhecer que um print ou uma manchete isolada é um recorte
intencional, e intervir com a fonte completa em vez de opinião.

\begin{decisao}
\textbf{Esta missão será substituída.} Em 23 de setembro de 2026 foi decidido
que a missão \enquote{A Seleção} (que imita uma rede social de fotos) ocupa a
Unidade 1 do 7º ano no lugar do Eco no Grupo. O Eco sai primeiro, numa frente
própria, e só depois a Seleção começa. Ver \id{F06.10}.
\end{decisao}

\begin{historia}{H06.5.1}{Preparar a mídia da missão}{\pronto}
\hu{Professor}{gerar a foto e o áudio da missão uma vez}{que a turma tenha a
experiência completa}
\aceite
\begin{itemize}
  \item A preparação leva cerca de 2 minutos e vale para a escola inteira.
  \item Sem mídia preparada, a missão continua jogável pela transcrição.
\end{itemize}
\end{historia}

\begin{historia}{H06.5.2}{Investigar o recorte e intervir}{\pronto}
\hu{Estudante do 7º ano}{descobrir o que o print cortou e responder no
grupo}{intervir com a fonte, não com a raiva}
\aceite
\begin{itemize}
  \item Aula 1: no grupo \enquote{Comunidade Alerta} ouço o áudio, leio as
    mensagens e abro o print; a reportagem completa mostra o parágrafo
    cortado.
  \item Aula 2: escolho entre 3 respostas; responder sem ter aberto a matéria
    completa é recusado dentro da própria ficção.
  \item Quem cita a matéria vê o administrador fixar a mensagem; quem age no
    impulso pode \enquote{apagar para todos} e tentar de novo.
  \item Depois da decisão, posso conversar com Amara (até 3 mensagens).
\end{itemize}
\end{historia}

\begin{historia}{H06.5.3}{Conduzir as duas aulas e ler a conduta}{\pronto}
\hu{Professor}{conduzir as duas aulas e ver como cada estudante chegou à
decisão}{avaliar o processo, não o acerto}
\aceite
\begin{itemize}
  \item Três botões: abrir a Aula 2, congelar para o debate e liberar o
    \enquote{Meu percurso}.
  \item A folha de conduta mostra quem abriu a matéria antes de responder,
    quem ouviu o episódio, quem apagou e tentou de novo e quem entregou.
  \item A distribuição das respostas aparece em porcentagem, sem pontos.
  \item Da conversa com Amara vejo quantas mensagens houve, nunca o conteúdo.
\end{itemize}
\end{historia}

% ---------------------------------------------------------------------------
\feature{F06.6}{Instanooria --- O conteúdo que não chegou}{\pronto}

\textbf{8º ano EF, Unidade 1} \textperiodcentered{} imita uma rede social de
fotos, em moldura de celular \textperiodcentered{} uma aula
\textperiodcentered{} IA no preparo e nas reações ao vivo.

\textbf{Proposta.} Mostrar como um feed decide o que cada pessoa vê e o que
isso faz com uma opinião em formação (a bolha). É a única missão com
\enquote{atividade} configurável pelo professor.

\begin{historia}{H06.6.1}{Criar a atividade a partir da minha aula}{\pronto}
\hu{Professor}{criar uma atividade sobre o tema que estou trabalhando}{que o
feed discuta um assunto da minha disciplina}
\aceite
\begin{itemize}
  \item \enquote{Nova atividade} pede Título, Disciplina, \enquote{Tema /
    recorte}, turma e \enquote{O que já foi trabalhado em aula (opcional)}.
  \item A pergunta-chave não é digitada: a IA propõe candidatas, cada uma com
    uma nota de equilíbrio entre os dois lados.
  \item Escolher uma pergunta desequilibrada exige que eu justifique por
    escrito.
  \item Valido os dois lados (Tese, \enquote{Argumentos (um por linha)},
    \enquote{Tom típico}).
\end{itemize}
\end{historia}

\begin{historia}{H06.6.2}{Revisar o conteúdo antes de ativar}{\pronto}
\hu{Professor}{revisar os posts gerados antes da aula}{garantir que nada
inadequado chegue à turma}
\aceite
\begin{itemize}
  \item A geração acontece em fila (Perfis, Posts, Imagens, Revisão) e a
    atividade passa por Rascunho, Aguardando validação, Gerando conteúdo, Em
    revisão, Ativa e Encerrada.
  \item Filtro os posts por lado, perfil, intensidade, categoria e imagem;
    posso regerar texto ou imagem, ou excluir.
  \item Imagem que falhou recebe uma arte substituta, e a atividade não trava.
  \item Com conteúdo de menos, a plataforma recusa ativar e lista o que falta.
  \item Nenhuma criança entra antes de \enquote{Ativar}.
\end{itemize}
\end{historia}

\begin{historia}{H06.6.3}{Decidir de que lado cada estudante começa}{\pronto}
\hu{Professor}{usar o que a turma revelou na Vila do Lume para definir o lado
de largada}{que cada estudante comece numa bolha parecida com a dele}
\aceite
\begin{itemize}
  \item \enquote{Trazer perfis da Vila do Lume} define o lado de cada
    estudante; quem não jogou a Vila entra num sorteio equilibrado.
  \item \enquote{Redistribuir a turma} permite ajustar à mão, e a plataforma
    não desfaz o meu ajuste.
\end{itemize}
\end{historia}

\begin{historia}{H06.6.4}{Investigar a pergunta pelo feed}{\pronto}
\hu{Estudante do 8º ano}{navegar pelo feed como numa rede social de
verdade}{formar uma opinião sobre a pergunta da missão}
\aceite
\begin{itemize}
  \item Escolho avatar e identidade; uso Stories, curto, comento, salvo,
    compartilho e publico.
  \item O feed se adapta ao que eu presto atenção.
  \item Personagens da rede podem reagir ao que eu escrevo (o professor pode
    desligar isso: \enquote{Perfis respondem ao estudante}).
  \item Nenhum avatar tem rosto realista.
\end{itemize}
\end{historia}

\begin{historia}{H06.6.5}{Descobrir o que o feed escondeu de mim}{\pronto}
\hu{Estudante do 8º ano}{ver o que o feed deixou de me mostrar}{entender como
a minha opinião foi moldada}
\aceite
\begin{itemize}
  \item Escolho um lado e escrevo a conclusão; no debate, a tela trava nos
    posts que salvei.
  \item A explicação mostra como o meu feed foi montado; \enquote{O que
    perderam} mostra o feed do outro lado, em cinza.
  \item Escrevo uma reflexão e recebo o meu relatório (o dossiê), com a
    comparação com a turma, que posso salvar como imagem.
\end{itemize}
\end{historia}

\begin{historia}{H06.6.6}{Conduzir a aula por etapas e acompanhar ao vivo}{\pronto}
\hu{Professor}{mover a turma de etapa em etapa e ver onde cada um está}{que a
turma chegue junta ao debate}
\aceite
\begin{itemize}
  \item Etapas: Começar investigação \seta{} Encerrar investigação \seta{}
    Iniciar debate \seta{} Encerrar debate \seta{} Mostrar o que perderam
    \seta{} Abrir reflexão \seta{} Encerrar atividade.
  \item O monitor mostra \enquote{Onde a turma está}, \enquote{Onde a bolha
    fechou}, \enquote{Viraram o jogo} e \enquote{Tempo no lado}, com alerta de
    estudantes parados.
  \item Posso abrir o relatório de cada estudante e \enquote{Apagar as sessões
    desta turma?}.
\end{itemize}
\end{historia}

\begin{historia}{H06.6.7}{Analisar a turma depois da aula}{\pronto}
\hu{Professor}{ver o que a bolha fez com a turma}{preparar a aula seguinte}
\aceite
\begin{itemize}
  \item A matriz \enquote{Conclusão $\times$ lado em que o feed nasceu};
    \enquote{A saída que existia} (quem usou a busca); e \enquote{Estudante por
    estudante}.
  \item \enquote{Baixar CSV}.
\end{itemize}
\pendencia{a IA desta missão (geração e reações) ainda não conta no limite
diário de IA (\id{H09.4.2}).}
\end{historia}

% ---------------------------------------------------------------------------
\feature{F06.7}{O Manifesto da IA --- Nossos direitos em código}{\pronto}

\textbf{8º ano EF, Unidade 4} (fecha o ano) \textperiodcentered{} imita um
assistente de IA \textperiodcentered{} uma aula \textperiodcentered{} IA só no
preparo.

\textbf{Proposta.} Identificar as questões éticas de uma IA na escola e propor
uma regra que garanta direitos. Não existe resposta certa.

\begin{historia}{H06.7.1}{Personalizar o caso com um dilema da escola}{\pronto}
\hu{Professor}{trocar o caso da missão por um dilema real da minha
escola}{aproximar a discussão da turma}
\aceite
\begin{itemize}
  \item É opcional: gero com IA, leio as réplicas da máquina e aprovo; sem
    aprovação, nada chega à turma.
  \item Se eu não fizer nada, roda o pacote pronto da Unooria.
\end{itemize}
\end{historia}

\begin{historia}{H06.7.2}{Auditar a máquina}{\pronto}
\hu{Estudante do 8º ano}{questionar o sistema que a escola vai comprar}{
descobrir que algoritmo não é neutro}
\aceite
\begin{itemize}
  \item Três questionamentos em cartões prontos, sem digitação livre; preciso
    disparar pelo menos 2, e a mesma pergunta não entra duas vezes.
  \item No segundo, a máquina admite que herda punições do passado.
  \item Amara só aparece se eu ficar 2 minutos parado.
\end{itemize}
\end{historia}

\begin{historia}{H06.7.3}{Fixar uma cláusula no manifesto}{\pronto}
\hu{Estudante do 8º ano}{escolher a regra que eu exigiria}{ver a minha ideia
num documento da turma}
\aceite
\begin{itemize}
  \item Quatro cláusulas: transparência dos dados, auditoria de viés, veto
    humano ou anonimização. A escolha é irreversível.
  \item Sem ter contestado a máquina duas vezes, não consigo fixar.
  \item O sistema monta o \enquote{Manifesto de Direitos Digitais dos
    Estudantes} com a minha cláusula em destaque, que posso copiar.
\end{itemize}
\end{historia}

\begin{historia}{H06.7.4}{Escrever o meu percurso sem expor dado pessoal}{\parcial}
\hu{Estudante do 8º ano}{registrar o que aprendi}{fechar a missão com as minhas
palavras}
\aceite
\begin{itemize}
  \item Duas palavras-chave e um parágrafo de 20 a 800 caracteres.
  \item Se eu escrever telefone, CPF, e-mail, CEP, cartão ou perfil de rede
    social, o texto é recusado, não é guardado, e a tela diz o que tirar.
  \item Selo \enquote{Guardião de Direitos Digitais} e episódio de fechamento.
\end{itemize}
\pendencia{os áudios (alerta de Amara e episódio) não foram gravados --- o
botão de ouvir aparece desabilitado, com o motivo. E o texto do percurso ainda
é visível à gestão (\id{H07.1.4}).}
\end{historia}

\begin{historia}{H06.7.5}{Pausar para o debate e ver a turma}{\pronto}
\hu{Professor}{parar a turma para debater e ver quais cláusulas foram
escolhidas}{discutir as escolhas sem transformá-las em acerto}
\aceite
\begin{itemize}
  \item \enquote{Pausar para Debate} (opcional) e o inverso.
  \item A distribuição das 4 cláusulas aparece em contagem, sem taxa de
    acerto.
  \item \enquote{Zerar a turma}.
\end{itemize}
\end{historia}

% ---------------------------------------------------------------------------
\feature{F06.8}{unooGPT --- Missão Tempestade}{\pronto}

\textbf{2ª série EM, Unidade 1} \textperiodcentered{} imita um assistente de
IA durante toda a missão \textperiodcentered{} uma aula \textperiodcentered{}
IA conversa com o estudante.

\textbf{Proposta.} Letramento de dados: perguntar melhor, perceber que fontes
discordam e decidir sob incerteza, sem confundir \enquote{não sabemos} com
\enquote{zero}. Nenhuma pergunta produz a resposta certa.

\begin{historia}{H06.8.1}{Dar a largada para a turma}{\pronto}
\hu{Professor}{controlar o relógio da turma}{que todos comecem e terminem
juntos}
\aceite
\begin{itemize}
  \item O cronômetro fica parado até \enquote{Dar a largada}.
  \item Também: +2 min, encerrar, modo livre e redefinir.
\end{itemize}
\end{historia}

\begin{historia}{H06.8.2}{Investigar perguntando}{\pronto}
\hu{Estudante do Médio}{fazer perguntas ao assistente sobre a tempestade que
chega}{decidir para onde mandar as equipes de resgate}
\aceite
\begin{itemize}
  \item Tenho 3 perguntas em 5 minutos para um \enquote{Cientista de Dados} que
    consulta 8 fontes, cada uma com um viés que eu não vejo.
  \item O mapa só acende o que as fontes consultadas cobrem; bairro em branco
    quer dizer \enquote{não sabemos}.
  \item Cada resposta leva de 10 a 30 segundos e usa o meu limite diário de
    IA.
\end{itemize}
\end{historia}

\begin{historia}{H06.8.3}{Decidir e ver o resultado como faixa}{\pronto}
\hu{Estudante do Médio}{mandar as equipes e ver o que aconteceu}{entender que
uma decisão sob incerteza tem vários desfechos possíveis}
\aceite
\begin{itemize}
  \item Duas equipes para quatro bairros; justifico com as fontes que abri ou
    com \enquote{foi intuição}.
  \item O resultado aparece como uma faixa de 20 repetições da mesma decisão,
    mais um fato concreto e a lista do que eu nunca abri. Sem pontuação.
  \item Posso recomeçar.
\end{itemize}
\end{historia}

\begin{historia}{H06.8.4}{Ler a turma e cada estudante}{\pronto}
\hu{Professor}{ver como cada estudante usou o tempo e as fontes}{discutir o
processo de decisão, não o placar}
\aceite
\begin{itemize}
  \item A folha da turma mostra a conduta, sem coluna de \enquote{resgatados}.
  \item O relatório de cada estudante mostra \enquote{Rodadas enfrentadas},
    \enquote{Até a primeira fala}, \enquote{Mudanças de destino},
    \enquote{O que a missão mediu}, \enquote{Como usou o tempo} e a conversa
    inteira.
  \item O guia traz a tabela dos 8 vieses, com o aviso de não projetá-la.
\end{itemize}
\pendencia{a ficha e o guia não passaram por teste em sala; estão marcados
como ponto de partida para revisão pedagógica.}
\end{historia}

% ---------------------------------------------------------------------------
\feature{F06.9}{Vila do Lume --- A Noite Longa}{\pronto}

\textbf{Para qualquer ano} (6 a 13 anos) \textperiodcentered{} jogo 3D próprio,
em primeira pessoa \textperiodcentered{} uma aula.

\textbf{Proposta.} Na camada visível, reconhecer um golpe que não tem os sinais
que se costuma ensinar a procurar. Na camada invisível, um retrato de valores
feito a partir de escolhas que custam algo --- e que define de que lado cada
estudante começa no Instanooria.

\begin{historia}{H06.9.1}{Acender as fogueiras da vila}{\pronto}
\hu{Estudante}{explorar a vila e acender as oito fogueiras}{salvar a vila antes
que a noite acabe}
\aceite
\begin{itemize}
  \item Uma instrução só: \enquote{A vila apagou. Acenda as oito fogueiras
    antes que a noite acabe.}
  \item Pelo caminho há gestos opcionais (ajudar um aldeão, acender um
    lampião, gravar o nome na pedra, tocar o sino), que não valem ponto.
  \item O nome da criança nunca é gravado.
\end{itemize}
\end{historia}

\begin{historia}{H06.9.2}{Receber mensagens no meio do jogo}{\pronto}
\hu{Estudante}{receber mensagens enquanto jogo}{descobrir, no debate, em quais
eu confiei}
\aceite
\begin{itemize}
  \item Chegam 5 mensagens sem pausar o jogo: verdadeiras e um golpe principal
    num de três níveis (grosseiro, autoridade com urgência, colega pedindo um
    passo pequeno).
  \item O nível é distribuído de forma equilibrada na turma; do 1º ao 5º ano o
    nível intermediário não é usado.
  \item O jogo não tem tela de resultado: a revelação é do professor.
\end{itemize}
\end{historia}

\begin{historia}{H06.9.3}{Ler a turma sem rotular ninguém}{\pronto}
\hu{Professor}{entender o que aconteceu com as mensagens na
turma}{conduzir o debate \enquote{Em quem dava pra confiar?}}
\aceite
\begin{itemize}
  \item Na revelação, projeto as mensagens das três versões lado a lado.
  \item Resultados: \enquote{As três bolhas nesta turma}, \enquote{O que
    aconteceu com as mensagens} e \enquote{Estudante por estudante} --- sem
    nenhum valor pessoal na lista.
  \item O relatório de uma criança abre pelo nível do golpe que ela recebeu e
    mostra \enquote{Os dez valores}; só há número quando há dois ou mais
    sinais, e \enquote{não teve a oportunidade} é diferente de \enquote{teve e
    não foi}.
  \item Partida abandonada não gera retrato, e a tela explica por quê.
  \item \enquote{Zerar a turma} (\enquote{Isso não tem volta.}) também
    redistribui as versões.
  \item A gestão vê o resumo agregado; a linha de uma criança, não.
\end{itemize}
\pendencia{o retrato de valores é uma hipótese de desenho, não um instrumento
validado. As reações dentro do jogo ainda não contam no limite diário de IA.}
\end{historia}

% ---------------------------------------------------------------------------
\feature{F06.10}{As missões que faltam}{\afazer}

A trilha tem 48 unidades (12 anos $\times$ 4). Estão publicadas 8, mais a Vila
do Lume. O 2º, o 3º e o 9º ano do Fundamental e a 1ª e a 3ª série do Médio
ainda não têm nenhuma; o 5º ano ficou sem a sua de 25 de setembro de 2026 até a
versão nova de Quem mandou isso? voltar (\id{F06.11}). As unidades
em produção aparecem para o professor como \enquote{em produção}, sem título
inventado. A próxima já decidida é
\textbf{A Seleção}, no 7º ano, Unidade 1, no lugar do Eco no Grupo.

\begin{historia}{H06.10.1}{Cada missão nova nasce completa}{\afazer}
\hu{PO}{que toda missão nova entre na trilha pelo mesmo caminho}{que a décima
missão custe menos que a nona, e nenhuma quebre uma promessa do produto}
\aceiteproposta
\begin{itemize}
  \item Tem ano e unidade, e aparece na trilha e na grade da sala.
  \item Só abre na turma em que o professor liberou.
  \item Aparece no painel da Unooria e \textbf{desliga} quando desligada
    (\id{H10.1.2}).
  \item Tem ficha e guia da aula com as nove seções, e a habilidade da BNCC.
  \item A correção é da plataforma; os resultados não têm ranking nem coluna de
    acerto.
  \item O texto livre da criança é lido só pelo professor da turma.
  \item Se usa IA: nada gerado chega à turma sem aprovação, há material pronto
    se a IA falhar, e o uso conta no limite diário.
  \item Toda a interface em português, exceto os termos do aplicativo imitado.
  \item Tem teste automatizado que percorre a missão como o estudante e como o
    professor.
\end{itemize}
\end{historia}

\begin{historia}{H06.10.2}{A Seleção (7º ano, Unidade 1)}{\afazer}
\hu{Estudante do 7º ano}{uma missão na rede social de fotos para a Unidade 1
do 7º ano}{que a trilha do 7º ano siga o material do 1º bimestre}
\aceiteproposta
\begin{itemize}
  \item Substitui por completo o Eco no Grupo; os dados do Eco são limpos antes.
  \item As imagens dos posts são geradas no preparo do professor, com
    alternativa pronta se a geração falhar.
  \item O detalhamento vem do documento de requisitos da missão.
\end{itemize}
\end{historia}

\begin{historia}{H06.10.3}{Quem mandou isso?, versão nova (5º ano, Unidade 1)}{\parcial}
\hu{Estudante do 5º ano}{a versão nova da missão do grupo de mensagens na
Unidade 1 do 5º ano}{aprender a conferir antes de repassar, agora também em
equipe}
\aceite
\begin{itemize}
  \item Substitui por completo a primeira versão, que saiu da trilha com os
    dados dela em 25 de setembro de 2026, e volta ao mesmo lugar da trilha.
  \item As mensagens do grupo se dividem em verdadeiras ou falsas e em
    emocionais ou neutras, e o estudante escolhe exatamente três para repassar.
  \item Depois de escolher, ele diz numa enquete o que pesou na escolha, e
    fecha com a regra \enquote{Se \ldots, então \ldots} e duas perguntas da
    Nina.
  \item A missão registra o que cada estudante demonstrou das habilidades
    (\id{E14}) e aceita equipes num computador dividido (\id{E13}).
  \item O detalhe, história por história, está em \id{F06.11}.
\end{itemize}
\pendencia{o teste que percorre a missão como o estudante e como o professor
foi escrito e ainda não rodou; a régua das habilidades aguarda a validação da
equipe pedagógica; e os textos escritos pela equipe aguardam o PO.}
\end{historia}

% ===========================================================================
%  E07 — Acompanhar a aprendizagem e a rede
% ===========================================================================
\chapter{E07 \textperiodcentered{} Acompanhar a aprendizagem e a rede}
\label{ep:e07}

\begin{fichaepico}
\begin{campos}
  \item[Objetivo] Cada nível enxerga o que precisa para agir: o professor, a
    turma e cada estudante; a diretoria, a adoção nas escolas; a rede, o
    agregado --- nunca a criança.
  \item[Para quem] Professor; Coordenação regional; Administrador do contrato.
  \item[Valor] É o que transforma uma missão jogada em aula avaliável, e o que
    a secretaria mostra quando presta contas.
  \item[Estado] \parcial{} --- os resultados da turma estão prontos; os
    relatórios da rede têm ajustes pendentes.
\end{campos}
\end{fichaepico}

\begin{tabularx}{\linewidth}{@{}L{16mm}X L{20mm}@{}}
\toprule
\headrow \thd{ID} & \thd{Feature} & \thd{Estado}\\
\midrule
\id{F07.1} & Resultados da turma & \parcial\\
\zebra \id{F07.2} & Relatório da diretoria & \parcial\\
\id{F07.3} & Análise de produto na rede & \parcial\\
\bottomrule
\end{tabularx}

% ---------------------------------------------------------------------------
\feature{F07.1}{Resultados da turma}{\parcial}

As regras abaixo valem para os resultados de todas as missões. O que é
específico de cada uma está no épico \id{E06}.

\begin{historia}{H07.1.1}{Ler a turma sem ranking}{\pronto}
\hu{Professor}{ver como a turma se saiu na missão}{planejar o debate e a aula
seguinte}
\aceite
\begin{itemize}
  \item Em \enquote{Resultados da turma}, escolho a turma e vejo a folha dela.
  \item Nenhuma folha tem coluna de pontos, de acerto ou ordena crianças.
  \item Sem turma escolhida, a tela explica o que fazer.
  \item A condução ao vivo da aula (largada, Aula 2, debate, revelação) fica na
    mesma tela dos resultados.
\end{itemize}
\end{historia}

\begin{historia}{H07.1.2}{Zerar a turma para repetir a missão}{\pronto}
\hu{Professor}{apagar os dados de uma missão na minha turma}{rodar a missão de
novo com a mesma turma}
\aceite
\begin{itemize}
  \item \enquote{Zerar a turma} pede confirmação, porque não tem volta.
  \item Só aquela missão, naquela turma, é zerada; nenhuma outra missão é
    tocada.
\end{itemize}
\end{historia}

\begin{historia}{H07.1.3}{Ler um estudante sem rotulá-lo}{\pronto}
\hu{Professor}{abrir o relatório de um estudante}{entender o caminho dele na
missão}
\aceite
\begin{itemize}
  \item Existe hoje no unooGPT, na Vila do Lume e no Instanooria (o dossiê).
  \item Uma criança por vez; o relatório nunca afirma omissão a partir de
    ausência de dado, e um valor com uma única medida não recebe número.
  \item A Coordenação e a rede não têm acesso ao relatório de uma criança.
\end{itemize}
\end{historia}

\begin{historia}{H07.1.4}{O texto da criança é lido só pelo professor dela}{\parcial}
\hu{família do estudante}{ter certeza de que o que a criança escreve na
missão é lido só pelo professor da turma}{que ela possa escrever com
sinceridade}
\aceite
\begin{itemize}
  \item O texto livre da criança aparece só para o professor da turma.
  \item A Coordenação, a regional e a rede veem apenas \enquote{entregou / não
    entregou}.
\end{itemize}
\pendencia{no Rastro do Boato e no Manifesto da IA, o texto do \enquote{Meu
percurso} ainda é visível à gestão. A versão nova de Quem mandou isso? segue a
regra: a regra \enquote{Se \ldots, então \ldots} da criança só aparece para quem
dá aula na turma (\id{H06.11.5}). Ver a decisão \id{D03}.}
\end{historia}

% ---------------------------------------------------------------------------
\feature{F07.2}{Relatório da diretoria}{\parcial}

\begin{historia}{H07.2.1}{Acompanhar a adoção na minha regional}{\parcial}
\hu{Coordenação regional}{ver quais escolas estão usando as missões}{apoiar
as escolas que estão atrás}
\aceite
\begin{itemize}
  \item \enquote{Relatório} mostra o consolidado por missão e por escola, com
    \enquote{Turmas abertas}, \enquote{Turmas encerradas} e
    \enquote{Estudantes com registro}.
  \item Só a própria regional, sem seletor.
  \item Sem nome de criança, sem bem-estar, sem ranking.
  \item \enquote{Baixar CSV}.
\end{itemize}
\pendencia{a coluna \enquote{Missão} mostra o código interno (por exemplo
\enquote{exp08ef1b}) em vez do nome da missão.}
\end{historia}

\begin{historia}{H07.2.2}{Comparar as regionais}{\pronto}
\hu{Administrador do contrato}{ver o mesmo relatório para qualquer
regional}{acompanhar a rede inteira}
\aceite
\begin{itemize}
  \item O mesmo relatório, com um seletor de regional.
\end{itemize}
\end{historia}

% ---------------------------------------------------------------------------
\feature{F07.3}{Análise de produto na rede}{\parcial}

\begin{historia}{H07.3.1}{Saber se a missão funciona entre turmas}{\parcial}
\hu{Administrador do contrato}{comparar o envolvimento e a eficácia da missão
entre turmas}{saber onde ela funciona e onde precisa de apoio}
\aceite
\begin{itemize}
  \item \enquote{Análise de produto} mostra, por atividade, Sessões,
    Envolvimento, Tempo de feed, Inversão, \enquote{Achou a busca},
    \enquote{Conteúdo morto} e \enquote{Imagens ok}, com o detalhe de cada
    atividade.
  \item Nenhum nome de estudante.
\end{itemize}
\pendencia{só se chega a esta tela digitando o endereço (não há item de
menu); cobre só o Instanooria; a escola é digitada à mão; e o CSV existe mas
não tem botão na tela.}
\end{historia}

\begin{historia}{H07.3.2}{Ver os números do dia}{\afazer}
\hu{Administrador do contrato}{ver quantas turmas e estudantes usaram a
plataforma hoje}{acompanhar a implantação sem esperar o fim do mês}
\aceiteproposta
\begin{itemize}
  \item Um painel com turmas com missão aberta, estudantes que entraram e
    missões concluídas no dia, por regional.
  \item Nenhum dado individual.
\end{itemize}
\end{historia}

% ===========================================================================
%  E08 — Ferramentas do professor
% ===========================================================================
\chapter{E08 \textperiodcentered{} Ferramentas do professor}
\label{ep:e08}

\begin{fichaepico}
\begin{campos}
  \item[Objetivo] Tirar do professor o trabalho repetitivo do planejamento e
    cuidar de quem cuida da turma.
  \item[Para quem] Professor.
  \item[Valor] São três produtos que a rede pode contratar separadamente do
    Entre Redes: Gerador de Plano de Aula, Planejamento de Ensino e Como você
    está?.
  \item[Estado] \pronto.
\end{campos}
\end{fichaepico}

\begin{tabularx}{\linewidth}{@{}L{16mm}X L{20mm}@{}}
\toprule
\headrow \thd{ID} & \thd{Feature} & \thd{Estado}\\
\midrule
\id{F08.1} & Gerador de Plano de Aula & \pronto\\
\zebra \id{F08.2} & Planejamento de Ensino & \pronto\\
\id{F08.3} & Como você está? & \pronto\\
\bottomrule
\end{tabularx}

\begin{nota}
\textbf{Dois planejadores, dois públicos.} O Gerador de Plano de Aula é para a
rede que já tem currículo próprio e material didático adotado: ele parte desse
material. O Planejamento de Ensino é para quem parte da BNCC e quer descer até
a aula de cada semana. Os dois convivem; um não substitui o outro.
\end{nota}

% ---------------------------------------------------------------------------
\feature{F08.1}{Gerador de Plano de Aula}{\pronto}

\begin{historia}{H08.1.1}{Gerar o plano anual a partir do material da rede}{\pronto}
\hu{Professor}{gerar um plano anual da minha disciplina e série}{não começar
o ano com a página em branco}
\aceite
\begin{itemize}
  \item \enquote{Meus planos de aula} \seta{} \enquote{Novo plano de aula}:
    escolho disciplina e série.
  \item Aparecem só as séries com material cadastrado pela rede; disciplinas
    sem material aparecem como \enquote{em breve}.
  \item A geração leva alguns minutos, com dicas na tela enquanto espero;
    posso fechar a tela, e ela continua.
  \item O plano começa pela fundamentação e pelos objetivos do ano, mostra o
    percurso do ano numa tabela (a pergunta, a aprendizagem central e o produto de
    cada bimestre) e, em cada bimestre, traz a pergunta norteadora, os conteúdos, as
    habilidades, as estratégias, as formas de avaliação e o produto de síntese.
  \item As formas de avaliação dizem o que o estudante mostra que aprendeu, e
    cada uma se liga a uma habilidade do bimestre.
  \item Um erro aparece como uma frase que eu entendo.
  \item Usa o meu limite diário de IA.
\end{itemize}
\pendencia{depende de a rede ter enviado o material de referência.}
\end{historia}

\begin{historia}{H08.1.2}{Editar e levar o plano para o Word}{\pronto}
\hu{Professor}{ajustar o plano gerado e exportar}{entregar à coordenação no
formato que ela pede}
\aceite
\begin{itemize}
  \item O plano abre como documento editável: salvar (Ctrl/Cmd+S), renomear,
    excluir.
  \item \enquote{Restaurar original} volta à versão que a IA gerou, que fica
    guardada intacta.
  \item Exportar em .doc (abre no Word).
  \item O plano é do professor: não fica público.
\end{itemize}
\end{historia}

% ---------------------------------------------------------------------------
\feature{F08.2}{Planejamento de Ensino}{\pronto}

\begin{historia}{H08.2.1}{Declarar a minha regência}{\pronto}
\hu{Professor}{informar o que eu leciono e a minha grade}{que o planejamento
tenha as datas reais do meu ano}
\aceite
\begin{itemize}
  \item Informo componente, série, ano letivo e data de início, grade de
    horários, minutos por aula e divisão do ano (bimestre, trimestre ou
    semestre).
  \item A carga horária sugerida vem da mediana de redes reais e aparece como
    sugestão, que posso mudar.
  \item Hora-aula não é hora-relógio: a conta respeita os minutos por aula.
\end{itemize}
\end{historia}

\begin{historia}{H08.2.2}{Distribuir as habilidades da BNCC no ano}{\pronto}
\hu{Professor}{colocar cada habilidade da BNCC num período do ano}{garantir
que nenhuma fique de fora}
\aceite
\begin{itemize}
  \item Cada habilidade vai para um período.
  \item No Ensino Médio, a tela avisa quando as habilidades são por área.
  \item Nenhuma habilidade fora do período entra no plano daquele período.
\end{itemize}
\end{historia}

\begin{historia}{H08.2.3}{Descer do plano anual à aula da semana}{\pronto}
\hu{Professor}{ter o plano anual e o plano de cada semana}{chegar à sala
sabendo o que fazer em cada aula}
\aceite
\begin{itemize}
  \item O plano anual é gerado e eu posso reescrevê-lo.
  \item As 40 semanas são planejadas uma por vez, com tema, objetivo, quatro
    momentos, recursos, avaliação, tarefa e adaptações.
  \item As datas vêm do calendário real: feriado remove a aula. A IA preenche
    o conteúdo, mas nunca decide quantas aulas cabem nem inventa código de
    habilidade.
  \item Marco \enquote{AULA CUMPRIDA} e escrevo observações.
\end{itemize}
\end{historia}

\begin{historia}{H08.2.4}{Exportar a semana, a quinzena ou o bimestre}{\pronto}
\hu{Professor}{baixar o planejamento no recorte que a coordenação
pede}{entregar sem reformatar}
\aceite
\begin{itemize}
  \item \enquote{Baixar esta semana}, \enquote{Esta e a próxima (quinzena)} e
    \enquote{Baixar o Nº Bimestre inteiro}, em .doc; ou imprimir.
  \item Semanas cuja geração falhou não entram no documento.
  \item Posso excluir o planejamento daquele componente e série.
\end{itemize}
\end{historia}

% ---------------------------------------------------------------------------
\feature{F08.3}{Como você está?}{\pronto}

Um acompanhamento curto do bem-estar do próprio professor: três frases por dia
de um questionário reconhecido (DASS-21, que olha depressão, ansiedade e
estresse), e sete dias fecham o retrato da semana. Leva menos de um minuto.

\begin{historia}{H08.3.1}{Responder hoje em menos de um minuto}{\pronto}
\hu{Professor}{responder três frases sobre como estou}{acompanhar o meu
bem-estar sem tomar tempo do dia}
\aceite
\begin{itemize}
  \item O cartão no Início mostra se já respondi hoje.
  \item \enquote{Responder hoje}: três frases numa escala e, se eu quiser, um
    texto livre; \enquote{Enviar respostas}; recebo uma mensagem de
    acolhimento.
  \item Uma resposta por dia.
  \item A tela diz com clareza que nenhuma pessoa da escola ou da rede lê o
    texto livre, e que uma IA lê.
  \item As frases do questionário nunca são reescritas.
\end{itemize}
\end{historia}

\begin{historia}{H08.3.2}{Ver o retrato da minha semana}{\pronto}
\hu{Professor}{ver o resultado ao completar sete dias}{perceber como estou ao
longo do tempo}
\aceite
\begin{itemize}
  \item \enquote{Falta N dias para o retrato da semana}; no sétimo dia,
    \enquote{Retrato da sua semana} com a pontuação de cada dimensão e a faixa
    correspondente.
  \item Um ciclo que fica 21 dias sem fechar expira, sem pontuação.
  \item \enquote{Seu histórico} mostra os ciclos concluídos e o que escrevi em
    cada dia.
\end{itemize}
\end{historia}

\begin{historia}{H08.3.3}{Receber apoio quando preciso}{\pronto}
\hu{Professor que está passando por um momento difícil}{receber orientação de
onde buscar ajuda}{não ficar sozinho com isso}
\aceite
\begin{itemize}
  \item Se as respostas ou o texto indicam risco, aparece uma mensagem
    \emph{fixa} com o CVV (188), o SAMU (192) e a orientação de procurar o
    CAPS ou a UBS.
  \item Nesse caso, a IA não é chamada.
\end{itemize}
\end{historia}

\begin{historia}{H08.3.4}{Só meu, e apagável}{\pronto}
\hu{Professor}{ter certeza de que ninguém da gestão vê o meu bem-estar, e poder
apagar tudo}{responder com sinceridade}
\aceite
\begin{itemize}
  \item Não existe tela de gestão, número agregado nem exportação desse dado.
  \item \enquote{Apagar todas as suas respostas} remove tudo, de todos os
    ciclos --- e funciona mesmo se a rede desligar o produto.
  \item O dado de bem-estar nunca entra no perfil portátil.
\end{itemize}
\end{historia}

% ===========================================================================
%  E09 — Inteligência artificial com limites
% ===========================================================================
\chapter{E09 \textperiodcentered{} Inteligência artificial com limites}
\label{ep:e09}

\begin{fichaepico}
\begin{campos}
  \item[Objetivo] Usar IA onde ela poupa trabalho do professor ou enriquece a
    missão --- sem que ela decida nota, conheça a criança ou seja condição
    para a aula acontecer. A única vez em que ela dá um nível é a exceção
    declarada da A Seleção, e quem decide continua sendo o professor
    (\id{H09.3.1}).
  \item[Para quem] Professor e Estudante (quem usa); rede (quem precisa
    confiar).
  \item[Valor] IA na escola pública só é aceitável com garantia e com custo
    previsível. Este épico é essa garantia.
  \item[Estado] \parcial{} --- as garantias estão prontas; duas missões ainda
    usam IA fora do limite diário.
\end{campos}
\end{fichaepico}

\begin{tabularx}{\linewidth}{@{}L{22mm}X X@{}}
\toprule
\headrow \thd{Onde} & \thd{O que a IA faz} & \thd{O que ela nunca faz}\\
\midrule
Preparo da aula & Adapta legendas, mensagens, casos e mídia ao tema do
  professor (Foto Conta, Manifesto, Instanooria, Quem mandou isso?, A Seleção
  --- o tema, os 12 posts, as imagens e as três falas do podcast que dependem
  do tema); gera a voz do podcast da A Seleção. &
  Chegar à turma sem aprovação.\\
\zebra Durante a missão & Responde ao estudante no unooGPT; reações de
  personagens no Instanooria e na Vila; na A Seleção, lê a legenda publicada e
  dá o nível dela numa habilidade da BNCC. & Dar nota, mostrar nível à criança
  ou receber o nome dela. O nível da legenda da A Seleção é a exceção
  declarada: só o professor o vê, e ele confere e ajusta.\\
Planejamento & Gera o plano anual e o das semanas. & Decidir quantas aulas
  cabem ou inventar habilidade da BNCC.\\
\zebra Bem-estar & Escreve o texto em volta (abertura, acolhimento). &
  Reescrever o questionário, ou falar quando há sinal de risco.\\
\bottomrule
\end{tabularx}

\begin{tabularx}{\linewidth}{@{}L{16mm}X L{20mm}@{}}
\toprule
\headrow \thd{ID} & \thd{Feature} & \thd{Estado}\\
\midrule
\id{F09.1} & Nada gerado chega à turma sem o professor & \pronto\\
\zebra \id{F09.2} & Nenhuma aula depende da IA & \pronto\\
\id{F09.3} & A IA não dá nota nem conhece a criança & \pronto\\
\zebra \id{F09.4} & Limite diário por perfil & \parcial\\
\bottomrule
\end{tabularx}

% ---------------------------------------------------------------------------
\feature{F09.1}{Nada gerado chega à turma sem o professor}{\pronto}

\begin{historia}{H09.1.1}{Aprovar o que a IA preparou}{\pronto}
\hu{Professor}{revisar e aprovar tudo o que a IA preparou para a minha
aula}{responder pelo que a turma vê}
\aceite
\begin{itemize}
  \item Nas missões com preparo por IA, nada chega aos estudantes sem
    \enquote{Aprovar} (ou \enquote{Aprovar e publicar}); no Instanooria, sem
    \enquote{Ativar}.
  \item Posso revisar item a item, editar, regerar ou excluir.
  \item A tela diz: \enquote{nada gerado por inteligência artificial chega aos
    estudantes sem aprovação}.
\end{itemize}
\end{historia}

% ---------------------------------------------------------------------------
\feature{F09.2}{Nenhuma aula depende da IA}{\pronto}

\begin{historia}{H09.2.1}{Dar a aula mesmo se a IA falhar}{\pronto}
\hu{Professor}{que a aula aconteça mesmo se a IA estiver fora do ar ou sem
saldo}{não perder a aula por um problema que não é meu}
\aceite
\begin{itemize}
  \item Sem IA ou sem saldo, entra o material pronto da Unooria, completo
    (Foto Conta, Manifesto, Quem mandou isso?, A Seleção).
  \item Na A Seleção, sem imagem gerada o post mostra a descrição da cena; sem
    a voz do podcast, a transcrição continua na tela; e, se a leitura da
    legenda falhar, nada é marcado e a enquete continua medindo a habilidade.
  \item No Instanooria, se as reações ao vivo falharem, simplesmente ninguém
    responde --- a missão segue.
  \item No bem-estar, entra um texto fixo e o acompanhamento nunca deixa de
    abrir.
  \item O estudante nunca espera uma geração, exceto na conversa do unooGPT,
    que é a própria missão.
\end{itemize}
\end{historia}

% ---------------------------------------------------------------------------
\feature{F09.3}{A IA não dá nota nem conhece a criança}{\pronto}

\begin{historia}{H09.3.1}{Nota é da plataforma, não da IA}{\pronto}
\hu{rede}{que nenhuma nota, correção ou classificação venha da IA, a não ser
por uma exceção declarada e cercada de garantias}{que o resultado seja justo e
explicável}
\aceite
\begin{itemize}
  \item \textbf{A regra:} gabarito, correção e classificação são calculados
    pela plataforma com regras fixas, e a IA só redige textos.
  \item \textbf{A exceção declarada, e só ela:} na A Seleção, a IA lê a
    legenda que o estudante publica nos Stories e dá o nível dela na EF07CO11
    (atingiu, em processo ou não atingiu), porque só uma leitura do texto diz se
    a legenda explica o critério da escolha (\id{H06.12.5}).
  \item As garantias da exceção: a IA trabalha no modo mais previsível que tem,
    com uma tentativa só, de poucos segundos, e a leitura conta no limite
    diário do estudante. Se qualquer coisa falhar, nada é marcado e a
    habilidade fica \enquote{Não avaliado}; nada do que a IA escreve é guardado.
  \item A criança não vê o nível. O professor o vê no painel de habilidades,
    confere e ajusta, e a primeira enquete da Amara mede a mesma habilidade
    (vale a melhor das duas).
  \item Nas outras missões vale a regra: nenhuma usa a IA para dar nível.
  \item O texto da IA é conferido antes de aparecer (tamanho, termos
    proibidos); se não passar, é refeito ou substituído pelo texto pronto.
\end{itemize}
\end{historia}

\begin{historia}{H09.3.2}{A IA não sabe quem é a criança}{\pronto}
\hu{família do estudante}{que a IA não receba dados que identifiquem a
criança}{que o uso de IA não exponha o meu filho}
\aceite
\begin{itemize}
  \item A IA recebe no máximo a série da turma (\enquote{4º ano}), nunca nome,
    escola ou pontuação.
  \item As instruções da IA nunca vão para o aparelho do estudante.
\end{itemize}
\pendencia{a IA processa fora do Brasil. Isso é aceito, mas precisa constar
no contrato e no termo de tratamento de dados de cada rede (ver o
Capítulo~\ref{cap:backlog}).}
\end{historia}

% ---------------------------------------------------------------------------
\feature{F09.4}{Limite diário por perfil}{\parcial}

\begin{historia}{H09.4.1}{Saber que cheguei ao limite do dia}{\pronto}
\hu{Estudante ou Professor}{uma mensagem clara quando atingi o limite de IA do
dia}{entender o que aconteceu e quando volta}
\aceite
\begin{itemize}
  \item O saldo é conferido \emph{antes} de usar a IA; sem saldo, aparece
    \enquote{limite do dia atingido} e nada é consumido.
  \item O limite é por pessoa e por perfil (estudante ou professor), definido
    pela Unooria (\id{H10.4.1}).
\end{itemize}
\pendencia{o contador zera à meia-noite no horário universal, que no Brasil é
21h. Ver a decisão \id{D04}.}
\end{historia}

\begin{historia}{H09.4.2}{Toda IA conta no limite}{\afazer}
\hu{Administrador Unooria}{que todo uso de IA da plataforma conte no limite
diário}{que o custo de cada rede seja previsível}
\aceiteproposta
\begin{itemize}
  \item A geração e as reações do Instanooria e as reações da Vila do Lume
    passam a contar no limite, como as demais.
  \item O limite do professor comporta a preparação de uma atividade completa
    do Instanooria.
\end{itemize}
\end{historia}

\begin{historia}{H09.4.3}{Ser avisado antes do limite estourar}{\afazer}
\hu{Administrador Unooria}{ser avisado quando o consumo de IA de uma rede se
aproximar do previsto}{ajustar antes que alguém fique sem IA no meio da aula}
\aceiteproposta
\begin{itemize}
  \item Aviso quando o consumo do dia chega a 80\% do previsto para a rede.
\end{itemize}
\end{historia}

% ===========================================================================
%  E10 — Painel da Unooria
% ===========================================================================
\chapter{E10 \textperiodcentered{} Painel da Unooria}
\label{ep:e10}

\begin{fichaepico}
\begin{campos}
  \item[Objetivo] A equipe da Unooria decide o que existe em cada rede, qual
    versão de cada missão está no ar e quanto de IA cada perfil pode usar ---
    e recebe os interessados do site.
  \item[Para quem] Administrador Unooria. O Administrador do contrato só lê o
    histórico.
  \item[Valor] É o que permite vender produtos separadamente, desligar uma
    missão com problema em minutos e voltar uma versão sem esperar o time.
  \item[Estado] \pronto{} --- faltam desligar por turma ou por prazo e alguns
    avisos.
\end{campos}
\end{fichaepico}

\begin{tabularx}{\linewidth}{@{}L{16mm}X L{20mm}@{}}
\toprule
\headrow \thd{ID} & \thd{Feature} & \thd{Estado}\\
\midrule
\id{F10.1} & Ligar e desligar produtos e missões & \pronto\\
\zebra \id{F10.2} & Histórico que a rede também lê & \pronto\\
\id{F10.3} & Versões das missões & \pronto\\
\zebra \id{F10.4} & Uso de IA & \pronto\\
\id{F10.5} & Prospecção & \pronto\\
\zebra \id{F10.6} & Experiências arquivadas & \pronto\\
\bottomrule
\end{tabularx}

\begin{nota}
\textbf{Por que só a Unooria opera.} O Administrador do contrato pode quase
tudo na rede dele --- mas não pode ligar o que a rede não contratou. É a única
capacidade que a Unooria tem e a rede não.
\end{nota}

% ---------------------------------------------------------------------------
\feature{F10.1}{Ligar e desligar produtos e missões}{\pronto}

\begin{historia}{H10.1.1}{Desligar um produto ou uma missão}{\pronto}
\hu{Administrador Unooria}{desligar um produto ou uma missão numa rede, numa
regional ou numa escola}{respeitar o que foi contratado ou tirar do ar uma
missão com problema}
\aceite
\begin{itemize}
  \item Em \enquote{Recursos}, escolho o Escopo: \enquote{Rede (contrato
    inteiro)}, uma regional ou uma escola.
  \item Vejo os cinco produtos e, dentro de cada um, as missões (o agente de IA
    do professor não tem missões dentro). Cada
    interruptor tem três estados: \enquote{Ligado}, \enquote{Desligado} e
    \enquote{Herda de cima} (na rede, \enquote{Padrão (ligado)}).
  \item Por padrão tudo está ligado; o escopo mais específico vence (escola,
    depois regional, depois rede).
  \item O produto é um teto: desligado, as missões dentro dele ficam
    desligadas, com o motivo ao lado.
  \item Desligar exige motivo; na rede, o aviso diz \enquote{Isto vale para a
    REDE INTEIRA}.
  \item A missão desligada sai da grade da turma (não fica com cadeado), e
    para a criança a mensagem é a mesma de missão ainda não aberta.
  \item Religar devolve a sala como estava: as liberações do professor nunca
    são apagadas.
  \item Com o produto desligado, cada pessoa continua podendo apagar os
    próprios dados.
\end{itemize}
\end{historia}

\begin{historia}{H10.1.2}{Missão nova já nasce com interruptor}{\pronto}
\hu{Administrador Unooria}{que toda missão ou ferramenta nova apareça sozinha
no painel}{nunca ter uma missão no ar que eu não consiga desligar}
\aceite
\begin{itemize}
  \item A missão nova aparece dentro do produto a que pertence, sem ninguém
    precisar cadastrá-la à mão.
  \item O teste de aceite é desligá-la e ver que ela some da sala e do
    estudante.
\end{itemize}
\end{historia}

\begin{historia}{H10.1.3}{Desligar por turma ou por prazo}{\afazer}
\hu{Administrador Unooria}{desligar um recurso só para uma turma, ou até uma
data}{atender pilotos e contratos temporários sem mexer na escola inteira}
\aceiteproposta
\begin{itemize}
  \item O Escopo passa a incluir turma.
  \item Ao desligar, posso informar até quando; na data, o interruptor volta a
    herdar de cima e isso fica no histórico.
\end{itemize}
\end{historia}

\begin{historia}{H10.1.4}{Desligar o agente de IA do professor numa rede}{\parcial}
\hu{Administrador Unooria}{ligar e desligar o agente de IA do professor como
qualquer produto}{oferecê-lo só a quem contratou, ou tirá-lo do ar num
problema}
\aceite
\begin{itemize}
  \item O agente aparece em \enquote{Recursos} como um produto, com o mesmo
    interruptor de três estados e o mesmo escopo (rede, regional, escola).
  \item Desligado, o botão do agente some das telas do professor naquela
    escola, e nenhuma conversa nova nem mensagem passa.
  \item Mesmo desligado, o professor continua podendo apagar as próprias
    conversas.
\end{itemize}
\pendencia{o interruptor já funciona; o agente em si ainda não está no ar nos
ambientes da rede (\id{F08.4}).}
\end{historia}

% ---------------------------------------------------------------------------
\feature{F10.2}{Histórico que a rede também lê}{\pronto}

\begin{historia}{H10.2.1}{Saber quem mudou o quê, e por quê}{\pronto}
\hu{Administrador do contrato}{ver tudo o que a Unooria ligou, desligou ou
mudou na minha rede}{explicar à secretaria por que algo sumiu}
\aceite
\begin{itemize}
  \item O histórico registra quem, o quê, quando e o motivo; o nome de quem
    mudou fica como era no dia.
  \item Ninguém altera o histórico.
  \item O Administrador do contrato lê em \enquote{Recursos da rede}, só
    leitura: \enquote{Quem opera estes interruptores é a Unooria; esta tela é
    o registro.}
  \item As coordenações não veem esta tela.
\end{itemize}
\end{historia}

% ---------------------------------------------------------------------------
\feature{F10.3}{Versões das missões}{\pronto}

\begin{historia}{H10.3.1}{Voltar para a versão anterior}{\pronto}
\hu{Administrador Unooria}{voltar uma missão para a versão anterior}{resolver
em minutos um problema que apareceu numa versão nova}
\aceite
\begin{itemize}
  \item A janela Versão mostra as versões publicadas, com a atual marcada
    \enquote{(no ar)}, e \enquote{Ativar para a rede}.
  \item Só se volta para versões que já existiram; a volta restaura também as
    configurações daquela versão.
  \item A tela diz se conferiu que a versão está de fato publicada.
  \item Enquanto isso, a missão aparece com o aviso \enquote{Esta experiência
    está em rollback}.
\end{itemize}
\end{historia}

\begin{historia}{H10.3.2}{Testar uma versão nova numa escola antes}{\pronto}
\hu{Administrador Unooria}{fixar uma versão numa escola ou regional}{testar
uma novidade num piloto antes de levar para a rede}
\aceite
\begin{itemize}
  \item Fora do escopo da rede: \enquote{Fixar neste escopo} e \enquote{Soltar
    a fixação}.
\end{itemize}
\end{historia}

\begin{historia}{H10.3.3}{Ser lembrado de uma versão fixada há muito tempo}{\afazer}
\hu{Administrador Unooria}{ser avisado quando uma versão fixada ficar muito
tempo sem ser solta}{não deixar uma escola esquecida numa versão antiga}
\aceiteproposta
\begin{itemize}
  \item O painel destaca fixações com mais de 30 dias.
\end{itemize}
\end{historia}

% ---------------------------------------------------------------------------
\feature{F10.4}{Uso de IA}{\pronto}

\begin{historia}{H10.4.1}{Definir quanto de IA cada perfil pode usar}{\pronto}
\hu{Administrador Unooria}{definir o limite diário de IA dos estudantes e dos
professores}{manter o custo previsível sem travar a aula}
\aceite
\begin{itemize}
  \item Dois cartões, \enquote{Estudantes} e \enquote{Professores}, com o
    limite em \enquote{Tokens por dia} traduzido em uso real (partidas da
    Missão Tempestade, planos anuais).
  \item Valores iniciais: cerca de 3 partidas por dia para o estudante e 6 a 9
    planos anuais para o professor.
  \item Mudar exige motivo e avisa \enquote{Vale para a rede inteira}; vale na
    próxima chamada, sem esperar o time.
  \item Zero é recusado: \enquote{para fechar a IA de um produto, desligue o
    produto em Recursos}.
  \item Toda mudança vai para o histórico. A rede não vê nem altera.
\end{itemize}
\end{historia}

% ---------------------------------------------------------------------------
\feature{F10.5}{Prospecção}{\pronto}

\begin{historia}{H10.5.1}{Receber os interessados do site}{\pronto}
\hu{Administrador Unooria}{ver quem se cadastrou no site}{entrar em contato
com as redes interessadas}
\aceite
\begin{itemize}
  \item \enquote{Prospecção} lista Quando, Nome, E-mail, Telefone,
    \enquote{Rede, secretaria ou escola} e Cargo, dos mais recentes aos mais
    antigos (até 5 mil).
  \item \enquote{Atualizar} e \enquote{Baixar CSV}.
  \item Só a Unooria vê esta lista.
\end{itemize}
\end{historia}

\begin{historia}{H10.5.2}{Eliminar um cadastro a pedido}{\afazer}
\hu{Administrador Unooria}{apagar o cadastro de quem pedir}{atender o direito
de eliminação da LGPD sem depender do time técnico}
\aceiteproposta
\begin{itemize}
  \item Cada linha tem \enquote{Eliminar a pedido do titular}, com
    confirmação; o cadastro nunca é editado, só eliminado.
\end{itemize}
\pendencia{hoje o pedido é atendido por um caminho manual.}
\end{historia}

% ---------------------------------------------------------------------------
\feature{F10.6}{Experiências arquivadas}{\pronto}

Quando uma missão nova ocupa o lugar de outra na trilha, a antiga não é apagada:
ela é \textbf{arquivada}. Sai da sala, do catálogo e do painel de recursos, e os
dados que as turmas produziram ficam guardados.

\begin{historia}{H10.6.1}{Ver o que saiu do produto sem ser apagado}{\pronto}
\hu{Administrador Unooria}{ver quais missões foram arquivadas e o que ficou
guardado delas}{saber o que a plataforma já teve e responder por aqueles dados}
\aceite
\begin{itemize}
  \item \enquote{Arquivados} lista cada missão arquivada com o nome, o código de
    origem, onde ficava na trilha, a última versão, os dados guardados (turmas e
    estudantes), o motivo e a data.
  \item Os números de dados guardados são os do dia do arquivamento, e a tela
    diz isso.
  \item A lista só mostra: ninguém joga uma missão arquivada, e trazê-la de
    volta é trabalho do time técnico.
  \item Só a Unooria vê esta lista.
\end{itemize}
\end{historia}

% ===========================================================================
%  E11 — Site público e captação
% ===========================================================================
\chapter{E11 \textperiodcentered{} Site público e captação}
\label{ep:e11}

\begin{fichaepico}
\begin{campos}
  \item[Objetivo] Quem ainda não é cliente entende a proposta do Entre Redes e
    deixa contato; quem já é encontra a entrada e os documentos legais.
  \item[Para quem] Visitante (secretarias, escolas, docentes e famílias).
  \item[Valor] É a porta comercial do produto e o lugar onde a rede e a família
    leem o que a plataforma promete.
  \item[Estado] \parcial{} --- site e cadastro prontos; faltam os caminhos
    para o login e para os documentos legais.
\end{campos}
\end{fichaepico}

\begin{tabularx}{\linewidth}{@{}L{16mm}X L{20mm}@{}}
\toprule
\headrow \thd{ID} & \thd{Feature} & \thd{Estado}\\
\midrule
\id{F11.1} & Site ENTRE REDES & \parcial\\
\zebra \id{F11.2} & Cadastro de interesse & \pronto\\
\id{F11.3} & Política de Privacidade e Termos de Uso & \parcial\\
\bottomrule
\end{tabularx}

% ---------------------------------------------------------------------------
\feature{F11.1}{Site ENTRE REDES}{\parcial}

\begin{historia}{H11.1.1}{Conhecer a proposta}{\pronto}
\hu{Visitante de uma secretaria de educação}{entender o que é o Entre Redes e
como ele funciona}{decidir se vale conversar com a Unooria}
\aceite
\begin{itemize}
  \item O site tem nove seções: Programa, O desafio, Nossa resposta, O método,
    As missões, Gestão, Currículo, Docentes e Famílias.
  \item O botão \enquote{Explorar} abre o índice; os prints das missões podem
    ser ampliados.
  \item No celular, o convite encurta para \enquote{Cadastro}.
  \item Quem já está com sessão aberta e abre o site vai direto para a própria
    área.
\end{itemize}
\end{historia}

\begin{historia}{H11.1.2}{Achar a entrada e os documentos a partir do site}{\afazer}
\hu{professor da rede que chegou pelo site}{encontrar o botão de entrar e os
documentos legais}{não depender de saber o endereço da tela de entrada}
\aceiteproposta
\begin{itemize}
  \item O cabeçalho do site tem \enquote{Entrar}.
  \item O rodapé leva à Política de Privacidade e aos Termos de Uso.
\end{itemize}
\end{historia}

% ---------------------------------------------------------------------------
\feature{F11.2}{Cadastro de interesse}{\pronto}

\begin{historia}{H11.2.1}{Deixar o meu contato}{\parcial}
\hu{Visitante}{deixar os meus dados para a Unooria me procurar}{conhecer a
solução com calma}
\aceite
\begin{itemize}
  \item Os três botões \enquote{Quero conhecer a solução} abrem a mesma janela:
    Nome, E-mail profissional, Telefone / WhatsApp, \enquote{Rede, secretaria
    ou escola}, Seu cargo e a caixa \enquote{Autorizo o contato da Unooria
    sobre o ENTRE REDES}.
  \item Depois de \enquote{Enviar cadastro}: \enquote{Cadastro enviado. A
    equipe da Unooria entrará em contato.}
  \item A resposta é sempre a mesma --- inclusive para cadastro repetido ou
    barrado ---, para não revelar quem já se cadastrou.
  \item O texto exato da autorização é guardado junto com o cadastro.
  \item Há proteção contra robôs e limite de envios por hora.
\end{itemize}
\pendencia{a janela de cadastro não leva à Política de Privacidade.}
\end{historia}

% ---------------------------------------------------------------------------
\feature{F11.3}{Política de Privacidade e Termos de Uso}{\parcial}

\begin{historia}{H11.3.1}{Ler o que a plataforma faz com os dados}{\parcial}
\hu{família ou encarregado de dados da rede}{ler em linguagem clara o que a
plataforma coleta e por quê}{confiar na plataforma antes de a criança usar}
\aceite
\begin{itemize}
  \item A Política de Privacidade tem 17 seções e os Termos de Uso 16, com
    sumário lateral que acompanha a leitura.
  \item Os Termos têm \enquote{As experiências são simulações --- leia esta
    parte}, \enquote{O que cada perfil pode fazer} e \enquote{A inteligência
    artificial na plataforma}.
  \item Os dois documentos só afirmam o que a plataforma faz de verdade, e
    terminam com \enquote{O que ainda está sendo definido}.
  \item A Política de Privacidade está na versão 1.1, em vigor desde 25 de
    setembro de 2026. Os Termos de Uso, ainda um rascunho para revisão
    jurídica, foram atualizados no mesmo dia.
  \item Os dois já dizem as regras da A Seleção: o colega com quem a
    publicação foi pareada no Mural vê os três posts, a legenda e o primeiro
    nome do estudante; e a IA lê a legenda para sugerir ao professor um nível,
    que a criança não vê.
\end{itemize}
\pendencia{razão social e CNPJ da Unooria ainda em branco; e estão declarados
como indefinidos os prazos de guarda e eliminação, o registro de acesso
individual, o canal para responsáveis e o relatório de impacto. Os Termos ainda
citam o UnooTube, arquivado, como exemplo de aplicativo imitado: a frase está
com o PO para revisar.}
\end{historia}

% ===========================================================================
%  E12 — Confiança: privacidade, segurança e operação
% ===========================================================================
\chapter{E12 \textperiodcentered{} Confiança: privacidade, segurança e operação}
\label{ep:e12}

\begin{fichaepico}
\begin{campos}
  \item[Objetivo] A rede, a escola e a família podem confiar na plataforma:
    ela protege o dado de criança, aguenta a manhã de aula, e a equipe sabe
    quando algo quebra.
  \item[Para quem] Secretaria, encarregado de dados da rede, família --- e,
    indiretamente, todos os papéis.
  \item[Valor] São as histórias que nenhuma tela mostra, mas que um edital, um
    jurídico ou um incidente cobram. Várias delas bloqueiam a entrada em
    produção.
  \item[Estado] \parcial{} --- isolamento e proteção de menores prontos;
    teste de carga, monitoramento em produção, retenção e acessibilidade
    pendentes.
\end{campos}
\end{fichaepico}

\begin{tabularx}{\linewidth}{@{}L{16mm}X L{20mm}@{}}
\toprule
\headrow \thd{ID} & \thd{Feature} & \thd{Estado}\\
\midrule
\id{F12.1} & Uma escola nunca vê outra & \pronto\\
\zebra \id{F12.2} & Dado de criança protegido & \pronto\\
\id{F12.3} & Guarda e eliminação de dados & \afazer\\
\zebra \id{F12.4} & Aguentar as 8h da manhã & \parcial\\
\id{F12.5} & A equipe sabe quando algo quebra & \parcial\\
\zebra \id{F12.6} & Acessibilidade & \afazer\\
\id{F12.7} & Português do Brasil & \pronto\\
\bottomrule
\end{tabularx}

% ---------------------------------------------------------------------------
\feature{F12.1}{Uma escola nunca vê outra}{\pronto}

\begin{historia}{H12.1.1}{Isolamento garantido entre escolas}{\pronto}
\hu{secretaria}{garantia de que uma escola nunca vê dado de criança de outra
escola}{atender a LGPD e o edital}
\aceite
\begin{itemize}
  \item A separação é garantida na própria base de dados, não só nas telas:
    vale mesmo que haja um erro de programação.
  \item As únicas exceções são as configurações da rede (disponibilidade,
    limites de IA) e os cadastros do site, que não são dado de estudante.
  \item Existe teste automatizado que tenta atravessar escolas e falha.
  \item Cada rede tem instalação própria: um incidente numa rede não atinge
    outra.
\end{itemize}
\end{historia}

% ---------------------------------------------------------------------------
\feature{F12.2}{Dado de criança protegido}{\pronto}

\begin{historia}{H12.2.1}{Nada que identifique a criança onde não precisa}{\pronto}
\hu{encarregado de dados da rede}{que os registros de uso e os relatórios não
identifiquem a criança sem necessidade}{reduzir o risco de vazamento}
\aceite
\begin{itemize}
  \item Os registros de uso carregam um identificador opaco, nunca nome ou
    e-mail; o CPF nunca aparece em registro técnico.
  \item Os relatórios de gestão são agregados; a linha de uma criança é só do
    professor dela.
  \item Onde a missão pede texto, há limite de tamanho, e dado pessoal é
    recusado (Manifesto da IA).
  \item Na Vila do Lume o nome da criança nunca é gravado.
  \item A exceção declarada é A Seleção, em que um colega lê o que a criança
    publicou: no Mural, o colega com quem a publicação foi pareada vê os três
    posts, a legenda e o primeiro nome dela (ou o nome da equipe). Por isso a
    legenda recusa link, telefone, e-mail, marcação de pessoa e palavrão, e a
    Política de Privacidade e os Termos de Uso dizem isso.
\end{itemize}
\pendencia{uma ferramenta de gravação de sessão de terceiros é usada nas telas
e precisa constar no termo de tratamento de dados de cada rede --- ou ser
retirada das telas de menores.}
\end{historia}

% ---------------------------------------------------------------------------
\feature{F12.3}{Guarda e eliminação de dados}{\afazer}

\begin{historia}{H12.3.1}{Prazos de guarda definidos e cumpridos}{\afazer}
\hu{encarregado de dados da rede}{saber por quanto tempo cada dado de criança
é guardado, e que ele seja apagado no prazo}{cumprir a LGPD e responder à
família}
\aceiteproposta
\begin{itemize}
  \item Um prazo definido para cada tipo de dado (resultado de missão, texto
    livre, registros de uso).
  \item Uma rotina automática apaga o que passou do prazo.
  \item A Política de Privacidade passa a declarar esses prazos.
\end{itemize}
\pendencia{decisão de produto e jurídica anterior ao desenvolvimento (\id{D09}).}
\end{historia}

% ---------------------------------------------------------------------------
\feature{F12.4}{Aguentar as 8h da manhã}{\parcial}

\begin{historia}{H12.4.1}{A rede inteira entrando ao mesmo tempo}{\parcial}
\hu{secretaria}{que a plataforma aguente a rede inteira começando a aula ao
mesmo tempo}{que nenhuma turma perca a aula às 8h}
\aceite
\begin{itemize}
  \item 2 a 4 mil pessoas entrando e abrindo missão em 10 minutos, sem erro e
    sem lentidão perceptível.
  \item Inclui turmas inteiras conversando com a IA ao mesmo tempo (unooGPT).
\end{itemize}
\pendencia{a plataforma está configurada para isso, mas o teste de carga
nunca foi feito. É o maior risco aberto antes de novembro.}
\end{historia}

% ---------------------------------------------------------------------------
\feature{F12.5}{A equipe sabe quando algo quebra}{\parcial}

\begin{historia}{H12.5.1}{Saber do problema antes da escola ligar}{\parcial}
\hu{equipe da Unooria}{ser avisada quando a plataforma falha}{resolver antes que
a escola perceba}
\aceite
\begin{itemize}
  \item Meta: 99,5\% das ações sem erro.
  \item Verificações automáticas a cada minuto, de três regiões; alertas por
    e-mail; registro dos erros que acontecem no navegador de quem usa,
    inclusive na tela de entrada.
  \item Nenhum dado pessoal nesses registros.
\end{itemize}
\pendencia{tudo isso está funcionando só no ambiente de desenvolvimento;
falta levar para homologação e produção, testar o canal de e-mail e ter um
alerta de custo.}
\end{historia}

% ---------------------------------------------------------------------------
\feature{F12.6}{Acessibilidade}{\afazer}

\begin{historia}{H12.6.1}{Usar a plataforma com deficiência}{\afazer}
\hu{estudante com deficiência visual ou motora}{usar a plataforma com leitor
de tela e teclado}{participar da aula como os colegas}
\aceiteproposta
\begin{itemize}
  \item Contraste, foco visível, navegação por teclado e leitor de tela
    verificados em todas as telas da plataforma e nas missões.
  \item A transcrição de todo áudio continua sempre visível (isso já existe).
\end{itemize}
\pendencia{hoje não há verificação sistemática; a correção deve ser feita no
componente visual compartilhado, para valer em todas as telas de uma vez.
Até lá, não é possível prometer acessibilidade a uma rede.}
\end{historia}

% ---------------------------------------------------------------------------
\feature{F12.7}{Português do Brasil}{\pronto}

\begin{historia}{H12.7.1}{Tudo na minha língua}{\pronto}
\hu{qualquer pessoa}{que toda a plataforma fale português do
Brasil}{não tropeçar em termo estrangeiro}
\aceite
\begin{itemize}
  \item Toda a interface da plataforma está em português do Brasil.
  \item Dentro do aplicativo imitado, ficam os termos do aplicativo real
    (Feed, Stories, Salvos), e os nomes Unoovista, Instanooria e unooGPT não se traduzem.
  \item \enquote{fake news}, \enquote{cyberbullying} e \enquote{print} ficam
    como estão, por decisão.
  \item Datas no formato brasileiro e no fuso da rede.
\end{itemize}
\end{historia}

% ===========================================================================
%  E13 — Equipes na missão
% ===========================================================================
\chapter{E13 \textperiodcentered{} Equipes na missão}
\label{ep:e13}

\begin{fichaepico}
\begin{campos}
  \item[Objetivo] Quando os estudantes dividem um computador, o professor
    organiza a turma em equipes, e a plataforma passa a saber quem fez a
    missão junto.
  \item[Para quem] Professor (monta as equipes e lê a turma) e Estudante (vê a
    própria equipe). A Coordenação consulta, sem mexer.
  \item[Valor] A maioria das escolas não tem um computador para cada
    estudante. Na sala com computador dividido, quem entra na missão
    representa a equipe --- e, sem saber disso, a plataforma conta a sala como
    se cada criança tivesse a própria tela: \enquote{12 não começaram}, com os
    12 sentados ao lado de quem começou.
  \item[Estado] \parcial{} --- montar, mostrar e contar por computador estão
    prontos e em revisão; falta o Instanooria formar equipes do mesmo lado do
    feed.
\end{campos}
\end{fichaepico}

\begin{tabularx}{\linewidth}{@{}L{16mm}X L{20mm}@{}}
\toprule
\headrow \thd{ID} & \thd{Feature} & \thd{Estado}\\
\midrule
\id{F13.1} & Montar as equipes & \pronto\\
\zebra \id{F13.2} & O aluno sabe a sua equipe & \pronto\\
\id{F13.3} & Acompanhar por computador & \pronto\\
\zebra \id{F13.4} & Instanooria: equipe de lado igual & \afazer\\
\bottomrule
\end{tabularx}

\section{O que é uma equipe aqui}

\textbf{Equipe é quem divide um computador numa missão}: dois ou três
estudantes diante da mesma tela, com o login de um deles. Não é grupo
pedagógico nem equipe fixa da turma --- vale para uma missão numa turma, e a
aula seguinte pode ter outras equipes. Também não se confunde com as equipes de
resgate da Missão Tempestade, que são parte do jogo.

\begin{nota}
\textbf{Pronto, em revisão.} As histórias marcadas \pronto{} neste épico
funcionam de ponta a ponta e têm o teste que percorre a tela, mas em 25 de
setembro de 2026 ainda estavam em revisão pelo time, antes de entrar na
plataforma, e aguardam a homologação. Se a revisão pedir ajuste, a história
volta a \parcial{}, com a pendência escrita.
\end{nota}

\section{O que fica fora desta onda}

\begin{decisao}
\textbf{Decidido em 25 de setembro de 2026.}
\begin{itemize}
  \item \textbf{Só três missões aceitam equipe}: Instanooria, Missão
    Tempestade e UnooTube. Nas outras, a ficha nem mostra a seção de equipes.
  \item \textbf{A Vila do Lume fica fora}: o retrato de valores que ela
    produz é uma medida de cada criança, e é ele que define o lado de largada
    no Instanooria (\id{H06.6.3}). Duas crianças numa tela dariam um retrato
    que não é de nenhuma das duas.
  \item \textbf{O Eco no Grupo fica fora}: A Seleção o substitui
    (\id{H06.10.2}) e já nasce com equipes.
  \item \textbf{O Que a Foto Conta, Quem mandou isso?, O Rastro do Boato e O
    Manifesto da IA ficam fora}, porque serão retirados da trilha.
  \item \textbf{O atrasado não entra numa equipe automaticamente}: numa sala
    com computador dividido, quem chega atrasado senta num computador que já
    entrou na missão. Se o professor quiser registrar, põe o estudante na
    equipe à mão (\id{H13.1.3}).
  \item \textbf{A evidência pedagógica por habilidade da BNCC, valendo para
    todos da equipe}, é outro épico: \id{E14} Evidência pedagógica
    (p.~\pageref{ep:e14}), com a base pronta e à espera da primeira missão com
    régua.
\end{itemize}
\end{decisao}

% ---------------------------------------------------------------------------
\feature{F13.1}{Montar as equipes}{\pronto}

As equipes se montam na ficha da missão, aberta pela sala (\enquote{Conduzir
com esta turma}), na seção \enquote{Equipes desta turma}: \enquote{Para quando
alunos dividem um computador: quem entra na missão representa a equipe.}

\begin{historia}{H13.1.1}{Dizer quantos computadores há e quem faltou}{\pronto}
\hu{Professor}{dizer quantos computadores a turma vai usar e quem faltou
hoje}{que as equipes sejam feitas só com quem está na sala}
\aceite
\begin{itemize}
  \item \enquote{Quantos computadores vocês vão usar?} aceita de 1 até o
    número de estudantes, e logo abaixo a tela diz o que isso dá (ex.:
    \enquote{25 estudantes em 10 computadores: equipes de 2 ou 3.}).
  \item Com um computador para cada um, a tela diz \enquote{cada um no seu}
    --- e isso também vale.
  \item \enquote{Quem faltou hoje?} abre a lista da turma: quem é marcado fica
    fora do sorteio e aparece em \enquote{Sem equipe}.
  \item Marcar a falta aqui \textbf{não é chamada}: não vale como registro de
    frequência e serve só para montar as equipes desta missão.
\end{itemize}
\end{historia}

\begin{historia}{H13.1.2}{Escolher entre sortear e separar eu mesmo}{\pronto}
\hu{Professor}{escolher, a cada vez, se a plataforma sorteia as equipes ou se
eu mesmo separo}{usar o que eu sei da minha turma quando for preciso}
\aceite
\begin{itemize}
  \item A plataforma \textbf{sempre pergunta} \enquote{Como montar as
    equipes?}, com \enquote{Sortear as equipes} e \enquote{Eu separo}, sem
    nenhuma opção já marcada.
  \item \enquote{Sortear as equipes} divide quem está presente em equipes de
    tamanho parecido, e dá para mudar quem fica onde antes de salvar;
    \enquote{Sortear de novo} refaz o sorteio.
  \item Os tamanhos diferem em no máximo um: 7 estudantes em 3 computadores
    viram equipes de 3, 2 e 2 --- nunca 3, 3 e 1, com alguém sobrando
    sozinho.
  \item \enquote{Eu separo}: \enquote{A plataforma cria 10 equipes vazias e
    você põe cada estudante numa delas} --- uma equipe vazia por computador.
\end{itemize}
\end{historia}

\begin{historia}{H13.1.3}{Ajustar as equipes à mão}{\pronto}
\hu{Professor}{mudar as equipes antes de salvar}{separar quem não trabalha bem
junto, ou juntar quem precisa de ajuda}
\aceite
\begin{itemize}
  \item Posso mudar cada estudante de equipe, renomear uma equipe, criar outra
    (\enquote{Nova equipe}) e remover uma equipe.
  \item Quem ainda não está em nenhuma equipe aparece em \enquote{Sem equipe};
    com todos distribuídos, a tela diz \enquote{Todo mundo está numa equipe.}
  \item Quem foi marcado como falta continua escolhível: o atrasado que chegar
    pode ser posto numa equipe à mão.
  \item A plataforma não deixa salvar equipe vazia (\enquote{Há uma equipe sem
    ninguém. Ponha alguém nela ou remova a equipe.}), equipe sem nome
    (\enquote{Toda equipe precisa de um nome.}) nem duas equipes com o mesmo
    nome.
\end{itemize}
\end{historia}

\begin{historia}{H13.1.4}{Salvar, ver na sala e desfazer}{\pronto}
\hu{Professor}{salvar as equipes e saber, pela sala, que elas estão
prontas}{chegar na aula com a turma já organizada}
\aceite
\begin{itemize}
  \item Ao salvar, a tela confirma \enquote{Equipes salvas.} e mostra um
    cartão por equipe.
  \item Na sala da turma, a missão ganha a etiqueta \enquote{5 equipes
    prontas} --- só informa, não é um atalho.
  \item Cada estudante fica em uma equipe só por missão.
  \item Se alguém salvou equipes para a turma enquanto eu editava, a tela
    avisa \enquote{Estas equipes mudaram em outro lugar} e oferece
    \enquote{Recarregar}: nada é sobrescrito sem que eu veja.
  \item \enquote{Desfazer equipes} pede confirmação (\enquote{Desfazer as
    equipes?}): as equipes desta missão nesta turma deixam de existir, e dá
    para montar de novo quando quiser.
  \item As equipes continuam editáveis mesmo depois de a missão ser
    encerrada.
\end{itemize}
\end{historia}

\begin{historia}{H13.1.5}{Repetir as equipes da última missão}{\pronto}
\hu{Professor}{reaproveitar as equipes que a turma já usou}{não montar tudo de
novo a cada missão}
\aceite
\begin{itemize}
  \item Quando a turma já tem equipes salvas em outra missão, aparece uma
    terceira opção: \enquote{Manter as equipes da última missão}, que diz de
    qual missão elas vêm e quando foram salvas.
  \item Quem faltou hoje sai da equipe e fica em \enquote{Sem equipe}; quem saiu
    da turma desde então não volta, e a tela diz quantos são.
  \item Nada fica valendo antes de eu salvar: posso ajustar como em
    \id{H13.1.3}.
\end{itemize}
\end{historia}

\begin{historia}{H13.1.6}{Mostrar as equipes para a turma}{\pronto}
\hu{Professor}{projetar as equipes para a sala}{que cada estudante ache o
computador dele sem eu ler a lista em voz alta}
\aceite
\begin{itemize}
  \item \enquote{Mostrar equipes} abre uma janela em tela cheia com o nome de
    cada equipe e o \textbf{primeiro nome} de cada estudante.
  \item Nada de faltas, de desempenho nem de quem ficou \enquote{Sem equipe}:
    é uma tela que a turma inteira vê.
  \item Quando dois estudantes da turma têm o mesmo primeiro nome, entra a
    inicial do sobrenome (\enquote{Maria S.}). Se a inicial também empata,
    entra o último sobrenome inteiro (\enquote{Gabriela Cardoso} e
    \enquote{Gabriela Correia}): cada estudante acha o próprio nome, e o nome
    completo só aparece quando nada mais curto separa os dois.
\end{itemize}
\end{historia}

\begin{historia}{H13.1.7}{Consultar as equipes, sem mexer}{\pronto}
\hu{Coordenação da escola}{ver as equipes que o professor montou}{acompanhar a
aula sem interferir nela}
\aceite
\begin{itemize}
  \item A Coordenação vê as equipes salvas, ou \enquote{Nenhuma equipe montada
    para esta turma.}
  \item Não há botão de montar, editar nem desfazer: só quem dá aula na turma
    monta as equipes.
\end{itemize}
\end{historia}

% ---------------------------------------------------------------------------
\feature{F13.2}{O aluno sabe a sua equipe}{\pronto}

\begin{historia}{H13.2.1}{Ver com quem divido o computador}{\pronto}
\hu{Estudante}{ver, ao abrir a missão, qual é a minha equipe e com quem eu
sento}{achar o meu lugar sem depender do quadro}
\aceite
\begin{itemize}
  \item Uma faixa de uma linha, acima da missão, diz \enquote{Sua equipe:
    Equipe 3 \textperiodcentered{} com Ana e Bruno}.
  \item Os colegas aparecem só pelo primeiro nome.
  \item Quando o primeiro nome repete na equipe (o meu inclusive), entra a
    inicial do sobrenome (\enquote{Gabriela C.}); se a inicial também empata,
    o último sobrenome inteiro (\enquote{Gabriela Cardoso}). É a mesma regra
    da tela projetada (\id{H13.1.6}): o nome completo só aparece quando nada
    mais curto separa os dois.
  \item A faixa é da plataforma, fora do aplicativo imitado: a missão em si não
    muda.
  \item Sem equipe nesta missão, a faixa não aparece e a missão abre como
    sempre.
\end{itemize}
\end{historia}

% ---------------------------------------------------------------------------
\feature{F13.3}{Acompanhar por computador}{\pronto}

\begin{historia}{H13.3.1}{Ler a turma por computador}{\pronto}
\hu{Professor}{que os resultados do UnooTube e da Missão Tempestade contem
computadores quando a turma está em equipes}{saber quem de fato não começou,
em vez de ver a sala inteira como atrasada}
\aceite
\begin{itemize}
  \item Com equipes, cada computador é uma linha: a coluna passa a se chamar
    \enquote{Estudante ou equipe} e mostra o nome da equipe, quem senta nela e
    quem entrou.
  \item Os totais contam computadores (\enquote{Computadores na turma}, no
    lugar de \enquote{Estudantes na turma}), com o aviso \enquote{Contado por
    computador: cada equipe desta turma é uma linha, com os dados de quem
    entrou na missão.}
  \item Se mais de um da equipe entrou, a linha mostra quem entrou por último.
  \item Sem equipes, os resultados continuam como antes: um estudante por
    linha.
\end{itemize}
\end{historia}

\begin{historia}{H13.3.2}{Guias que explicam a aula em equipe}{\pronto}
\hu{Professor}{que o guia da aula me diga como conduzir a missão com
computador dividido}{não descobrir as regras no meio da aula}
\aceite
\begin{itemize}
  \item Os guias do UnooTube e da Missão Tempestade apontam para
    \enquote{Equipes desta turma}, na ficha da missão (Minhas turmas \seta{} a
    turma \seta{} Conduzir com esta turma).
  \item O da Missão Tempestade avisa que as 3 perguntas e o limite diário de IA
    são da conta de quem entrou: a equipe negocia em voz alta o que perguntar,
    e o que for escrito fica no nome de quem entrou.
  \item O do UnooTube lembra que os 3 vídeos salvos são da equipe.
\end{itemize}
\end{historia}

% ---------------------------------------------------------------------------
\feature{F13.4}{Instanooria: equipe de lado igual}{\afazer}

No Instanooria, cada estudante começa num lado do feed (\id{H06.6.3}). Hoje o
professor já consegue montar equipes para essa missão, mas o sorteio ainda não
olha o lado: uma equipe pode juntar estudantes dos dois lados, e quem entra
recebe o próprio lado. Esta é a próxima frente do épico.

\begin{historia}{H13.4.1}{Sortear equipes do mesmo lado do feed}{\afazer}
\hu{Professor}{que o sorteio junte na mesma equipe estudantes do mesmo lado do
feed}{que cada computador viva uma bolha só}
\aceiteproposta
\begin{itemize}
  \item Quem tem perfil da Vila do Lume vai para uma equipe do seu lado.
  \item Quem não tem perfil é \enquote{neutro} e é sorteado para completar as
    equipes.
  \item Quando a conta não fecha, a tela diz quantas equipes ficaram mistas.
  \item Os tamanhos continuam diferindo em no máximo um (\id{H13.1.2}).
\end{itemize}
\end{historia}

\begin{historia}{H13.4.2}{Quem entra recebe o lado da equipe}{\afazer}
\hu{Professor}{que qualquer um da equipe que entre na missão receba o mesmo
lado do feed}{que a bolha da equipe não dependa de quem pegou o teclado}
\aceiteproposta
\begin{itemize}
  \item O estudante que entra recebe o lado da equipe, e não o próprio.
  \item Se outro da equipe entrar depois, recebe o mesmo lado.
\end{itemize}
\end{historia}

\begin{historia}{H13.4.3}{Sem redistribuir o lado quando há equipe}{\afazer}
\hu{Professor}{que a turma com equipes não tenha o lado redistribuído}{que a
redistribuição não separe o lado de quem divide a mesma tela}
\aceiteproposta
\begin{itemize}
  \item Com equipes montadas, \enquote{Redistribuir a turma} fica desligado, e
    a tela diz por quê: o lado passa a ser da equipe.
  \item Sem equipes, a redistribuição funciona como hoje (\id{H06.6.3}).
\end{itemize}
\end{historia}

\begin{historia}{H13.4.4}{Ver as equipes no monitor da aula}{\afazer}
\hu{Professor}{ver as equipes no monitor do Instanooria}{acompanhar cada
computador ao vivo, como já faço no UnooTube e na Missão Tempestade}
\aceiteproposta
\begin{itemize}
  \item O monitor mostra o nome de cada equipe e quem senta nela.
  \item Sem equipes, o monitor continua como hoje (\id{H06.6.6}).
\end{itemize}
\end{historia}

% ===========================================================================
%  90 — Backlog: o que está pela metade ou por fazer
% ===========================================================================
\chapter{Backlog: o que falta}
\label{cap:backlog}

Este capítulo junta, num lugar só, tudo o que as histórias marcaram como
\parcial{} ou \afazer{}, mais o que o levantamento encontrou fora das telas.
Está organizado por \textbf{horizonte}, não por épico. A coluna \enquote{Quem}
diz de onde vem o trabalho: nem tudo é desenvolvimento.

\begin{tabularx}{\linewidth}{@{}L{22mm}X@{}}
\toprule
\headrow \thd{Quem} & \thd{Significado}\\
\midrule
Dev & Trabalho do time de desenvolvimento.\\
\zebra Conteúdo & Material pedagógico ou de mídia: áudio, foto, revisão de
  guia, teste em sala.\\
Jurídico & Contrato, termo de tratamento de dados, documentos legais.\\
\zebra Produto & Decisão do PO antes de qualquer desenvolvimento (ver o
  Capítulo~\ref{cap:decisoes}).\\
Operação & Configuração ou publicação em cada rede.\\
\zebra Externo & Serviço contratado fora da Unooria, como a revisão de
  segurança.\\
\bottomrule
\end{tabularx}

\section{Bloqueia a entrada em produção (novembro de 2026)}

\begin{tabularx}{\linewidth}{@{}L{15mm}X L{20mm}@{}}
\toprule
\headrow \thd{Ref.} & \thd{O que falta} & \thd{Quem}\\
\midrule
\endhead
\id{H12.4.1} & Teste de carga: 2 a 4 mil pessoas entrando e abrindo missão em
  10 minutos, com turmas inteiras usando IA ao mesmo tempo. Nunca foi feito. &
  Dev\\
\zebra \id{H12.5.1} & Monitoramento e alertas em homologação e produção (hoje
  só em desenvolvimento), com teste do canal de e-mail e alerta de custo. &
  Dev\\
\id{H12.3.1} & Prazos de guarda e eliminação do dado de criança, e a rotina
  que os cumpre. & Produto, Jurídico, Dev\\
\zebra --- & Revisão de segurança externa e teste de intrusão. & Externo\\
\id{H09.3.2} & Contrato e termo de tratamento de dados de cada rede declarando
  que a IA processa fora do Brasil. & Jurídico\\
\zebra \id{H12.2.1} & Ferramenta de gravação de sessão de terceiros nas telas:
  declarar no termo de cada rede ou retirar das telas de menores (\id{D13}). &
  Produto, Jurídico\\
\id{H01.5.1} & Login único da rede (WSO2) validado em homologação. & Dev,
  Operação\\
\zebra \id{H11.3.1} & Razão social e CNPJ nos documentos legais. & Jurídico\\
\bottomrule
\end{tabularx}

\section{Esperado na primeira rede}

\begin{tabularx}{\linewidth}{@{}L{15mm}X L{20mm}@{}}
\toprule
\headrow \thd{Ref.} & \thd{O que falta} & \thd{Quem}\\
\midrule
\endhead
\id{H07.1.4} & Texto livre só para o professor também no Rastro do Boato e no
  Manifesto da IA (\id{D03}). & Dev\\
\zebra \id{H09.4.2} & IA do Instanooria e da Vila do Lume dentro do limite
  diário. & Dev\\
\id{H12.6.1} & Acessibilidade verificada (contraste, foco, teclado, leitor de
  tela) --- costuma ser exigência de edital. & Dev\\
\zebra \id{H01.4.1} \id{H01.4.2} & Trocar a senha por conta própria e recuperar
  a senha esquecida (\id{D05}). & Dev, Produto\\
\id{H03.4.2} & Escolher a escola numa lista, em vez de digitar o código. & Dev\\
\zebra \id{H07.2.1} & Nome da missão (e não o código) no relatório da
  diretoria. & Dev\\
\id{H07.3.1} & Análise de produto no menu, com botão de exportar e seleção de
  escola. & Dev\\
\zebra \id{H11.1.2} & Botão \enquote{Entrar} e links para os documentos legais
  no site; link da Política de Privacidade no cadastro (\id{H11.2.1}). & Dev\\
\id{H01.3.1} & Trocar de escola também no celular. & Dev\\
\zebra \id{H03.3.1} & Filtro Papel completo em Usuários. & Dev\\
\id{H01.2.2} & Retirar o item \enquote{Iniciativas}, que não faz nada, dos
  menus laterais. & Dev\\
\zebra \id{H06.1.1} & Fotos reais e locuções do O Que a Foto Conta. & Conteúdo\\
\id{H06.7.4} \id{H06.4.4} & Áudios do Manifesto da IA e do Rastro do Boato. &
  Conteúdo\\
\zebra \id{H06.8.4} \id{H04.4.1} & Revisão pedagógica e teste em sala do
  guia do unooGPT; autor e revisor do gabarito do Rastro. &
  Conteúdo\\
\id{H10.5.2} & Eliminar um cadastro de prospecção pela tela. & Dev\\
\zebra \id{H06.10.2} \id{H06.12.1} \id{H06.12.4} \id{H06.12.5} & A Seleção (7º ano,
  Unidade 1, \id{F06.12}), já construída e percorrida pelo teste: revisar as
  imagens geradas, conferir com a IA de verdade os níveis que ela dá à legenda
  e revisar os textos escritos pela equipe. & Conteúdo, Dev, Produto\\
\id{H06.10.5} \id{H06.12.6} & Revisar a voz gerada do podcast da A
  Seleção; se não passar, gravar o episódio em estúdio, com o mesmo roteiro. &
  Conteúdo\\
\zebra \id{H06.10.3} \id{H06.11.1} \id{H06.11.4} \id{H06.11.5} & Quem mandou isso?
  versão nova (5º ano, Unidade 1, \id{F06.11}): revisar os textos escritos pela
  equipe e decidir o equilíbrio do pacote pronto (\id{D14}). & Dev, Conteúdo,
  Produto\\
\id{H06.10.4} & Jornal da Escola no lugar do UnooTube (4º ano, Unidade
  1), para o 4º ano voltar a ter missão. & Dev, Conteúdo\\
\zebra \id{H11.3.1} & Revisar a frase dos Termos de Uso que ainda cita o UnooTube,
  arquivado, como aplicativo imitado. & Produto\\
\id{H14.2.3} \id{H14.3.1} \id{H14.3.2} \id{H14.3.3} \id{H14.4.2}
  \id{H14.4.3} \id{H14.6.1} & Teste de ponta a ponta da evidência pedagógica
  que passe pelo que o de hoje não percorre: as equipes, a aula projetada
  (\enquote{a revisar} e \enquote{Confirmar para a turma}), o ajuste de vários
  estudantes e o histórico, o \enquote{Não participou}, a ordem da lista e um
  perfil com conteúdo. & Dev\\
\zebra \id{F13.4} \id{H13.4.1} \id{H13.4.2} \id{H13.4.3} \id{H13.4.4} & Equipes do
  mesmo lado do feed no Instanooria: o sorteio agrupa pelo lado, quem entra
  recebe o lado da equipe, a redistribuição fica desligada e o monitor mostra
  as equipes. É a próxima frente de \id{E13}. & Dev\\
\bottomrule
\end{tabularx}

\section{Evolução}

\begin{tabularx}{\linewidth}{@{}L{15mm}X L{20mm}@{}}
\toprule
\headrow \thd{Ref.} & \thd{O que falta} & \thd{Quem}\\
\midrule
\endhead
\id{H04.2.3} & Liberar a missão com data de abertura e fechamento (\id{D12}). &
  Dev\\
\zebra \id{H10.1.3} & Desligar recurso por turma ou até uma data. & Dev\\
\id{H10.3.3} & Aviso de versão fixada há muito tempo. & Dev\\
\zebra \id{H02.4.2} & Virada do ano assistida (a primeira real é em fevereiro de
  2027). & Dev\\
\id{H09.4.3} & Aviso quando o consumo de IA de uma rede chegar a 80\%. & Dev\\
\zebra \id{H07.3.2} & Números do dia para a rede. & Dev\\
\id{H05.2.2} & Guardar os registros de uso sem conexão. & Dev\\
\zebra \id{F06.10} \id{H06.10.1} & As 41 unidades da trilha ainda em produção,
  cada uma entrando pelo mesmo caminho, já completa (\id{H06.10.1}). &
  Conteúdo, Dev\\
\id{H08.3.4} & Visão agregada e anônima do bem-estar para a rede --- só se
  decidida (\id{D11}). & Produto, Jurídico\\
\zebra \id{H14.1.1} & A régua pedagógica nas outras missões. A primeira ---
  Quem mandou isso?, versão nova --- já entrou e tornou visível ao professor o
  que \id{E14} faz, e A Seleção veio em seguida; o Instanooria é a próxima
  candidata. & Conteúdo, Dev\\
\id{F14.1} & O unooGPT refeito com atividades, para entrar na
  evidência pedagógica. & Dev, Conteúdo\\
\zebra \id{H14.4.4} & A tela em que a coordenação consulta a evidência (a
  leitura já é permitida). & Dev\\
\id{H14.6.2} & O painel que junta as habilidades de todas as missões, com a
  regra definida pela equipe pedagógica. & Produto, Conteúdo, Dev\\
\bottomrule
\end{tabularx}

% ===========================================================================
%  91 — Decisões em aberto para o produto
% ===========================================================================
\chapter{Decisões em aberto para o produto}
\label{cap:decisoes}

O levantamento encontrou perguntas que o código não responde e que não são de
engenharia: são escolhas de produto. Cada uma tem o contexto, as opções e,
quando há, uma recomendação. Nenhuma foi decidida neste documento.

\newcommand{\decisaoaberta}[4]{%
  \begin{historia}{#1}{#2}{}
  \textbf{Contexto.} #3\par\smallskip
  #4
  \end{historia}}

\decisaoaberta{D01}{Quem cria o ano letivo?}{A tela \enquote{Ano letivo} diz
\enquote{O calendário é do contrato, não da escola}, mas a Coordenação da
escola consegue criar período (\id{H02.4.1}).}{\textbf{Opções:} (a) só o
Administrador do contrato cria, e a escola consulta; (b) a escola cria, e a
frase da tela muda. \textbf{Recomendação:} (a), porque um calendário por
escola quebra o relatório da rede.}

\decisaoaberta{D02}{A Coordenação da escola pode transferir estudante?}{O botão
\enquote{Transferir de escola} aparece só para o Administrador do contrato; a
Coordenação consegue o mesmo efeito importando o CPF (\id{H03.6.1}).}{\textbf{Opções:}
(a) mostrar o botão também para as coordenações; (b) manter só no contrato e
tratar a importação como caminho oficial da escola.}

\decisaoaberta{D03}{O texto livre da criança é só do professor em todas as
missões?}{No Rastro do Boato e no Manifesto da IA a gestão ainda vê o texto;
a versão nova de Quem mandou isso? já segue a regra --- a gestão vê que a regra
foi enviada, e só quem dá aula na turma lê o texto (\id{H07.1.4},
\id{H06.11.5}).}{\textbf{Recomendação:} alinhar todas à regra de \id{H07.1.4}
--- a gestão vê só \enquote{entregou / não entregou} ---, e tornar a regra
obrigatória para toda missão nova (\id{H06.10.1}).}

\decisaoaberta{D04}{A que horas o limite de IA recomeça?}{O contador de cada
pessoa zera à meia-noite no horário universal, que no Brasil é 21h
(\id{H09.4.1}). Um professor que planeja à noite pode ver o limite
\enquote{renascer} no meio do trabalho.}{\textbf{Opções:} (a) meia-noite no
fuso da rede; (b) manter e explicar na tela.}

\decisaoaberta{D05}{Como o estudante recupera a senha?}{Hoje só a coordenação
reseta. Estudante dos anos iniciais quase nunca tem e-mail
(\id{H01.4.2}).}{\textbf{Opções:} (a) recuperação por e-mail só para adultos, e
o estudante segue pela coordenação; (b) o professor também pode resetar a
senha dos estudantes da turma dele. Com o login único da rede (\id{H01.5.1}),
a pergunta pode deixar de existir.}

\decisaoaberta{D06}{A conta terá perfil e foto?}{O menu do avatar tem só
\enquote{Sair}, e isso foi deliberado.}{Confirmar que continua fora do escopo
--- foto de criança é dado sensível e não traz valor pedagógico evidente.}

\decisaoaberta{D07}{O estudante terá um histórico fora das missões?}{Hoje ele
vê só a grade da turma; não há portfólio nem \enquote{o que já fiz}.}{Decidir
se entra no produto e, se entrar, o que ele mostra sem virar nota.}

\decisaoaberta{D08}{Toda missão terá relatório individual?}{Só unooGPT, Vila
do Lume e Instanooria têm (\id{H07.1.3}).}{Decidir se passa a ser exigência
da esteira de missões novas (\id{H06.10.1}) ou se fica a critério de cada
missão.}

\decisaoaberta{D09}{Por quanto tempo guardamos o dado de criança?}{Não há prazo
definido, e a Política de Privacidade declara isso como pendência
(\id{H12.3.1}).}{Definir, com o jurídico, um prazo por tipo de dado --- é
pré-requisito para o desenvolvimento da rotina de eliminação.}

\decisaoaberta{D10}{Qual é o modelo das próximas missões?}{O Instanooria tem
\enquote{atividade} configurável (tema, geração, revisão, etapas), e a versão
nova de Quem mandou isso? tem o pacote de mensagens que o professor prepara e
aprova para a turma, que é a atividade dela (\id{H06.11.1}); as outras seguem
\enquote{liberar e conduzir}. Em 25 de setembro de
2026, com a evidência pedagógica (\id{E14}), ficou decidido que cada missão
define as próprias atividades e que o unooGPT e o UnooTube serão refeitos com
atividades numa próxima onda; o modelo das demais continua em aberto.}{Decidir
o modelo padrão: a atividade configurável dá mais controle ao professor e custa
bem mais para construir e para preparar a aula --- e é por atividade que a
evidência pedagógica é registrada.}

\decisaoaberta{D11}{A rede verá algum agregado do bem-estar?}{Hoje não existe
nenhuma visão de gestão, por desenho (\id{H08.3.4}).}{\textbf{Opções:} (a)
nunca; (b) um agregado anônimo, com grupos mínimos e termo específico com a
rede. Qualquer caminho diferente de (a) precisa do jurídico antes.}

\decisaoaberta{D12}{A liberação com data entra antes de novembro?}{Hoje o
professor libera na hora da aula, e é fácil esquecer (\id{H04.2.3}).}{Decidir
a prioridade.}

\decisaoaberta{D13}{Gravação de sessão por terceiros em tela de criança}{Uma
ferramenta externa de gravação de sessão é usada nas telas
(\id{H12.2.1}).}{\textbf{Opções:} (a) declarar no termo de tratamento de dados
de cada rede; (b) retirar das telas de estudante. \textbf{Recomendação:} (b)
nas telas de menores.}

\decisaoaberta{D14}{No pacote pronto de Quem mandou isso?, imagem e link
equilibrados?}{No material de origem do pacote pronto \enquote{O Mistério do
Passeio ao Museu}, os dois links eram falsos --- e a turma aprenderia que link
é sempre mentira, o atalho que a missão existe para desmontar. Até a resposta,
vale o equilíbrio: a plataforma exige que a imagem e o link tenham uma
verdadeira e uma falsa cada, e no pacote pronto a quinta mensagem (verdadeira)
passou a ser link e a sétima (falsa) passou a ser texto, sem mudar nenhum texto
(\id{H06.11.1}).}{\textbf{Opções:} (a) manter o equilíbrio, para o pacote pronto
e para o gerado; (b) voltar ao material original e aceitar que o formato
entregue a resposta. \textbf{Recomendação:} (a).}


\appendix
% ===========================================================================
%  Apêndice A — As missões e os seus códigos
% ===========================================================================
\chapter{As missões e os seus códigos}
\label{ap:missoes}

O time de desenvolvimento chama cada missão por um código que diz onde ela
mora na trilha: \cd{exp} + ano + etapa + unidade + \cd{b}. Assim,
\cd{exp08ef1b} é o 8º ano do Ensino Fundamental, Unidade 1. O que não é missão
da trilha leva \cd{app-} e um número. A tabela é a ponte entre o nome que o
PO usa e o código que aparece num card técnico.

\begin{tabularx}{\linewidth}{@{}X L{46mm}L{30mm}L{24mm}@{}}
\toprule
\headrow \thd{Nome na tela} & \thd{Código} & \thd{Ano / unidade} & \thd{Produto}\\
\midrule
O Que a Foto Conta & \cd{exp01ef1b} & 1º EF, Unidade 1 & Entre Redes\\
\zebra Jornal da Escola: A Edição de Bento & \cd{exp04ef1b} & 4º EF, Unidade 1 & Entre Redes\\
UnooTube (arquivada) & \cd{backup-exp04ef1b} & 4º EF, Unidade 1 (até 25/09/2026) & Entre Redes\\
\zebra Quem mandou isso? & \cd{exp05ef1b} & 5º EF, Unidade 1 & Entre Redes\\
O Rastro do Boato & \cd{exp06ef2b} & 6º EF, Unidade 2 & Entre Redes\\
\zebra A Seleção & \cd{exp07ef1b} & 7º EF, Unidade 1 & Entre Redes\\
Unoovista & \cd{exp08ef1b} & 8º EF, Unidade 1 & Entre Redes\\
\zebra O Manifesto da IA & \cd{exp08ef4b} & 8º EF, Unidade 4 & Entre Redes\\
UnooriaAI (Missão Tempestade) & \cd{exp02em1b} & 2ª EM, Unidade 1 & Entre Redes\\
\zebra Vila do Lume & \cd{app-70-vila-do-lume} & para qualquer ano & Entre Redes\\
Gerador de Plano de Aula & \cd{app-90-plano-de-aula} & --- & Plano de Aula\\
\zebra Como você está? & \cd{app-91-bem-estar} & --- & Bem-estar\\
Planejamento de Ensino & \cd{app-92-plano-de-ensino} & --- & Plano de Ensino\\
\bottomrule
\end{tabularx}

\begin{atencao}
Até setembro de 2026 as missões tinham outro código (\cd{exp-01}, \cd{exp-02}
\ldots, pela ordem de criação), e a troca embaralhou os números: a antiga
\cd{exp-01} é hoje o Unoovista (\cd{exp08ef1b}; até outubro de 2026, Instanooria), e a antiga \cd{exp-04} é O
Que a Foto Conta (\cd{exp01ef1b}). Documento ou card antigo pode usar o código
velho. E um mesmo código pode ter mudado de missão: \cd{exp04ef1b} foi o
UnooTube até 25 de setembro de 2026 e é hoje o Jornal da Escola --- o UnooTube
guardado atende por \cd{backup-exp04ef1b}. E a Missão Tempestade (\cd{exp02em1b})
se chamou unooGPT até outubro de 2026, quando trocou o clone do ChatGPT pelo
UnooriaAI.
\end{atencao}

% ===========================================================================
%  Apêndice B — Onde cada épico mora na documentação técnica
% ===========================================================================
\chapter{Onde cada épico mora na documentação técnica}
\label{ap:rastreio}

Para o time de desenvolvimento: o detalhe técnico de cada épico está na
documentação técnica (PU-DOC-001). A tabela aponta o capítulo e as decisões
registradas (ADR) que explicam por que a feature é como é.

\begin{tabularx}{\linewidth}{@{}L{10mm}L{44mm}X@{}}
\toprule
\headrow \thd{Épico} & \thd{Capítulos da PU-DOC-001} & \thd{Decisões (ADR)}\\
\midrule
\endhead
\id{E01} & Multi-tenancy e segurança; Contratos de integração & \adrr{ADR-03}
  (OIDC genérico), \adrr{ADR-15} (o hub assina a sessão da missão),
  \adrr{ADR-29} (senha provisória), \adrr{ADR-85} (trocar de escola)\\
\zebra \id{E02} & Multi-tenancy e segurança; Modelo de dados & \adrr{ADR-80}
  (seis papéis e a cadeia), \adrr{ADR-81} (panorama da diretoria)\\
\id{E03} & Multi-tenancy e segurança; Modelo de dados; Especificação da API &
  \adrr{ADR-54} (pessoa nasce por importação), \adrr{ADR-59} (o balcão delega
  ao importador), \adrr{ADR-32} (transferência), \adrr{ADR-35} (perfil
  portátil)\\
\zebra \id{E04} & Catálogo de experiências & \adrr{ADR-69} (liberação por
  turma), \adrr{ADR-86} (guia da aula)\\
\id{E05} & Contratos de integração; Telemetria & \adrr{ADR-01} (missão como
  arquivo estático), \adrr{ADR-09} (correção no servidor)\\
\zebra \id{E06} & Catálogo de experiências & \adrr{ADR-09}, \adrr{ADR-12} (sem
  número com uma medida só), \adrr{ADR-22}, \adrr{ADR-42}, \adrr{ADR-67}
  (mídia preparada pelo professor), \adrr{ADR-72} (a revelação),
  \adrr{ADR-110} (código da missão), \adrr{ADR-108} e \adrr{ADR-114} (a missão
  que nasce vazia, e a primeira versão de Quem mandou isso? retirada),
  \adrr{ADR-118} (a versão nova), \adrr{ADR-120} (a régua v2 dela),
  \adrr{ADR-113} (A Seleção no lugar do Eco no Grupo), \adrr{ADR-115} (A
  Seleção), \adrr{ADR-119} (o podcast da A Seleção), \adrr{ADR-116} (o UnooTube
  arquivado, e não apagado), \adrr{ADR-117} (o Jornal da Escola: a redação é
  da Unooria, o professor abre o resultado e a revisão, e a IA só escolhe o
  layout)\\
\id{E07} & Catálogo de experiências; Multi-tenancy e segurança & \adrr{ADR-12},
  \adrr{ADR-81}, \adrr{ADR-82} (relatório da diretoria)\\
\zebra \id{E08} & Catálogo de experiências; Camada de IA & \adrr{ADR-10}
  (bem-estar sem gestor), \adrr{ADR-76} (planejamento semanal)\\
\id{E09} & Camada de inteligência artificial & \adrr{ADR-07} (IA só pela
  plataforma), \adrr{ADR-67}, \adrr{ADR-74}, \adrr{ADR-105} (limite de IA
  editável), \adrr{ADR-115} (a IA gradua a legenda da A Seleção, exceção
  declarada), \adrr{ADR-119} (a voz gerada do podcast), \adrr{ADR-117} (a IA em sala que
  só monta a página do Jornal)\\
\zebra \id{E10} & Multi-tenancy e segurança; Engenharia de entrega &
  \adrr{ADR-87} (disponibilidade), \adrr{ADR-88} (painel só da Unooria),
  \adrr{ADR-89} (histórico que a rede lê), \adrr{ADR-01} (voltar versão),
  \adrr{ADR-105}, \adrr{ADR-116} (experiências arquivadas)\\
\id{E11} & Multi-tenancy e segurança & \adrr{ADR-70} (páginas legais),
  \adrr{ADR-107} (site e cadastro de interesse)\\
\zebra \id{E12} & Multi-tenancy e segurança; Telemetria; Observabilidade;
  Riscos e evolução & \adrr{ADR-04} (instalação por rede), \adrr{ADR-08}
  (isolamento no banco), \adrr{ADR-91} (idioma), \adrr{ADR-98} e
  \adrr{ADR-99} (fuso)\\
\id{E13} & Arquitetura de software; Multi-tenancy e segurança; Modelo de
  dados; Especificação da API & \adrr{ADR-111} (equipe é quem divide um
  computador), \adrr{ADR-118} (Quem mandou isso? nasce com equipes),
  \adrr{ADR-115} (o par do Mural da A Seleção é do computador),
  \adrr{ADR-117} (o Jornal da Escola também nasce com equipes)\\
\zebra \id{E14} & Arquitetura de software; Multi-tenancy e segurança; Modelo de
  dados; Especificação da API & \adrr{ADR-112} (a evidência é por atividade e
  rodada, e a melhor rodada vale por habilidade), \adrr{ADR-111} (a equipe
  que recebe a cópia), \adrr{ADR-69} (a pergunta só abre com a missão
  liberada), \adrr{ADR-09} (gabarito no servidor), \adrr{ADR-89} (histórico
  append-only), \adrr{ADR-118} (a primeira régua: a atividade é o pacote, a
  rodada é a reaplicação), \adrr{ADR-120} (a régua v2 de Quem mandou isso?),
  \adrr{ADR-115} (a segunda régua: a atividade é a cobertura, e a IA gradua a
  legenda), \adrr{ADR-117} (a terceira régua, a do Jornal: três habilidades,
  e a atividade é o pacote das oito notícias)\\
\bottomrule
\end{tabularx}

\begin{nota}
Os casos de uso da PU-DOC-001 (\ucr{UC-01} a \ucr{UC-28}) cobrem uma parte das
histórias daqui, com outra granularidade: um caso de uso descreve um fluxo do
sistema; uma história descreve a necessidade de uma pessoa. As duas visões
convivem --- este documento não substitui o capítulo de casos de uso.
\end{nota}

% ===========================================================================
%  Apêndice C — Fonte Mermaid dos diagramas
% ===========================================================================
\chapter{Fonte Mermaid dos diagramas}
\label{ap:mermaid}

Todo diagrama deste documento tem uma versão em Mermaid, para que qualquer
pessoa consiga alterá-lo sem saber TikZ.

\textbf{Figura~\ref{fig:caminho} --- O caminho de uma aula}

\begin{lstlisting}[language=mermaid]
flowchart LR
  p1["1. Disponibilizar<br/>o produto e as missões existem nesta rede<br/><i>Administrador Unooria</i>"]
  p2["2. Cadastrar<br/>escolas, pessoas e turmas<br/><i>Coordenação / Adm. do contrato</i>"]
  p3["3. Preparar e liberar<br/>a missão para a turma<br/><i>Professor</i>"]
  p4["4. Jogar<br/>a missão, no ritmo da aula<br/><i>Estudante</i>"]
  p5["5. Conduzir e ler<br/>a turma: debate, revelação, resultados<br/><i>Professor</i>"]
  p6["6. Acompanhar<br/>a adoção, sem identificar criança<br/><i>Coordenação regional / Adm. do contrato</i>"]
  p1 --> p2 --> p3 --> p4 --> p5 --> p6
\end{lstlisting}


\end{document}
